\documentclass[12pt,a4paper]{article}

% Encoding and fonts
\usepackage[utf8]{inputenc}
\usepackage[T1]{fontenc}
\usepackage{lmodern}

% Page layout
\usepackage[top=1in, bottom=1in, left=1in, right=1in]{geometry}

% Typography
\usepackage{microtype}
\usepackage{parskip}

% Language
\usepackage[english]{babel}

% Headers and footers
\usepackage{fancyhdr}
\pagestyle{fancy}
\fancyhf{}
\fancyhead[L]{\leftmark}
\fancyhead[R]{\thepage}
\fancyfoot[C]{\thepage}
\renewcommand{\headrulewidth}{0.4pt}
\renewcommand{\footrulewidth}{0pt}

% Section styling
\usepackage{titlesec}
\titleformat{\section}{\Large\bfseries}{\thesection}{1em}{}
\titleformat{\subsection}{\large\bfseries}{\thesubsection}{1em}{}
\titleformat{\subsubsection}{\normalsize\bfseries}{\thesubsubsection}{1em}{}
\titlespacing*{\section}{0pt}{2.5ex plus 1ex minus .2ex}{1.5ex plus .2ex}
\titlespacing*{\subsection}{0pt}{2ex plus 1ex minus .2ex}{1ex plus .2ex}
\titlespacing*{\subsubsection}{0pt}{1.5ex plus 1ex minus .2ex}{0.8ex plus .2ex}

% List formatting
\usepackage{enumitem}
\setlist[itemize]{topsep=0.5ex, itemsep=0.25ex, parsep=0pt}
\setlist[enumerate]{topsep=0.5ex, itemsep=0.25ex, parsep=0pt}

% Essential packages
\usepackage{graphicx}
\usepackage[table]{xcolor}
\usepackage{booktabs}
\usepackage{array}
\usepackage{tabularx}
\usepackage{longtable}
\usepackage{amsmath}
\usepackage{amssymb}
\usepackage[normalem]{ulem}
\rowcolors{2}{gray!10}{white}
\renewcommand{\arraystretch}{1.3}

% Cross-references
\usepackage{hyperref}
\usepackage{cleveref}

% Hyperref setup
\hypersetup{
    colorlinks=true,
    linkcolor=blue,
    filecolor=magenta,
    urlcolor=cyan,
}

% Code listings
\usepackage{listings}
\definecolor{codebg}{RGB}{248,248,248}
\definecolor{codeframe}{RGB}{200,200,200}
\definecolor{codekeyword}{RGB}{0,0,180}
\definecolor{codecomment}{RGB}{0,128,0}
\definecolor{codestring}{RGB}{163,21,21}
\lstset{
    basicstyle=\ttfamily\small,
    breaklines=true,
    breakatwhitespace=false,
    frame=single,
    framerule=0.5pt,
    rulecolor=\color{codeframe},
    backgroundcolor=\color{codebg},
    keepspaces=true,
    showstringspaces=false,
    tabsize=4,
    xleftmargin=1em,
    xrightmargin=1em,
    aboveskip=1em,
    belowskip=1em,
    keywordstyle=\color{codekeyword}\bfseries,
    commentstyle=\color{codecomment}\itshape,
    stringstyle=\color{codestring},
}

% YAML language definition for listings
\lstdefinelanguage{YAML}{
    keywords={true,false,null,y,n},
    keywordstyle=\color{blue}\bfseries,
    sensitive=false,
    comment=[l]{\#},
    morecomment=[s]{/*}{*/},
    commentstyle=\color{gray}\ttfamily,
    stringstyle=\color{red}\ttfamily,
    morestring=[b]',
    morestring=[b]",
}

% TOML language definition for listings
\lstdefinelanguage{TOML}{
    keywords={true,false},
    keywordstyle=\color{blue}\bfseries,
    sensitive=true,
    comment=[l]{\#},
    commentstyle=\color{gray}\ttfamily,
    stringstyle=\color{red}\ttfamily,
    morestring=[b]',
    morestring=[b]",
}

% Dockerfile language definition for listings
\lstdefinelanguage{Dockerfile}{
    keywords={FROM,RUN,CMD,LABEL,EXPOSE,ENV,ADD,COPY,ENTRYPOINT,VOLUME,USER,WORKDIR,ARG,ONBUILD,STOPSIGNAL,HEALTHCHECK,SHELL},
    keywordstyle=\color{blue}\bfseries,
    sensitive=false,
    comment=[l]{\#},
    commentstyle=\color{gray}\ttfamily,
    stringstyle=\color{red}\ttfamily,
    morestring=[b]',
    morestring=[b]",
}

% JSON language definition for listings
\lstdefinelanguage{JSON}{
    keywords={true,false,null},
    keywordstyle=\color{blue}\bfseries,
    sensitive=true,
    commentstyle=\color{gray}\ttfamily,
    stringstyle=\color{red}\ttfamily,
    morestring=[b]",
}

% Code listing float environment
\usepackage{float}
\newfloat{listing}{htbp}{lol}[section]
\floatname{listing}{Listing}

% Admonition boxes
\usepackage{tcolorbox}
\tcbuselibrary{skins,breakable}

% Admonition style definitions
\newtcolorbox{admonitionbox}[2][]{
    enhanced,
    breakable,
    colback=#2!5,
    colframe=#2!80!black,
    fonttitle=\bfseries,
    title=#1,
    left=2mm,
    right=2mm,
}

% Synthon tuple boxes
\newtcolorbox{synthonbox}{
    enhanced,
    colback=blue!5,
    colframe=blue!75!black,
    fontupper=\large,
    center upper,
    boxrule=1pt,
    arc=4pt,
}

\begin{document}

\title{SynthOmnicon: A Constraint Algebra for Self-Organizing Systems}
\date{\today}

\maketitle

\begin{center}
\rule{0.5\textwidth}{0.4pt}
\end{center}

\subsection{I. The Framework}

\textbf{The central observation.} Systems that self-organize --- that enforce constraints on the states of their partners through reversible or irreversible interactions --- share an ordinal structure regardless of substrate. The imine condensation, the kinase-substrate recognition event, the Cooper pair condensation, and the liquid-to-gel transition in a condensate are not related by physics. They are related by constraint grammar: each specifies a fidelity, a kinetic character, a granularity of control, a grammar of partner selection, and a criticality class. SynthOmnicon's claim is precise: this shared structure is algebraic, and the algebra is predictive.

\textbf{What the algebra produces.} Encoded as eleven-primitive tuples $\langle D; T; R; P; F; K; G; \Gamma; \Phi; S; \Omega \rangle$, systems admit composition under seven operations --- meet ($\sqcap$), join ($\sqcup$), tensor ($\otimes$), lift, path, pipeline, and decomposition. From ordinal structure alone, without numerical parameterization, the algebra produces: the correct competitive displacement ordering in CB{[}7{]} host-guest chemistry (6/6 predictions from $F$-rank alone); the $d = 0.000$ identity between mechanically-primed angiosperm Hv1 and constitutively-active gymnosperm Hv1, collapsing a phylogenetic argument into a single number; the four-conflict isolation of quantum gravity from the Standard Model at the $G = G_{\text{LOCAL}}$ boundary; the $\Gamma$-only conflict driving condensate liquid-to-gel transition, with $F$, $K$, and $G$ all determined downstream; the $+2.303$ nat ($= \ln 10$) criticality-lift cost appearing identically across topological phase transitions, protein folding barriers, and Landauer information bounds. These are numerically conserved values across physically unrelated systems. A coarse-graining that arbitrarily discretised continuous measurements would not produce such coincidences.

\textbf{What the framework is not.} SynthOmnicon is not a Theory of Everything. It makes no ontological claim about what reality is at bottom and derives no specific law from a proposed fundamental substrate. Its claim is more precise and more limited: given any system with internal structure, certain conditional relationships hold --- about what states are accessible, at what cost, and in what order. The primitives identify what a system \emph{is conditional on}, not why it exists. A wrong prediction falsifies the encoding, not the algebra. This is the formal content of \emph{universal conditional logic} (UCL): the same conditional structure appears across domains because those domains share a constraint grammar, not because they share a physical substrate. In the same way that Boolean algebra is not a theory of everything about circuits but is the universal conditional logic of two-valued systems, SynthOmnicon claims to be the Boolean algebra of self-organising systems.

\textbf{Definition.} A \emph{Synthon} is a directed relational operator: a minimal specification of constraint-enforcement capacity defined entirely by its interactions with a compatible context. No primitive in the tuple $\langle D; T; R; P; F; K; G; \Gamma; \Phi; S; \Omega \rangle$ describes an intrinsic property of an isolated object. $F$ is competitive displacement rank --- there is no ''$F_\hbar$ in isolation,'' only ''$F_\hbar$ relative to a specified competitor set.'' $K$ is a barrier relative to an environmental driving force. $\Gamma$ is partner-selection logic that presupposes a partner. $\Phi$ is a phase class across a catalog of interacting configurations. $\Omega$ encodes topological protection against perturbations, meaningless without an environment to perturb. A synthon tuple encodes \emph{interaction affordances} --- what constraints it can enforce, in what order, against which partners, at what scale --- not the constitution of any substance. A tuple without a context is interaction potential; the unit of physical content is the tuple-in-context.

This is not a philosophical gloss. It is a \textbf{type-system requirement}: you cannot assign $F$, $K$, $\Gamma$, or $\Omega$ without specifying an interaction context. The algebra enforces this structurally: every operation in meet, join, tensor, lift, path, and pipeline requires at least one additional operand. There are no unary information generators. The algebra cannot process ''nothing but the object.'' This is why the framework is domain-agnostic by construction --- the primitives are relational, and relations are substrate-independent. The systematic asymmetry of the algebra ($\text{path}(A \to B) \neq \text{path}(B \to A)$, the $F$-floor ratchet is directed, lift has no inverse) places SynthOmnicon in the tradition of structural realism: the world's causal structure is relational but ordered, and the ordering is the load-bearing part.

\textbf{Predictiveness.} A classification system need only assign labels; a predictive grammar must compose primitives and derive non-obvious consequences about assembled system behavior. The primitive basis is evaluated against five criteria: \emph{composability} (primitives combine to generate predictions), \emph{orthogonality} (each captures a distinct dimension of relational structure), \emph{completeness} (every real synthon is encodable), \emph{productivity} (novel encodings yield novel predictions), and \emph{falsifiability} (predictions are disprovable). The constraint algebra is binary by requirement: no operation generates physical information from a single operand. This is a consequence of the definition, not an implementation detail --- a relational operator with one operand is undefined.

\textbf{Falsifiability structure.} The framework is falsifiable at two independent levels. First: if an algebraically-derived prediction fails, the primitive assignment is wrong --- the algebra is tautological given correct assignments, and assignments are empirically determinable (the \emph{algorithmic assignment project}, §XXII.2, specifies the convergence protocol). Second: whether the primitives are \emph{natural joints} --- tracking real scale separations in nature rather than useful conventional bins --- is an open empirical question. The cross-domain numerical coincidences (ln 10, $d = 0.000$, four SM/QG conflicts) provide non-trivial evidence for natural joints without proving them. Both levels are independently evaluable. The UCL claim is stable under either outcome; only the \emph{depth} of the naturalness changes.

\begin{center}
\rule{0.5\textwidth}{0.4pt}
\end{center}

\subsection{II. The Eleven Primitives}

SynthOmnicon encodes any system as an eleven-element tuple. The first ten primitives cover classical and quantum systems; the eleventh ($\Omega$) is a quantum extension active only for topologically protected states:

\[\langle D \;;\; T \;;\; R \;;\; P \;;\; F \;;\; K \;;\; G \;;\; \Gamma \;;\; \Phi \;;\; S \;;\; \Omega \rangle\]

$\Omega$ is optional; classical synthons carry $\Omega_0$ (trivial) or leave the field unset.

\begin{table}[htbp]
\centering
\small
\rowcolors{1}{white}{white}
\begin{tabularx}{\textwidth}{>{\centering\arraybackslash}X >{\centering\arraybackslash}X >{\centering\arraybackslash}X}
\toprule
\rowcolor{gray!25}\textbf{Primitive} & \textbf{Description} & \textbf{Values} \\
\midrule
\rowcolors{1}{white}{gray!10}
\textbf{Dimensionality ($D$)} & Coordinate set along which the synthon operates & $D_{\wedge}$ molecular $\cdot$ $D_{\bigtriangleup}$ supramolecular $\cdot$ $D_{\infty}$ temporal $\cdot$ hybrid sets $\cdot$ \textbf{$D_{\text{holo}}$} holographic (bulk-boundary correspondence, AdS/CFT) \\
\textbf{Topology ($T$)} & Internal connectivity pattern of the synthon's minimal motif & $T_{\bowtie}$ cyclic $\cdot$ $T_{\ggg}$ chain $\cdot$ $T_{\square}$ hub/node $\cdot$ $T_{\square\square}$ cage $\cdot$ $T_{\cup}$ bowl $\cdot$ \$T\_\{ \\
\textbf{Recognition Mode ($R$)} & Physical mechanism enabling reliable constraint propagation & $R_{\subseteq}$ covalent $\cdot$ $R_{\supseteq}$ non-covalent $\cdot$ $R_{\ddagger}$ catalytic $\cdot$ $R_{\Leftrightarrow}$ mechanical $\cdot$ covalent-dynamic \\
\textbf{Polarity ($P$)} & Directional character of the interaction & $P_{+}$ acceptor $\cdot$ $P_{-}$ donor $\cdot$ $P_{\pm}^{\text{sym}}$ self-complementary symmetric $\cdot$ $P_{\pm}^{\psi}$ self-complementary pseudosymmetric $\cdot$ $P_{+-}$ directional donor-acceptor \\
\textbf{Fidelity ($F$)} & Thermodynamic reliability of the synthon, anchored to $\xi_{CP}$ & $F_{\hbar}$ high ($\xi_{CP} \leq 8.5$ nats) $\cdot$ $F_{\eth}$ medium (8.5--11.0 nats) $\cdot$ $F_{\ell}$ low ($> 11.0$ nats) \\
\textbf{Kinetic Character ($K$)} & Activation barrier and pathway multiplicity for constraint propagation & $K_{\text{fast}}$ ($\Delta G^{\ddagger} < 60$ kJ/mol) $\cdot$ $K_{\text{mod}}$ (60--100 kJ/mol) $\cdot$ $K_{\text{slow}}$ ($> 100$ kJ/mol) $\cdot$ $K_{\text{trap}}$ (pathway multiplicity) $\cdot$ \textbf{$K_{\text{MBL}}$} (many-body localization --- disorder-frozen, not barrier-limited) \\
\textbf{Granularity ($G$)} & Scale of control exerted by the synthon & $G_{\beth}$ local $\cdot$ $G_{\gimel}$ mesoscale $\cdot$ $G_{\aleph}$ global/network \\
\textbf{Interaction Grammar ($\Gamma$)} & Logic governing partner selection & $\Gamma_{\wedge}$ AND $\cdot$ $\Gamma_{\vee}$ OR $\cdot$ $\Gamma_{\to}$ SEQUENTIAL $\cdot$ \textbf{$\Gamma_{\downarrow}$ DISSIPATIVE} (irreversible loss); each qualified by tier: SPECIFIC $\cdot$ SELECTIVE $\cdot$ BROAD $\cdot$ \textbf{QUANTUM} (superposition-preserving) \\
\textbf{Criticality Phase ($\Phi$)} & Phase of the synthon relative to the $G$--$D$ criticality locus & $\Phi_{\text{sub}}$ subcritical $\cdot$ $\Phi_c$ critical $\cdot$ $\Phi_{\text{super}}$ supercritical \\
\textbf{Stoichiometry ($S$)} & Valency ratio of the recognition event & $1:1$ homodimeric $\cdot$ $n:n$ symmetric multimeric $\cdot$ $n:m$ asymmetric; constrains $T_{\bowtie}$ topology and $P$ polarity \\
\textbf{Topological Protection Index ($\Omega$)} & Symmetry class of topological protection (quantum extension) & $\Omega_0$ trivial (classical) $\cdot$ $\Omega_Z$ winding number (Kitaev chain, SSH) $\cdot$ $\Omega_{Z_2}$ (topological insulators) $\cdot$ $\Omega_C$ Chern number (quantum Hall) $\cdot$ \textbf{$\Omega_{NA}$} non-abelian anyons ($\nu$=5/2 FQH, Kitaev honeycomb) \\
\bottomrule
\end{tabularx}
\end{table}

Set notation accommodates hybrid systems: $D = \{D_{\bigtriangleup}, D_{\infty}\}$ denotes a MOF-catalyst architecture operating in both spatial and temporal dimensions simultaneously. $\Phi$ defaults to $\Phi_{\text{sub}}$ for entries where criticality has not been established.

\textbf{$D_{\text{holo}}$ (holographic, v0.4.4):} A twelfth dimensionality value for systems where bulk degrees of freedom are encoded on a lower-dimensional boundary --- the AdS/CFT correspondence and its generalisations. The AdS/CFT boundary encoding previously required a hybrid $D = \{D_{\bigtriangleup}, D_{\infty}\}$ approximation (see §XVII, Result 8); $D_{\text{holo}}$ makes the bulk-boundary correspondence a first-class primitive rather than a proxy. Key algebraic property: any transition from $D_{\text{holo}}$ to any bulk phase ($D_{\bigtriangleup}$, $D_{\wedge}$, etc.) is a 1st-order morphism with infinite primitive cost --- the bulk-boundary map is not a continuous HotSwap. \texttt{ads\_cft\_boundary} synthon: $\langle D_{\text{holo}}; T_{\in}; R_{\ddagger}; P_{\pm}^{\psi}; F_{\eth}; K_{\text{mod}}; G_{\aleph}; \Gamma_{\wedge}(\text{SELECTIVE}); \Phi_c \rangle$.

\subsubsection{Primitive notes}

\textbf{D, T, R} are well-specified, genuinely orthogonal, and participate cleanly in composition rules. Set notation for $D$ handles hybrid systems without multiplying primitive count. The $R_{\Leftrightarrow}$ mechanical bond mode is distinguished by a discontinuous steric-cliff barrier profile rather than a Morse-like continuous curve --- the subject of the planned Transformation \#8.

\textbf{$T$ topology values} span the full range from acyclic to fully enclosed three-dimensional frameworks. The twelve named values are: $T_{\bowtie}$ (cyclic), $T_{\ggg}$ (chain), $T_{\square}$ (hub/node), $T_{\square\square}$ (cage), $T_{\cup}$ (bowl), $T_{|}$ (linear), $T_{\perp}$ (branched), $T_{\in}$ (network, generic), $T_{\in}(\text{hex})$ (hexagonal), $T_{\in}(\text{mixed})$ (mixed-ring), $T_{\in}(\times 2)$ (interpenetrating), and $T_{\in}(\text{sym})$ (centrosymmetric). The three values with full grounding axioms are $T_{\bowtie}$, $T_{\square}$, and $T_{\square\square}$; $T_{\cup}$ carries its own partial-enclosure axiom (see below); the three acyclic connectivity classes $T_{|}$ (unbranched linear chain, no junction nodes), $T_{\perp}$ (branched acyclic, one or more junction nodes), and $T_{\in}$ (multiply-connected network, cycles permitted without full enclosure) extend the topology basis to cover polymers, dendrimers, and extended coordination networks. \textbf{The cage/bowl/hub triad encodes degrees of enclosure:} $T_{\square}$ (hub/node) coordinates partners with no enclosure; $T_{\cup}$ (bowl) provides a concave cavity with a single open portal --- guest enters and exits freely, $K_{\text{fast}}$ is the default; $T_{\square\square}$ (cage) encloses in all three spatial dimensions --- guest egress requires framework distortion, $K_{\text{slow}}$ or $K_{\text{trap}}$ is the default. This three-way distinction is non-conservative in HotSwap: $T_{\cup} \to T_{\square\square}$ changes the kinetic regime even when all other primitives match. The bowl topology was identified through catalog self-audit: 222 entries previously assigned $T_{\bowtie}$ (cyclic) were reclassified to $T_{\cup}$ when the audit revealed they described calixarene, resorcinarene, pillar{[}n{]}arene, and related open-cavity hosts --- none of which have a closing bond. The formal entry for $T_{\square\square}$ appears below; Axiom 7's closing-bond requirement is extended to a closing-face requirement for cage synthons (see Axiom 7).

\textbf{P} distinguishes $P_{\pm}^{\text{sym}}$ (identical donor and acceptor faces, true homodimer) from $P_{\pm}^{\psi}$ (geometrically cyclic but electronically asymmetric, as in the carboxylic acid dimer with inequivalent C=O and O--H faces). The acid homodimer is $P_{\pm}^{\psi}$; a symmetric urea-type dimer is $P_{\pm}^{\text{sym}}$. Axiom 1's fidelity floor applies to both subclasses.

\textbf{F} is the thermodynamic fidelity metric: how strongly the correct product is favored at equilibrium, quantified by $\xi_{CP}$ in nats. Tier boundaries are anchored at $\xi_{CP} \leq 8.5$ nats ($F_{\hbar}$), 8.5--11.0 nats ($F_{\eth}$), and $> 11.0$ nats ($F_{\ell}$). These boundaries are derived from the calibrated I(bits) pipeline (v2.2); the previous heuristic values (9.0/11.5 nats) were based on an undercalibrated I estimate of 4--6 bits. In supramolecular host--guest systems the tier boundary between $F_{\eth}$ and $F_{\ell}$ corresponds approximately to $K_a \sim 10^7$ M$^{-1}$ ($\Delta G \approx -40$ kJ/mol at 298 K): complexes below this threshold (e.g., CB{[}7{]}$\cdot$DABCO, $K_a \approx 2 \times 10^5$ M$^{-1}$) are assigned $F_{\ell}$; those above (e.g., CB{[}7{]}$\cdot$adamantane, $K_a \approx 4 \times 10^8$ M$^{-1}$) are $F_{\eth}$; and ultrastrong binders (CB{[}7{]}$\cdot$ferrocene, $K_a \approx 3 \times 10^{12}$ M$^{-1}$) are $F_{\hbar}$. This three-tier resolution is required to correctly predict the competitive displacement hierarchy: the $F$-floor hard constraint in HotSwap (§2.1, SYNTHONIC\_HOTSWAP.md) is verified against the CB{[}7{]} series (Kim 2001; Assaf \& Nau 2015) in the validation table (§XII). A synthon with high $F$ but a prohibitive $\Delta G^{\ddagger}$ is operationally inaccessible --- a condition encoded by $K$, not $F$. The two primitives are orthogonal and must not be conflated.

\textbf{K} encodes the kinetic character of constraint propagation, independent of $F$: the carboxylic acid homodimer is $F_{\hbar}$, $K_{\text{fast}}$; the gas-phase imine condensation proxy is $F_{\eth}$, $K_{\text{slow}}$; the aqueous imine is $F_{\eth}$, $K_{\text{mod}}$ --- the $K$ assignments are different while $F$ is identical in both imine cases. $K_{\text{trap}}$ encodes high pathway multiplicity, where kinetic products diverge from thermodynamic products, a condition not representable as a barrier height alone.

\textbf{G and D} are independent for the vast majority of synthons. Changing $G$ without changing $D$ produces different constraint propagation range predictions; changing $D$ without changing $G$ produces different assembly topology predictions. At the criticality locus, $G$ degenerates with $D$ --- encoded as $\Phi_c$ (see §X).

\textbf{$\Gamma$} encodes both the \emph{logic} of partner selection (three operators) and the \emph{breadth} of the partner set (three tiers):

\begin{itemize}
    \item $\Gamma_{\wedge}$ (AND): all partners required simultaneously --- ternary complex, cooperative allosteric system
    \item $\Gamma_{\vee}$ (OR): any one of a defined partner set suffices --- promiscuous binder, degenerate recognition
    \item $\Gamma_{\to}$ (SEQUENTIAL): partner A required before B can bind --- template-directed assembly, allosteric activation
\end{itemize}

Each operator is qualified by a tier: SPECIFIC (one partner), SELECTIVE (2--10), BROAD (10--100+). A full encoding takes the form $\Gamma_{\wedge}(\text{SELECTIVE})$.

\textbf{$\Phi$} encodes the phase of the synthon relative to the criticality locus. $\Phi_{\text{sub}}$ is the default for synthons where $G$ and $D$ are demonstrably independent. $\Phi_c$ marks the criticality locus: $\xi \to \infty$, scale-free behavior, $G/D$ degenerate, recursively self-encoding tuple (Axiom 5). $\Phi_{\text{super}}$ marks the post-assembly state where synthon identity is absorbed into the assembled material. Empirical anchoring of $\Phi_c$ remains an active research objective (see §X).

\textbf{$S$} encodes the stoichiometric valency of the recognition event: the ratio of partners required to form the minimal functional motif. $S$ is not redundant with $T$ or $\Gamma$: topology encodes the shape of the connectivity, grammar encodes the logic of partner selection, but neither encodes the numerical ratio. A $T_{\bowtie}$ cyclic dimer may be $1:1$ (homodimer, $P_{\pm}$) or $n:m$ (asymmetric host-guest, $P_{+-}$ with $\Gamma_{\vee}$). \textbf{Hard constraints:} $T_{\bowtie} + S = 1:1$ requires $P_{\pm}$ (asymmetric polarity with symmetric stoichiometry is a contradiction); $T_{\bowtie} + S = n:m$ ($n \neq m$) requires $\Gamma_{\vee}(\text{BROAD})$ or network topology. $S$ defaults to $1:1$ for $T_{\bowtie}$ entries where $P_{\pm}$ is present and no other stoichiometry is specified; enforced by Pass 4 audit. Weight in similarity scoring: 0.08 (contributing \textasciitilde{}6\% of total), raised to 0.12 with \texttt{--stoichiometry-aware} for valency-sensitive queries.

\subsubsection{$T_{\square\square}$ --- Cage Topology: Formal Entry}

A synthon whose minimal motif forms a three-dimensionally closed polyhedral framework, fully encapsulating a guest within a defined interior volume. Distinguished from $T_{\square}$ (hub/node) by the requirement of \emph{closure in all three spatial dimensions} --- a hub/node coordinates but does not enclose; a cage encloses and thereby sequesters.

\textbf{Examples:} Cryptand--metal complexes, Fujita-type $\text{Pd}_{12}\text{L}_{24}$ spheres, covalent organic cages (COCs), cucurbituril hosts, carcerands/hemicarcerands, metal-organic polyhedra (MOPs).

\textbf{Primitive interactions for $T_{\square\square}$ systems:}

\begin{table}[htbp]
\centering
\small
\rowcolors{1}{white}{white}
\begin{tabularx}{\textwidth}{>{\centering\arraybackslash}X >{\centering\arraybackslash}X >{\centering\arraybackslash}X}
\toprule
\rowcolor{gray!25}\textbf{Primitive} & \textbf{Typical value} & \textbf{Rationale} \\
\midrule
\rowcolors{1}{white}{gray!10}
\textbf{$R$} & $R_{\supseteq}$ or $R_{\subseteq}$ & Non-covalent (self-assembled cages) or covalent (COCs); $R_{\Leftrightarrow}$ excluded --- mechanical bonding requires a thread, not encapsulation \\
\textbf{$K$} & $K_{\text{mod}}$ or $K_{\text{slow}}$ & Guest exchange requires partial cage disassembly or portal opening; $K_{\text{fast}}$ only for hemicarcerands with defined portals larger than the guest \\
\textbf{$G$} & $G_{\gimel}$ or $G_{\aleph}$ & Mesoscale for discrete cages; $G_{\aleph}$ when cages tessellate into a lattice \\
\textbf{$F$} & $F_{\hbar}$ expected & Encapsulation enforces high geometric fidelity; $F_{\ell}$ in a $T_{\square\square}$ system is an Axiom 1 analogue violation (cage closure is the $T_{\square\square}$ equivalent of cyclic cooperativity) \\
\textbf{$S$} & $n{:}1$ (cage : guest) & Stoichiometry is typically $1:1$ or $1:n$ for multi-guest cages; host:guest ratio is the canonical $S$ assignment \\
\textbf{$\xi_{CP}$ baseline} & \textasciitilde{}8.5--9.2 nats & Higher than $T_{\square}$ ($\approx 7.8$ nats) due to stricter kinetic requirements on guest exchange; lower than $T_{\bowtie}$ + $R_{\ddagger}$ systems \\
\bottomrule
\end{tabularx}
\end{table}

\textbf{Axiom 7 grounding for $T_{\square\square}$:} The closing-bond requirement of Axiom 7 generalizes to a \emph{closing face} --- the final panelling or coordination event that seals the third dimension. Grounding text must contain an assembly indicator (self-assemble, condense, panelling, cage-close, cyclize) and a closure indicator (encloses, sequesters, encapsulates, portal, aperture, face-capped). Computationally enforced as Pass 2b.

\textbf{Distinction from $T_{\square}$ (hub/node):}

\begin{table}[htbp]
\centering
\small
\rowcolors{1}{white}{white}
\begin{tabularx}{\textwidth}{>{\centering\arraybackslash}X >{\centering\arraybackslash}X >{\centering\arraybackslash}X}
\toprule
\rowcolor{gray!25}\textbf{Feature} & \textbf{$T_{\square}$ hub/node} & \textbf{$T_{\square\square}$ cage} \\
\midrule
\rowcolors{1}{white}{gray!10}
Guest access & Open coordination sphere & Fully enclosed interior \\
Dimensional closure & 2D or partial & 3D complete \\
$K$ for guest exchange & $K_{\text{fast}}$--$K_{\text{mod}}$ & $K_{\text{mod}}$--$K_{\text{slow}}$ \\
Typical $G$ & $G_{\beth}$--$G_{\aleph}$ & $G_{\gimel}$--$G_{\aleph}$ \\
Axiom 7 closing unit & Named bond & Named closing face \\
\bottomrule
\end{tabularx}
\end{table}

\textbf{K-compatibility note for HotSwap:} Swapping $T_{\square} \to T_{\square\square}$ is not a conservative swap even when $D$, $R$, and $S$ match, because the cage imposes a kinetic barrier on guest exchange absent in the hub/node system. The K-compatibility check (§8.0, SYNTHONIC\_HOTSWAP.md §2.2) must count pathways through the cage portal explicitly --- hemicarcerands with narrow portals are $K_{\text{trap}}$ candidates regardless of their nominal $\Delta G^{\ddagger}$.

\textbf{Code:} \texttt{Topology.CAGE} (\texttt{''T\_cage''} / \texttt{''T\_\textbackslash{}square\textbackslash{}square''}). Complexity score: 4. Topology cooperativity factor: 1.8 (intermediate between $T_{\bowtie}$ at 1.5 and $T_{\square}$ at 2.0 --- cage closure produces cooperativity but the kinetic ceiling on guest exchange attenuates the effective amplification vs. a hub's open coordination sphere).

\begin{center}
\rule{0.5\textwidth}{0.4pt}
\end{center}

\subsection{III. The Kinetic Primitive $K$: Separation of Thermodynamic and Kinetic Fidelity}

Thermodynamic and kinetic fidelity are physically distinct quantities that can diverge substantially. A synthon may be thermodynamically favored ($F_{\hbar}$, low $\xi_{CP}$) yet kinetically inaccessible under synthesis conditions; conversely, a kinetically fast synthon may produce metastable assemblies that are thermodynamically disfavored. Encoding only $F$ leaves this distinction invisible.

$K$ encodes the kinetic character of constraint propagation:

\begin{itemize}
    \item $K_{\text{fast}}$: $\Delta G^{\ddagger} < 60$ kJ/mol; spontaneous on experimental timescales. Accessibility score: 0.95.
    \item $K_{\text{mod}}$: $\Delta G^{\ddagger} \approx 60$--100 kJ/mol; accessible with mild activation (heat, catalyst, solvent). Accessibility score: 0.70.
    \item $K_{\text{slow}}$: $\Delta G^{\ddagger} > 100$ kJ/mol; requires significant activation or is effectively irreversible on practical timescales. Accessibility score: 0.30.
    \item $K_{\text{trap}}$: pathway multiplicity high; kinetic products diverge substantially from thermodynamic products. Accessibility score: 0.50.
\end{itemize}

$K$ and $F$ are orthogonal by construction. All four combinations of $\{F_{\hbar}, F_{\ell}\} \times \{K_{\text{fast}}, K_{\text{slow}}\}$ have real chemical exemplars. The chelate node (Transformation \#3) is $F_{\hbar}$, $K_{\text{fast}}$ --- the bidentate geometry both raises the thermodynamic stability and accelerates re-formation, making $K$ a positive amplifier of effective fidelity in this case. The imine condensation (Transformation \#4) is $F_{\eth}$, $K_{\text{slow}}$ (gas phase) / $K_{\text{mod}}$ (aqueous) --- the same thermodynamic classification covering two operationally distinct systems.

The ten-primitive tuple:

\[\langle D \;;\; T \;;\; R \;;\; P \;;\; F \;;\; K \;;\; G \;;\; \Gamma \;;\; \Phi \;;\; S \rangle\]

Production-level $K$ validation requires full transition-state geometry optimization; single-point data are sufficient for directional tier assignment.

\begin{center}
\rule{0.5\textwidth}{0.4pt}
\end{center}

\subsection{IV. Composition Axioms: The Grammar's Production Rules}

A predictive grammar requires explicit production rules stated as falsifiable propositions. The following seven axioms are grounded in computational validation and operational use of the framework. They are working hypotheses, not proven theorems; their value lies in being falsifiable. Axioms 1--5 govern primitive composition. Axioms 6--7 govern physical grounding and are enforced at catalog registration time.

\textbf{Axiom 1 --- Cyclic closure amplifies fidelity (the $T_{\bowtie}$--$F$ rule).}
A synthon with $T_{\bowtie}$ and $P_{\pm}$ (either $P_{\pm}^{\text{sym}}$ or $P_{\pm}^{\psi}$) achieves $F \geq F_{\eth}$, provided $R_{\supseteq}$ or $R_{\subseteq}$. The cyclic motif creates a geometric constraint that prevents partial dissociation without full bond rupture, converting independent pairwise contacts into a single cooperative unit. \textbf{Hard constraint: $T_{\bowtie}$ + $P_{\pm}$ + $F_{\ell}$ is forbidden.} Falsified by: a cyclic self-complementary motif with $\xi_{CP} > 11.0$ nats.

\textbf{Axiom 2 --- Local grammar blocks network propagation (the $G_{\beth}$--$\Gamma$ barrier rule).}
A synthon with $G_{\beth}$ and $\Gamma_{\wedge}(\text{SPECIFIC})$ cannot propagate constraint beyond its immediate recognition pair, and cannot nucleate a network-scale ($G_{\aleph}$) assembly event without a change in at least one of $G$, $\Gamma$, or $T$. Falsified by: a single $G_{\beth}$/$\Gamma_{\wedge}(\text{SPECIFIC})$ synthon documented as the sole organizing element of a MOF, polymer, or oscillatory network.

\textbf{Axiom 3 --- Cooperative induction superlinearity signals a $G_{\beth} \to G_{\gimel}$ transition.}
When the induction component of $E_{\text{int}}$ (from SAPT decomposition) grows faster than linearly with recognition contact count, the system has crossed from $G_{\beth}$ to $G_{\gimel}$. The triple H-bond DAD$\cdot$ADA array (Transformation \#5) shows an induction ratio of 2.5--3.5$\times$ for single $\to$ triple contacts while electrostatics remain approximately additive. Any synthon array showing superlinear SAPT induction should be assigned $G_{\gimel}$ regardless of nominal contact count. Falsified by: an array with superlinear induction but demonstrably local (non-propagating) constraint.

\textbf{Axiom 4 --- Sequential grammar requires temporal or catalytic dimension (the $\Gamma_{\to}$--$D$ rule).}
$\Gamma_{\to}$ is physically realizable only if the synthon possesses $D_{\infty}$ or $R_{\ddagger}$, or both. Ordered recognition requires a mechanism by which partner A's binding changes the system state before B arrives --- a temporal cycle or catalytic transformation. Pure spatial synthons ($D_{\wedge}$ or $D_{\bigtriangleup}$ only, $R_{\supseteq}$) cannot encode sequential partner grammars without kinetic trapping, which is a $K$ effect, not a $\Gamma$ one. \textbf{Hard constraint at registration time.} Falsified by: a purely spatial, non-catalytic synthon with experimentally confirmed ordered partner binding in the thermodynamic limit.

\textbf{Axiom 5 --- Criticality contracts the primitive basis.}
A synthon at the criticality locus ($\Phi_c$: $G$ and $D$ degenerate, $\xi \to \infty$, scale-free behavior) requires fewer independent primitive values to specify its behavior. $G$ becomes redundant given $D$ at criticality because scale-invariance makes ''local vs. global'' an undefined distinction. A critical synthon's behavior at the molecular scale fully predicts its supramolecular and temporal behavior by the same production rules. Falsified by: a scale-free synthon requiring different $G$ assignments at molecular and supramolecular scales to correctly predict behavior at each.

\textbf{Axiom 6 --- Temporal grounding: $D_{\infty}$ requires either a named closed cycle or a continuously supplied dissipative flux.}
Any synthon assigned $D_{\infty}$ must satisfy one of two modes:

\begin{itemize}
    \item \textbf{Discrete reset (type: \texttt{''discrete''})} --- a stoichiometric or triggered reset step returns the system to its initial state after each turnover. Examples: proline-catalysed aldol (iminium hydrolysis $\to$ free proline), photo-switch (light-off thermal back-relaxation), rotaxane dethreading (stopper removal). Grounding requires \texttt{cycle\_steps} $\geq$ 2, or structured fields \texttt{initial\_state}, \texttt{transformation}, \texttt{work\_performed}, \texttt{reset\_mechanism}.
    \item \textbf{Continuous dissipative flux (type: \texttt{''continuous''})} --- an open, far-from-equilibrium system maintained by ongoing supply of a thermodynamic gradient (fuel inflow, photon flux, redox potential). No discrete reset step exists; the cycle terminates only when external supply is exhausted. Examples: Soai autocatalytic amplification (continuous iPr$_2$Zn + aldehyde supply), Belousov-Zhabotinsky oscillations (continuous redox reagent influx), metabolic loops (continuous electron/proton flux). Grounding requires \texttt{driving\_gradient.description} and \texttt{driving\_gradient.coupling} fields.
\end{itemize}

A directed transformation that has neither a recoverable discrete reset nor an identified sustaining flux is not a temporal synthon. \textbf{Enforced computationally (audit Pass 1):} $D_{\infty}$ entries are checked against the structured \texttt{synthon.grounding{[}''reset''{]}} block first; if absent, keyword scan detects reset/process indicators; entries failing both are flagged for mandatory review. The \texttt{grounding{[}''reset''{]}{[}''type''{]}} field controls which validation path is applied.

\textbf{Axiom 7 --- Cyclic and cage topology grounding: $T_{\bowtie}$ requires a named closing bond; $T_{\square\square}$ requires a named closing face.}
Any synthon assigned $T_{\bowtie}$ must have a physically specified closing interaction. A linear or chain-like system (cumulene, rod, allene, extended polymer) cannot be assigned $T_{\bowtie}$; the topological error propagates into every downstream analogy search since topology is a primary similarity dimension. \textbf{Enforced computationally:} $T_{\bowtie}$ without grounding text containing a closing-bond indicator (hydrogen bond, ring, macrocycle, chelate, rotaxane, catenane, dimer, C--C bond) is flagged. Linear/chain/axial/rod keywords in a $T_{\bowtie}$ description trigger mandatory review. This axiom is the basis for audit Pass 2.

\textbf{Extension to $T_{\square\square}$ (cage topology):} Any synthon assigned $T_{\square\square}$ must have a physically specified \emph{closing face} --- the final panelling, coordination, or condensation event that seals the third spatial dimension. Grounding text must contain: (a) an \textbf{assembly indicator} (self-assemble, condense, panelling, cage-close, cyclize) and (b) a \textbf{closure indicator} (encloses, sequesters, encapsulates, portal, aperture, face-capped). A $T_{\square\square}$ entry without an identified closing face is flagged \texttt{grounding\_status = unverified} identically to an unclosed $T_{\bowtie}$. \textbf{Enforced computationally (Pass 2b):} $T_{\square\square}$ without closing-face language triggers mandatory review. Falsified by: a three-dimensionally enclosed polyhedral host system correctly grounded with only hub/node language.

\textbf{Axiom 1 --- Quantum boundary condition.} Axiom 1 is a \emph{classical} axiom: it assumes cyclic self-complementary recognition is cooperative, amplifying fidelity above the $F_{\ell}$ floor. Quantum entanglement violates this assumption in a specific, diagnosable way. An entangled spin singlet is correctly $T_{\bowtie}$ (two particles forming a closed loop of mutual constraint) and $P_{\pm}^{\text{sym}}$ (symmetric partners), yet the LLM initially assigned $F_{\ell}$ --- triggering an Axiom 1 violation that persisted through 3 refinement iterations. The resolution reveals a hidden conflation:

\begin{itemize}
    \item The LLM assigned $F_{\ell}$ because quantum entanglement \emph{cannot transmit classical information} (no-communication theorem).
    \item But $F$ in the framework measures \emph{reliability of the constraint}, not Shannon channel capacity.
    \item For a spin singlet: measuring one spin guarantees the other is anti-parallel with 100\% reliability. The constraint fires perfectly --- $F_{\hbar}$, not $F_{\ell}$.
\end{itemize}

Corrected encoding: $\langle D_{\wedge}; T_{\bowtie}; R_{\supseteq}; P_{\pm}^{\text{sym}}; F_{\hbar}; K_{\text{trap}}; G_{\aleph}; \Gamma_{\wedge}(\text{SELECTIVE}); \Phi_{\text{sub}} \rangle$. Axiom 1 satisfied. The $K_{\text{trap}}$ (permanent, non-exchangeable) and $G_{\aleph}$ (non-local correlation propagates globally --- Bell non-locality) distinguish it from all classical $T_{\bowtie}$ systems, which have $K_{\text{fast}}$ or $K_{\text{mod}}$ and $G_{\beth}$ or $G_{\gimel}$. \textbf{Axiom 1 is a classical production rule; its violation by a quantum system is a diagnostic signal, not a falsification of the axiom.} The framework correctly identifies the quantum boundary by the pattern: $T_{\bowtie} + P_{\pm} + F_{\ell}$ is not just forbidden --- it is \emph{unreachable} in the classical domain, but reachable in the quantum domain where $F$ and Shannon capacity decouple.

Expanding the axiom set --- particularly axioms connecting $K_{\text{trap}}$ to selectivity outcomes and $\Gamma_{\wedge}$ behavior in ternary complexes --- is an active theoretical task.

\begin{center}
\rule{0.5\textwidth}{0.4pt}
\end{center}

\subsection{V. Theoretical Underpinnings: Constraint Propagation, Non-Equilibrium Thermodynamics, and the $\xi_{CP}$ Metric}

\subsubsection{Constraint propagation and dissipative structures}

The Unified Synthonicon is a practical instantiation of two interconnected theoretical frameworks. The first is the mathematics of constraint propagation, where a synthon functions as a local constraint --- its combination of dimensionality, topology, and recognition mode defines rules that limit the possible configurations of surrounding components, reducing the system's degrees of freedom and steering its evolution toward a lower-entropy target state. The second is the thermodynamics of dissipative structures, which provides the theoretical basis for temporal synthons: systems maintaining organized states by continuously consuming energy and exporting entropy, with the Kolmogorov-Sinai entropy (rate of phase-space information loss) directly proportional to the thermodynamic entropy production rate.

\subsubsection{$\eta_{CP}$ and $\xi_{CP}$: constraint propagation efficiency and inefficiency index}

\[\eta_{CP} = \frac{I \times F}{\Delta G / E_{\text{bit}}^{\text{molar}}}\]

where $I$ = information gain (bits) from configurational restriction, $F$ = fidelity (0--1) normalized to the strongest system in the comparison set, $\Delta G$ = free energy cost (kJ/mol), and $E_{\text{bit}}^{\text{molar}} \approx 1.72 \times 10^{-3}$ kJ/mol/bit (Landauer cost per bit at 298 K). $\eta_{CP} = 1$ represents perfect Landauer efficiency.

\[\xi_{CP} = -\ln(\eta_{CP}) \quad \text{(nats)}\]

$\xi_{CP}$ functions as a common currency across domains, enabling cross-domain comparison. As of v2.2, $I$ is derived from a first-principles degree-of-freedom (DOF) counting pipeline rather than a heuristic. The pipeline decomposes information content into $I_{\text{recognition}}$ (selectivity DOFs: contact distance, H-bond angles, torsional conformers), $I_{\text{orientation}}$ (rigid-body overhead), and $I_{\text{net}} = I_{\text{recognition}} - 0.3 \times I_{\text{orientation}}$. Live calibration (\texttt{syncon info-bits --calibrate}, v0.4.2) for three reference systems:

\begin{table}[htbp]
\centering
\small
\rowcolors{1}{white}{white}
\begin{tabularx}{\textwidth}{>{\centering\arraybackslash}X >{\centering\arraybackslash}X >{\centering\arraybackslash}X >{\centering\arraybackslash}X >{\centering\arraybackslash}X >{\centering\arraybackslash}X}
\toprule
\rowcolor{gray!25}\textbf{Target} & \textbf{$I_\text{rec}$} & \textbf{$I_\text{net}$} & \textbf{$I_\text{+solvent}$} & \textbf{Expected range} & \textbf{Status} \\
\midrule
\rowcolors{1}{white}{gray!10}
Carboxylic acid homodimer & 9.39 bits & 8.02 bits & 13.98 bits & 9.0--10.5 bits &  \\
Triple H-bond DAD$\cdot$ADA & 16.57 bits & 15.19 bits & 21.15 bits & 14.0--18.0 bits &  \\
Proline aldol cycle & 7.98 bits & 6.61 bits & 12.57 bits & 6.0--9.0 bits &  \\
\bottomrule
\end{tabularx}
\end{table}

Use $I_\text{rec}$ for propagation estimates; $I_\text{net}$ for selectivity-purified comparisons; $I_\text{+solvent}$ for thermodynamic budgeting. Calibrated domain-dependent ranges: $I_\text{rec}$ \textbf{8--17 bits}, $I_\text{net}$ \textbf{7--15 bits}, $I_\text{+solvent}$ \textbf{13--21 bits} (was 4--6 bit heuristic). Solvent correction adds \textasciitilde{}4.6--5.6 bits uniformly across the three reference systems. Cooperativity scaling: \textasciitilde{}4--5 bits/contact confirmed across 2--4 H-bond arrays. The primary remaining open task is anharmonic corrections to the harmonic well approximation and full QM benchmarking of the $\sigma$-hole angle window.

\textbf{$\xi_{CP}$ table (298 K, calibrated I values):}

\begin{table}[htbp]
\centering
\small
\rowcolors{1}{white}{white}
\begin{tabularx}{\textwidth}{>{\centering\arraybackslash}X >{\centering\arraybackslash}X >{\centering\arraybackslash}X >{\centering\arraybackslash}X >{\centering\arraybackslash}X >{\centering\arraybackslash}X >{\centering\arraybackslash}X}
\toprule
\rowcolor{gray!25}\textbf{System} & \textbf{$I_{\text{cal}}$ (bits)} & \textbf{$\xi_{CP}$ (nats)} & \textbf{Range} & \textbf{Tier} & \textbf{Domain} & \textbf{Primitive insight} \\
\midrule
\rowcolors{1}{white}{gray!10}
Acetic acid homodimer (AA:AA) & 9.4 & \textbf{6.66} & {[}6.56--6.77{]} & HIGH & Molecular & $F_{\hbar}$ reference, $P_{\pm}^{\psi}$; high I at low \$ \\
$\sigma$-hole dimer (I$\cdot$$\cdot$$\cdot$N) & 7.8 & \textbf{7.59} & {[}7.47--7.73{]} & HIGH & Molecular & Halogen-bond $F_{\hbar}$; narrow angle window $\pm$2.5$^{\circ}$ \\
Triple H-bond array (DAD$\cdot$ADA) & 16.0 & \textbf{7.65} & {[}7.59--7.72{]} & HIGH & Supramolec. & Cooperativity factor 1.25; Axiom 3 \\
$\sigma$-hole trimer & 11.2 & \textbf{8.40} & {[}8.31--8.49{]} & HIGH & Supramolec. & Network $\sigma$-hole \\
Acid--amide heterodimer & 8.2 & \textbf{8.19} & {[}8.07--8.32{]} & HIGH & Molecular & $P_{+-}$ directional; $\Delta G$ requires ITC confirmation \\
Adenine--thymine pair (A$\cdot$T) & \textasciitilde{}7.6 & \textasciitilde{}8.0--8.1 & --- & HIGH & Molecular & Canonical biological $F_{\hbar}$ \\
Zn--bpy chelate ($G_{\beth} \to G_{\aleph}$) & --- & \textasciitilde{}8.5 & --- & HIGH/MEDIUM & Supramolec. & Granularity amplification \\
Formamide homodimer & 6.8 & \textbf{8.70} & {[}8.56--8.86{]} & MEDIUM & Molecular & $P_{\pm}^{\text{sym}}$ weak; $\Delta G$ requires ITC confirmation \\
Proline aldol cycle (per cycle) & 7.5 & \textbf{9.21} & {[}9.09--9.36{]} & MEDIUM & Temporal & $D_{\infty}$, $F_{\eth}$, $K_{\text{mod}}$; ee prediction 70--85\% (exp. 74\%) \\
Zr-oxo balanced (triplet) & --- & 12.2 & --- & LOW & Prior proxy & Early framework reference \\
Zr-oxo tight singlet & --- & 15.0 & --- & LOW & Prior proxy & High geometric strain \\
Zr-oxo + toluene & --- & 16.8 & --- & LOW & Prior proxy & Confinement overhead \\
\bottomrule
\end{tabularx}
\end{table}

Tier boundaries: HIGH $\leq 8.5$ nats $\cdot$ MEDIUM 8.5--11.0 nats $\cdot$ LOW $> 11.0$ nats. Uncertainty ranges from $\pm$1 bit I calibration uncertainty.

The acetic acid homodimer (6.66 nats) occupies the lowest-waste position in the HIGH tier because it encodes high recognition information (I = 9.4 bits) at a modest free energy cost ($|\Delta G| = 12$ kJ/mol gas phase) --- the characteristic signature of an entropy-dominated association event where geometric pre-organisation does most of the selectivity work. The triple H-bond array (7.65 nats) and $\sigma$-hole dimer (7.59 nats) cluster at 7.6--7.7 nats; the cooperative gain in the triple array is captured in its elevated I rather than in a suppressed $\xi_{CP}$. The proline catalytic cycle (9.21 nats) demonstrates that a $D_{\infty}$ temporal synthon is thermodynamically competitive with spatial synthons on a per-bit basis. The HIGH-tier cluster spans 6.7--8.4 nats.

\begin{center}
\rule{0.5\textwidth}{0.4pt}
\end{center}

\subsection{VI. The Criticality Condition and the G--D Phase Diagram}

The convergence of $G$ and $D$ at the criticality locus generates the framework's most structurally novel prediction. In statistical physics, universality is the condition where a system's behavior near a phase transition is determined entirely by its symmetry and effective dimensionality, independent of microscopic details. At the criticality locus, a synthon loses its characteristic control length scale simultaneously in both $G$ and $D$: its influence propagates across all granularity levels, and its behavior at the molecular level predicts its supramolecular and temporal behavior by the same rules. This is the fractal self-encoding property described in Axiom 5.

\textbf{Primitive basis contraction.} At criticality, $G$ becomes degenerate with $D$ and cannot be independently assigned. The effective tuple contracts. Any synthon requiring independent $G$ and $D$ assignments is demonstrably not operating at criticality.

\textbf{Universality class membership.} Two synthons from entirely different chemical domains that share the same universality class at criticality will exhibit identical scaling exponents for their constraint propagation behavior --- even if their interaction energies, recognition modes, and topologies differ completely. This is the strongest cross-domain prediction the framework can make, testable via scaling analysis of $\xi_{CP}$ near cooperative phase transitions.

\textbf{The $\Phi$ phase primitive.} Criticality is encoded as a phase of the synthon itself --- $\Phi_c$ --- rather than a geometric point in the $G$-$D$ plane. The question of whether $\Phi$ represents a genuinely independent dimension or is derivable from existing primitives at criticality (Axiom 5 predicts derivability: $\Phi_c$ is signaled by $G/D$ degeneracy) requires at least one empirical anchor system for resolution.

\textbf{Candidate anchor systems.} The leading theoretical candidate remains the Varma quantum XY model, whose correlation function factorizes as a product of spatial and temporal parts --- a $G/D$ criticality locus signature --- with $\xi_r = \ln \xi_\tau$ (temporal correlation length growing exponentially relative to spatial), precisely the self-similar recursive encoding behavior Axiom 5 predicts. Encoding the Varma quantum critical point as a synthon is the first test of whether Axiom 5's criticality predictions match a rigorously derived condensed-matter result. Additional candidates include cooperative H-bond arrays near their percolation threshold, MOF systems near a structural phase transition, and the Transformation \#8 rotaxane dethreading scan. Among chemical systems, the \textbf{Soai autocatalytic reaction} (pyrimidyl-alkanol self-amplification, Soai 1995) is the highest-confidence experimental candidate: it encodes as $\langle D_{\infty}; T_{\bowtie}; R_{\ddagger}; P_{+-}; F_{\hbar}; K_{\text{mod}}; G_{\gimel}; \Gamma_{\to}(\text{SPECIFIC}); \Phi_{\text{sub}}; 1{:}1 \rangle$ and yields a Varma candidacy score of \textbf{0.920} (approaching $\Phi_c$) due to the Frank-model classical bifurcation fingerprint ($D_{\infty} + T_{\bowtie} + P_{\text{directional}} + F_{\hbar}$ co-present; pitchfork bifurcation at ee = 0; Gridnev 2010, Shibata 2009). In contrast, the proline-aldol cycle scores \textbf{0.380} ($\Phi_{\text{sub}}$): $\xi_r = 6.2$ (60 Å pair correlation), $\xi_\tau \approx 1.8 \times 10^{14}$ ($\omega_c = 10^{12}$ s$^{-1}$, solvent relaxation), ratio $= 0.189$ --- well below the critical ratio of 1.0 --- consistent with subcritical oscillatory chemistry without symmetry-breaking cooperation (Blackmond RPKA 2004; Houk/List DFT 2004). The DB24C8 pseudorotaxane scores \textbf{0.461} (steric-cliff proxy, Factor 6; Groppi 2020).

\subsubsection{Quantitative degeneracy measurement (v2.2)}

The Varma probe now provides quantitative rather than binary degeneracy classification. Two quantities are computed:

\textbf{Dynamic exponent $z_{\text{eff}}$.} Defined as $z_{\text{eff}}(\omega) = \ln \xi_\tau(\omega) / \ln \xi_r(\omega)$. In conventional criticality $z$ is a finite constant ($z \approx 1$--2). In the Varma QXY class, $\xi_\tau = \exp(\xi_r)$, so $z_{\text{eff}} = \xi_r / \ln \xi_r$ --- a function that diverges without bound as $\xi_r$ increases. This divergence is the distinguishing signature of logarithmic $G/D$ degeneracy. Reference values:

\begin{table}[htbp]
\centering
\small
\rowcolors{1}{white}{white}
\begin{tabularx}{\textwidth}{>{\centering\arraybackslash}X >{\centering\arraybackslash}X >{\centering\arraybackslash}X}
\toprule
\rowcolor{gray!25}\textbf{System} & \textbf{$z_{\text{eff}}$} & \textbf{Type} \\
\midrule
\rowcolors{1}{white}{gray!10}
Varma QXY ($\xi_r = 13.8$, $\xi_\tau = 10^6$) & 5.26, growing & logarithmic \\
$\xi_r = 2 \to 20$ sweep (Varma) & 2.89 $\to$ 6.68 (+2.3$\times$) & diverging  \\
2D H-bond percolation threshold & \textbf{1.330} & power-law ($z = 4/3$, exact) \\
\bottomrule
\end{tabularx}
\end{table}

The 2D percolation result exactly recovers the known theoretical value ($\nu = 4/3$), independently verifying the formula.

\textbf{Degeneracy strength score} ($s \in [0, 1]$). A composite four-component score classifying the regime:

\begin{table}[htbp]
\centering
\small
\rowcolors{1}{white}{white}
\begin{tabularx}{\textwidth}{>{\centering\arraybackslash}X >{\centering\arraybackslash}X >{\centering\arraybackslash}X}
\toprule
\rowcolor{gray!25}\textbf{Score range} & \textbf{Tier} & \textbf{Physical meaning} \\
\midrule
\rowcolors{1}{white}{gray!10}
0.00--0.30 & \texttt{none} & $G$ and $D$ fully independent \\
0.30--0.60 & \texttt{logarithmic} & Varma QXY --- $\xi_r \approx \ln \xi_\tau$ (weak Axiom 5) \\
0.60--0.85 & \texttt{power-law} & Conventional QCP --- $\xi_r \sim \xi_\tau^{1/z}$, $z$ finite \\
0.85--1.00 & \texttt{collapse} & Direct $G/D$ identity (strong Axiom 5) \\
\bottomrule
\end{tabularx}
\end{table}

These distinctions are actionable for the programmable matter roadmap: systems with $s \approx 0.60$--0.70 are near the logarithmic--power-law boundary and are the most tractable experimental targets for tuning toward collapse degeneracy. The score is computed by \texttt{degeneracy\_strength()} in \texttt{varma\_probe.py}; batch rankings are produced by \texttt{syncon criticality-probe --batch --export-candidates}.

\textbf{Scoring factors in \texttt{score\_phi\_c\_candidacy()}.} The probe evaluates eight independent heuristic factors, each contributing a weighted partial score:

\begin{itemize}
    \item \textbf{Factors 1--5 (Varma QXY structural heuristics):} Temporal dimensionality $D_{\infty}$ (w=0.35), catalytic recognition $R_{\ddagger}$ with moderate kinetics (w=0.25), mesoscale granularity $G_{\gimel}$ (w=0.20), and related structural proxies for quantum XY universality.
    \item \textbf{Factor 6 --- Steric-cliff proxy (w=0.65):} Reads \texttt{proxy\_degeneracy\_strength} from grounding metadata \texttt{phi\_c\_candidacy}. Fires when proxy score $\geq$ 0.50. This scores the mechanical-bond barrier-sharpness mechanism ($R_{\Leftrightarrow}$ systems such as DB24C8; universality class: steric-cliff proxy, distinct from Varma QXY temporal mechanism).
    \item \textbf{Factor 7 --- Frank-model classical bifurcation (w=0.25):} Fires when all four co-requisites are present simultaneously: $D_{\infty}$ (temporal), $T_{\bowtie}$ (cyclic), $P_{\text{directional}}$ (donor-acceptor), $F_{\hbar}$ (high fidelity). This identifies the Frank 1953 pitchfork bifurcation at ee = 0 --- a \emph{classical} symmetry-breaking criticality mechanism in enantiospecific closed autocatalytic cycles, universality class distinct from Varma QXY.
    \item \textbf{Factor 8 --- Quantum criticality fingerprint (w=0.20):} Fires when $G_{\aleph}$ (global granularity) + $F_{\hbar}$ (high fidelity) + $K_{\text{trap}}$ (frozen kinetics) are all present \emph{without} $D_{\infty}$ (no temporal dimension). This pattern identifies zero-temperature quantum phase transitions driven by ground-state degeneracy rather than thermal fluctuations --- transverse-field Ising model at $h = h_c$, heavy fermion compounds (CeCu$_{6-x}$Au$_x$, YbRh$_2$Si$_2$), quantum dots at charge degeneracy. Universality class: TFI/heavy-fermion, distinct from all three classical mechanisms. Falsifiable prediction: susceptibility divergence $\chi(T \to 0) \sim T^{-\gamma}$.
\end{itemize}

\textbf{Three-mechanism discrimination (validated v0.3.3):}

\begin{table}[htbp]
\centering
\small
\rowcolors{1}{white}{white}
\begin{tabularx}{\textwidth}{>{\centering\arraybackslash}X >{\centering\arraybackslash}X >{\centering\arraybackslash}X >{\centering\arraybackslash}X >{\centering\arraybackslash}X}
\toprule
\rowcolor{gray!25}\textbf{System} & \textbf{Score} & \textbf{Candidacy} & \textbf{Mechanism} & \textbf{Key Factor} \\
\midrule
\rowcolors{1}{white}{gray!10}
Soai autocatalytic cycle & \textbf{0.920} & approaching $\Phi_c$ & Frank-model bifurcation & Factor 7 (w=0.25) + proxy 0.60 \\
DB24C8 pseudorotaxane & \textbf{0.461} & approaching $\Phi_c$ & Steric-cliff proxy & Factor 6 (proxy = 0.71) \\
Proline-aldol cycle & \textbf{0.380} & $\Phi_{\text{sub}}$ & None detected & Structural heuristics only \\
\bottomrule
\end{tabularx}
\end{table}

The catalog currently contains no confirmed $\Phi_c$ entries. The Soai result is the highest-confidence genuine candidate: probe score 0.920, $\xi_r / \ln \xi_\tau = 0.94$ (from Gridnev 2010 / Shibata 2009 data), and the full Frank-model co-requisite tuple satisfied. A direct experimental measurement of the spatial correlation length divergence near the bifurcation point (SAXS/DLS at varying initial ee) would convert this from structural candidacy to a confirmed $\Phi_c$ anchor. Entries carrying $\Phi_c$ assignments without grounding data, named closing bonds, or $\xi_r$/$\xi_\tau$ measurements remain excluded from analogies as contamination artifacts.

\begin{center}
\rule{0.5\textwidth}{0.4pt}
\end{center}

\subsection{VII. The Relational Substrate}

The eleven primitives $\langle D; T; R; P; F; K; G; \Gamma; \Phi; S; \Omega \rangle$ are, without exception, \textbf{relational operators} --- they describe constraints between entities, capacities for interaction, and partner-selection logic. None encodes a monadic property of a system in isolation.

\begin{itemize}
    \item \textbf{$F$ (Fidelity):} The thermodynamic reliability of \emph{constraint satisfaction} relative to a binding partner or competitor. There is no ''intrinsic $F$''; $F_\hbar > F_\eth$ is a statement about which entity prevails in a competitive context (V-1, CB{[}7{]} series: 6/6 directional predictions from the $F$ ordinal alone, without knowing the intrinsic chemistry of the guests).
    \item \textbf{$K$ (Kinetic):} A barrier to \emph{rearrangement} --- a transition between states implies at least two states and an environment that drives or resists the transition.
    \item \textbf{$\Gamma$ (Grammar):} Explicitly defined as partner selection logic. A $\Gamma$ assignment with no partner is undefined.
    \item \textbf{$\Omega$ (TopoIndex):} Topological protection \emph{against perturbations}. The protection is meaningless without an environment to be protected from.
    \item \textbf{$\Phi$ (Criticality):} Defined via \texttt{degeneracy\_strength} --- a count of equivalent global configurations --- which requires a system of multiple interacting parts.
\end{itemize}

This is not a philosophical preference. It is a \textbf{hard constraint of the type system}: you cannot assign $F$, $K$, $\Gamma$, or $\Omega$ without specifying an interaction context. A synthon tuple without a context is a description of interaction potential --- never a description of isolated being.

\subsubsection{Directed Constraints, Not Mere Relations}

The framework is more precisely described as \textbf{constraint-based} rather than purely relational. Pure mathematical relations are symmetric: if A relates to B then B relates to A. The SynthOmnicon algebra is systematically asymmetric:

\begin{itemize}
    \item \texttt{path(A -\textgreater{} B) != path(B -\textgreater{} A)} --- the $F$-floor ratchet is directional; higher-$F$ guests displace lower-$F$ guests, never the reverse (V-1: 3 allowed / 3 blocked, all confirmed).
    \item \texttt{tensor(kitaev\_chain, qubit)} assigns $F_\ell$ as the bottleneck at the qubit interface; the causal direction flows from the weaker partner.
    \item \texttt{lift(critical)} is one-way --- there is no \texttt{descend} operation.
\end{itemize}

This asymmetry is load-bearing: it is what makes predictions directional rather than merely topological. The framework encodes \textbf{directed, ordered constraints}, and the ordering is empirically anchored. This places the framework within the tradition of structural realism (Ladyman, French): the world's causal structure is relational but not symmetrically so.

\subsubsection{The Algebra Has No Unary Information Generators}

The compositional algebra (\texttt{meet}, \texttt{join}, \texttt{tensor}, \texttt{path}, \texttt{lift}, \texttt{pipeline}) contains no unary operations that generate new physical information from a single isolated synthon. Every operation requires at least one additional operand:

\begin{itemize}
    \item \texttt{tensor(s1, s2)} --- mutual information discount between \emph{two} systems (P-11)
    \item \texttt{meet(s1, s2)}, \texttt{join(s1, s2)} --- common sub/super-structures between \emph{pairs}
    \item \texttt{path(src, dst)} --- transformation \emph{between} states; requires two endpoints
    \item \texttt{lift(s, target)} --- migration to a target context (\texttt{temporal}, \texttt{spatial}, \texttt{critical}, \texttt{molecular})
\end{itemize}

\texttt{lift} is the apparent exception --- it takes one synthon --- but requires a named \emph{target context} as its second argument, and is blocked unless the synthon satisfies the relational gate \texttt{F \textgreater{}= F\_\textbackslash{}hbar}. The gate is enforced by the environment's fidelity floor, not by any intrinsic property of the synthon alone.

\textbf{Implication:} the algebra cannot process ''nothing but the object.'' It requires a second operand --- environment, partner, or target state --- to compute anything new.

\subsubsection{Axioms as Relational Well-Formedness Constraints}

The seven axioms encode relational requirements:

\begin{itemize}
    \item \textbf{Axiom 1 (Fidelity Floor):} $T_{\bowtie} + P_{\pm} \to F \geq F_{\eth}$. The cycle requires \emph{both halves}. Isolate one and fidelity collapses.
    \item \textbf{Axiom 5 (Criticality):} $\Phi_c$ is defined by $G$--$D$ degeneracy across the catalog --- a relational condition, not an intrinsic property of any single entry.
    \item \textbf{Axiom 6 (Temporal Grounding):} $D_{\infty}$ requires a reset mechanism, defined relative to a driving gradient or external boundary condition.
\end{itemize}

\textbf{Axiom 1 as quantum boundary detector (V-5)} provides the sharpest illustration. The entangled spin system was encoded with $T_{\bowtie} + P_{\pm}^{\text{sym}}$ but initially $F_\ell$ --- violating Axiom 1. The violation persisted through three refinement iterations because it is irresolvable in the classical frame: the spin singlet achieves 100\% constraint reliability ($F_\hbar$) but this decouples from Shannon channel capacity. The framework correctly identified the boundary where classical relational constraints no longer apply. Axiom 1 is not falsified by the quantum world; it correctly diagnoses the domain of its own validity.

\subsubsection{Encodability Does Not Imply Isolation}

The photon encodes cleanly as $\langle D_{\infty}; T_{|}; R_{\supseteq}; K_\text{trap}; G_{\beth}; \Phi_\text{sub}; \ldots \rangle$. The tuple is not null; it describes the photon's \emph{interaction affordances} --- what constraints it can enforce (permanent entanglement, unidirectional propagation) and what it cannot (no dynamic exchange, no partner grammar). The framework does not refuse to describe isolated objects; it \textbf{redescribes them as interaction-ready nodes}, which is already a relational statement. The invalidity surfaces only when \texttt{tensor(photon)} is called without a second argument --- the algebra returns an error because no second operand exists. The tuple encodes relational potential; the algebra computes only when that potential is actualized.

\subsubsection{Ordinal Information Is Sufficient}

The Tier I confirmed predictions (P-1 through P-4 in \texttt{PRIMITIVE\_PREDICTIONS.md}) share a structural feature: \textbf{the causal driver in every case is an ordinal comparison, not an absolute quantity}. The CB{[}7{]} displacement hierarchy was predicted from $F_\hbar > F_\eth > F_\ell$ without knowing the intrinsic binding enthalpies. The Soai Factor 7 firing was predicted from the co-occurrence pattern $D_{\infty} + T_{\bowtie} + P_{+-} + F_{\hbar}$ without knowing the reaction mechanism in detail. The ice VI multi-ordering prediction derived from $K_\text{fast}$ alone, without the proton tunneling rates.

The perturbation universality (P-12) quantifies this: the $K_\text{trap} \to K_\text{MBL}$ cost is $+2.303$ nats $= \ln(10)$ across all three gap-protected topological phases regardless of their distinct $\Omega$ classes. The value is the log of the tier ratio (accessibility drops 10$\times$) --- not an independent physical constant --- but a \emph{discrete tier ratio in a purely ordinal encoding} correctly predicts a \emph{universal energy cost} across three physically distinct topological phases. The tier discretization is the invariant.

\textbf{The framework demonstrates that a classification scheme built entirely from directed relational ordinals is sufficient to generate correct quantitative predictions.} Intrinsic scalar properties --- binding energy, hydrophobicity, gap magnitude --- are not required as inputs. Whether intrinsic properties do not exist or are merely epistemically inert within this syntax is a question the predictions leave open. The distinction between ''not needed'' and ''not there'' is correctly not resolved by the framework.

\emph{Speculative extrapolation (noted):} If Tier III predictions hold --- specifically P-15b ($K_\text{trap} \to K_\text{MBL}$ energy cost confirmed experimentally) and P-19 ($\chi(T \to 0) \sim T^{-\gamma}$ for Factor-8 synthons) --- the +2.303 nat universality becomes a measured physical quantity derivable from ordinal tier structure alone. The implication (consistent with but operationally stronger than Wheeler's ''It from Bit'' and Rovelli's relational quantum mechanics) would be that systems do not have intrinsic kinetic characters that the ordinal approximates, but occupy discrete relational kinetic classes that the ordinal records exactly. The ordinal is not a coarse-graining of the continuous; it may be the correct description.

\begin{center}
\rule{0.5\textwidth}{0.4pt}
\end{center}

\subsection{VIII. Occam Targets --- Three Free Parameters Eliminated}

\emph{Confirmed by computation, 2026-03-17.}

Occam's razor is never fully satisfied, but three apparent free parameters of the framework were shown to be \textbf{derivable from the primitive structure itself} --- reducing the effective degree of freedom count to zero in each case.

\subsubsection{$\lambda$ Has No Free Value: It Is the Primitive Matching Fraction (P-20)}

The tensor formula $\xi_\text{ens} = \xi_1 + \xi_2 - \lambda \cdot I(s_1; s_2)$ uses $I(s_1;s_2) = \text{frac} \cdot \min(\xi_1,\xi_2)$ where frac is the fraction of matching primitive slots. The previously fixed $\lambda = 0.30$ is not a free parameter. Two boundary conditions uniquely determine $\lambda$:

\begin{itemize}
    \item \textbf{Idempotency:} $s \otimes s = s$ requires $\lambda = 1$ when $\text{frac} = 1$. Verified: $\xi(s \otimes s,\, \lambda=1) = \xi(s)$ exactly for all tested synthons.
    \item \textbf{Full synergy:} $\text{frac} = 0 \Rightarrow \lambda = 0$ (orthogonal synthons receive no MI discount).
\end{itemize}

Therefore $\lambda(s_1,s_2) = \text{frac}(s_1,s_2)$. The expected value over all 465 catalog pairs is $\langle \text{frac} \rangle = 0.3023$, within 0.002 of the fixed approximation.

\subsubsection{F-Tier Boundaries Are Integer Boltzmann Discrimination Ratios (P-21)}

The Fidelity values 0.40, 0.75, 0.95 are not empirical thresholds. They are the Boltzmann discrimination fractions $F = \sigma(\Delta\Delta G / kT)$ for three integer selectivity ratios:

\[F_\ell = \tfrac{2}{5} = \sigma\!\left(-\ln\tfrac{3}{2}\right), \qquad F_\eth = \tfrac{3}{4} = \sigma(\ln 3), \qquad F_\hbar = \tfrac{19}{20} = \sigma(\ln 19)\]

The equalities $\text{logit}(3/4) = \ln 3$ and $\text{logit}(19/20) = \ln 19$ are exact. $F_\eth = 3/4$ corresponds to a 3:1 selectivity (one classical cooperative bond, $\Delta\Delta G_0 = \ln 3 \approx 1.10\,kT$). $F_\hbar = 19/20$ corresponds to two cooperative bonds with a $+0.747\,kT$ enhancement over independent additivity --- the cooperative premium that separates the quantum tier from the classical one. $F_\ell = 2/5$ encodes a competitive environment with 1.5 natural competitors dominating the correct recognition event.

\subsubsection{$\Omega$ Is a Derived Label, Not an Independent Primitive (P-22)}

The topological index $\Omega$ is fully determined by five existing primitives $\{T, K, D, \Gamma, G\}$ via a five-rule decision tree with zero mismatches across all 32 catalog synthons:

\begin{table}[htbp]
\centering
\small
\rowcolors{1}{white}{white}
\begin{tabularx}{\textwidth}{>{\centering\arraybackslash}X >{\centering\arraybackslash}X >{\centering\arraybackslash}X}
\toprule
\rowcolor{gray!25}\textbf{Condition} & \textbf{$\Omega$} & \textbf{Physical meaning} \\
\midrule
\rowcolors{1}{white}{gray!10}
\texttt{T=BRAID} $\land$ \texttt{Gamma=QUANTUM\_AND} & $\Omega_{NA}$ & Non-Abelian anyons (FQH Moore-Read) \\
\texttt{T=LINEAR} $\land$ \texttt{K=TRAP} $\land$ \texttt{Gamma=QUANTUM\_AND} & $\Omega_Z$ & Kitaev chain, Majorana zero modes \\
\texttt{T=NETWORK} $\land$ \texttt{D=SUPRAMOLECULAR} $\land$ \texttt{K} $\in$ \{SLOW, TRAP\} & $\Omega_{Z_2}$ & Topological insulator, bulk-boundary \\
\texttt{Gamma=QUANTUM\_AND} or (\texttt{Gamma=SPECIFIC\_AND} $\land$ \texttt{G=GLOBAL}) & $\Omega_0$ & Quantum particles, no topo class \\
otherwise & None & Classical --- no topological protection \\
\bottomrule
\end{tabularx}
\end{table}

\textbf{The effective tuple is reducible from 11 to 10 independent primitives.} $\Omega$ is a convenience label for a five-primitive conjunction; it encodes no information not already present in $\{T, K, D, \Gamma, G\}$.

\subsubsection{Summary of Eliminated Free Parameters}

\begin{table}[htbp]
\centering
\small
\rowcolors{1}{white}{white}
\begin{tabularx}{\textwidth}{>{\centering\arraybackslash}X >{\centering\arraybackslash}X >{\centering\arraybackslash}X >{\centering\arraybackslash}X}
\toprule
\rowcolor{gray!25}\textbf{Parameter} & \textbf{Old status} & \textbf{Derived form} & \textbf{Zero-parameter condition} \\
\midrule
\rowcolors{1}{white}{gray!10}
$\lambda$ & Fixed at 0.30 & $\lambda = \text{frac}(s_1,s_2)$ & Idempotency + full-synergy boundary conditions \\
$F_\eth = 0.75$ & Empirical & $\sigma(\ln 3)$ & Integer selectivity ratio 3:1 \\
$F_\hbar = 0.95$ & Empirical & $\sigma(\ln 19)$ & Integer selectivity ratio 19:1 with cooperative enhancement \\
$F_\ell = 0.40$ & Empirical & $1/(1+\tfrac{3}{2})$ & 1.5 natural competitors in competitive regime \\
$\Omega$ & Independent primitive & $f(T, K, D, \Gamma, G)$ & 5-rule decision tree; 0 mismatches in 32-synthon catalog \\
\bottomrule
\end{tabularx}
\end{table}

\begin{center}
\rule{0.5\textwidth}{0.4pt}
\end{center}

\subsection{IX. Tuple Algebra and Compositional Design}

The ten-tuple $\langle D; T; R; P; F; K; G; \Gamma; \Phi; S \rangle$ is not merely a notation --- it is a mathematical object. Because each primitive takes values in a finite ordered or categorical set, the tuple space admits a rich algebraic structure: a metric space for measuring synthon similarity, a lattice for combining synthons, a directed graph for swap planning, a monoidal product for ensemble prediction, and a set of natural transformations for cross-domain migration. Together these operations form a \textbf{compositional design language} implemented in \texttt{synthomnicon/algebra.py} and exposed as CLI sub-commands.

\begin{center}
\rule{0.5\textwidth}{0.4pt}
\end{center}

\subsubsection{VI.1 The Tuple Quasi-Metric}

\textbf{Distance function.} Define a weighted quasi-metric $d: \mathcal{S} \times \mathcal{S} \to \mathbb{R}_{\geq 0}$ on the synthon space by

\[d(s_1, s_2) = \sum_{p \in \mathcal{P}} w_p \cdot \delta_p(s_1[p], s_2[p])\]

where $\mathcal{P}$ is the set of primitives, $w_p$ are non-negative weights summing to 1, and $\delta_p$ is a primitive-specific local distance:

\begin{itemize}
    \item \textbf{Ordered primitives} ($F$, $K$, $G$): $\delta_p(a, b) = |r(a) - r(b)| / \max\text{-rank}$, where $r$ is the ordinal rank.
    \item \textbf{Categorical primitives} ($D$, $T$, $R$, $P$, $\Gamma$, $\Phi$, $S$): $\delta_p(a, b) = 0$ if $a = b$, else 1.
\end{itemize}

\textbf{Asymmetry.} The metric is a \emph{quasi}-metric: $d(s_1, s_2) \neq d(s_2, s_1)$ in general, because categorical distances are symmetric but the $F$ floor constraint makes directed swaps non-symmetric. This is operationally correct: a swap from high-fidelity to low-fidelity is a different and worse operation than the reverse.

\textbf{Default weights} (matching HotSwap primitive importance):

\begin{table}[htbp]
\centering
\small
\rowcolors{1}{white}{white}
\begin{tabularx}{\textwidth}{>{\centering\arraybackslash}X >{\centering\arraybackslash}X >{\centering\arraybackslash}X}
\toprule
\rowcolor{gray!25}\textbf{Primitive} & \textbf{Weight} & \textbf{Rationale} \\
\midrule
\rowcolors{1}{white}{gray!10}
$D$ & 0.15 & Domain boundary --- swap across $D$ is structurally exceptional \\
$T$ & 0.12 & Topology drives assembly pattern \\
$R$ & 0.12 & Recognition mechanism --- hardest to substitute \\
$P$ & 0.10 & Polarity drives partner compatibility \\
$F$ & 0.18 & Fidelity is the primary thermodynamic constraint \\
$K$ & 0.12 & Kinetic accessibility \\
$G$ & 0.10 & Scale of control \\
$\Gamma$ & 0.08 & Partner selection logic \\
$\Phi$ & 0.07 & Criticality phase \\
$S$ & 0.08 & Stoichiometric ratio \\
\bottomrule
\end{tabularx}
\end{table}

CLI: \texttt{syncon distance S1 S2 {[}--symmetric{]}}

\begin{center}
\rule{0.5\textwidth}{0.4pt}
\end{center}

\subsubsection{VI.2 Lattice Operations: Meet and Join}

\textbf{Partial order.} The tuple space carries a product partial order: $s_1 \leq s_2$ if $s_1[p] \leq s_2[p]$ for all primitives (with the natural order on $F$: $F_{\hbar} \geq F_{\eth} \geq F_{\ell}$; on $K$: $K_{\text{fast}} \geq K_{\text{mod}} \geq K_{\text{slow}} \geq K_{\text{trap}}$; on $G$: $G_{\aleph} \geq G_{\gimel} \geq G_{\beth}$).

\textbf{Meet (greatest lower bound, $\sqcap$).} The meet $s_1 \sqcap s_2$ is the ''most conservative valid synthon below both'':

\[s_1 \sqcap s_2[p] = \begin{cases} \min(s_1[p], s_2[p]) & \text{if } p \in \{F, K, G\} \\ s_1[p] & \text{if } s_1[p] = s_2[p] \text{ (categorical)} \\ \text{CONFLICT} & \text{otherwise (categorical mismatch)} \end{cases}\]

The CONFLICT sentinel propagates through downstream operations and is reported to the user. A meet with a CONFLICT is still a valid result --- it identifies which primitives are incompatible and require resolution.

\textbf{Join (least upper bound, $\sqcup$).} The join $s_1 \sqcup s_2$ is the ''minimal synthon that subsumes both'':

\[s_1 \sqcup s_2[p] = \begin{cases} \max(s_1[p], s_2[p]) & \text{if } p \in \{F, K, G\} \\ s_1[p] & \text{if } s_1[p] = s_2[p] \text{ (categorical)} \\ \text{CONFLICT} & \text{otherwise (categorical mismatch)} \end{cases}\]

\textbf{$\Phi$ resolution rules.} The Criticality Phase primitive ($\Phi$) follows special rules in both operations: $\Phi_c$ (critical) is the maximum in both meet and join --- it propagates from either input to the output. This reflects the physical fact that a system encoding criticality retains that property under both conservative composition (meet) and promotion (join).

\textbf{Interpretation.} Meet is the ''design floor'' --- what properties you are guaranteed to retain if you use either synthon. Join is the ''design ceiling'' --- what properties the assembled system could express. A meet with no CONFLICTs means the two synthons are mutually substitutable within the ordering; any CONFLICTs identify the exact incompatible primitives.

CLI: \texttt{syncon meet S1 S2}, \texttt{syncon join S1 S2}

\begin{center}
\rule{0.5\textwidth}{0.4pt}
\end{center}

\subsubsection{VI.3 Path Algebra: The Valid-Swap Directed Graph}

\textbf{Swap graph.} Define the valid-swap directed graph $\mathcal{G} = (V, E)$ where:

\begin{itemize}
    \item $V$ = all synthons in the catalog with \texttt{grounding\_status = ''full''} or \texttt{''override''}
    \item $(s_1, s_2) \in E$ iff \texttt{HotSwapEngine.validate\_candidate(s1, s2)} returns APPROVED or CONDITIONAL
\end{itemize}

The graph is sparse because the HotSwap validator enforces $D$ and $T$ exact match, and the $F$ floor blocks downgrade edges. Crucially, the graph is \textbf{directed}: an edge $(s_1, s_2)$ exists without implying $(s_2, s_1)$.

\textbf{Path search.} A path $s_0 \to s_1 \to \cdots \to s_n$ gives a multi-step redesign sequence. The cumulative cost is $\xi_{\text{total}} = \sum_i |\Delta\xi_{CP}(s_i \to s_{i+1})|$. Unlike a single swap, the path metric is not transitive: $d(s_0, s_2)$ is not generally $d(s_0, s_1) + d(s_1, s_2)$ because intermediate synthons change the effective $F$ floor. This is the defining asymmetry of the swap graph as a quasi-metric space.

\textbf{BFS algorithm.} Path search uses breadth-first search restricted to the same $\{D, T\}$ cluster (since any cross-$D$ or cross-$T$ edge would trigger a hard HotSwap block), with early termination when cumulative $|\Delta\xi_{CP}|$ exceeds the tolerance. Default tolerance: 2.0 nats. Default max hops: 5.

CLI: \texttt{syncon path SOURCE DESTINATION {[}--max-hops N{]} {[}--xi-tolerance F{]}}

\begin{center}
\rule{0.5\textwidth}{0.4pt}
\end{center}

\subsubsection{VI.4 Tensor Product: Ensemble Prediction}

\textbf{Definition.} The tensor product $s_1 \otimes s_2$ models the effective primitive tuple of an assembled ensemble containing both synthons as active components. The composition rules are:

\begin{table}[htbp]
\centering
\small
\rowcolors{1}{white}{white}
\begin{tabularx}{\textwidth}{>{\centering\arraybackslash}X >{\centering\arraybackslash}X >{\centering\arraybackslash}X}
\toprule
\rowcolor{gray!25}\textbf{Primitive} & \textbf{Rule} & \textbf{Rationale} \\
\midrule
\rowcolors{1}{white}{gray!10}
$D$ & Set union ($D_1 \cup D_2$) & Ensemble operates in all dimensions of its components \\
$T$ & Promote: $T_{\square\square} > T_{\square} > T_{\bowtie} > T_{\ggg}$ & Most complex topology dominates assembly geometry \\
$R$ & Promote: dynamic \textgreater{} static (covalent-dynamic $>$ non-covalent, mechanical, catalytic $>$ covalent) & Dominant interaction governs constraint propagation \\
$P$ & Donor-acceptor ($P_{+-}$) if both are present; else self-complementary if both symmetric & Mixing donor and acceptor produces a directional system \\
$F$ & Min ($F_{\text{eff}} = \min(F_1, F_2)$) & \textbf{Bottleneck principle}: the lower-fidelity component limits the ensemble \\
$K$ & Min ($K_{\text{eff}} = \min(K_1, K_2)$) & Slowest step governs kinetics \\
$G$ & Max ($G_{\text{eff}} = \max(G_1, G_2)$) & Larger-scale control propagates to system level \\
$\Gamma$ & AND-composition: $(op_1 \otimes op_2, \max\text{-tier})$ & Partner requirements are additive \\
$\Phi$ & $\Phi_c$ propagates (if either is critical) & Criticality is contagious --- any critical component makes the ensemble critical \\
$S$ & $S_1 \times S_2$ (stoichiometric product) & Combined assembly has product valency \\
\bottomrule
\end{tabularx}
\end{table}

\textbf{Effective $\xi_{CP}$ correction.} The ensemble $\xi_{CP}$ is not simply $\xi_1 + \xi_2$ --- overlap between the two information channels reduces the total constraint-propagation cost:

\[\xi_{\text{ens}} = \xi_1 + \xi_2 - \lambda \cdot I(s_1; s_2)\]

where $I(s_1; s_2)$ is the mutual information between the two tuple vectors (estimated from shared primitive values), and $\lambda \in [0, 1]$ is a discount parameter (default 0.5). This is the classical ensemble overlap correction in information-theoretic terms.

CLI: \texttt{syncon tensor S1 S2 {[}--lambda F{]}}

\begin{center}
\rule{0.5\textwidth}{0.4pt}
\end{center}

\subsubsection{VI.5 Natural Transformations: Cross-Domain Lifts}

\textbf{Definition.} A \emph{natural transformation} $\eta: F \Rightarrow G$ between two endofunctors on the synthon category is a family of morphisms $\eta_s: F(s) \to G(s)$ that commutes with all HotSwap morphisms. In practice, the four implemented lifts are \textbf{primitive-rewriting maps} that migrate a synthon from one domain to another while preserving as much structure as possible.

\textbf{Lift 1: Temporal ($D_{\wedge} \to D_{\infty}$).} Models what a molecular synthon becomes when embedded in a catalytic cycle:

\[D_{\wedge} \to D_{\infty}, \quad R_{\subseteq} \to R_{\ddagger}, \quad K_{\text{fast}} \to K_{\text{mod}}, \quad \Gamma_{\text{op}} \to \Gamma_{\to}(\text{SEQUENTIAL})\]

\textbf{Lift 2: Spatial ($D_{\wedge} \to D_{\bigtriangleup}$).} Models migration from molecular to crystal-packing domain:

\[D_{\wedge} \to D_{\bigtriangleup}, \quad T \to T_{\square} \text{ if not already spatial}, \quad G_{\beth} \to G_{\gimel}\]

\textbf{Lift 3: Criticality ($\Phi_{\text{sub}} \to \Phi_c$, eligibility-gated).} A promotion to the critical phase, conditional on $F \geq F_{\hbar}$ (the criticality eligibility gate):

\[\Phi_{\text{sub}} \to \Phi_c, \quad G \to G_{\aleph} \text{ (global degeneracy)}\]

This lift is blocked if $F < F_{\hbar}$: a low-fidelity synthon cannot carry a reliable criticality signature. This is the direct operationalization of Axiom 5 in the algebra.

\textbf{Lift 4: Forgetful / Molecular ($D_{\infty} \to D_{\wedge}$).} The forgetful functor --- projects a temporal or supramolecular synthon back to its molecular-level description by dropping the cycle-specific primitives:

\[D_{\infty} \to D_{\wedge}, \quad R_{\ddagger} \to R_{\subseteq}, \quad \Gamma_{\to} \to \Gamma_{\wedge}(\text{SELECTIVE})\]

\textbf{Naturality.} All four lifts commute with HotSwap morphisms. Formally: if $h: s_1 \to s_2$ is a valid HotSwap, then $\eta_{s_2} \circ h = h' \circ \eta_{s_1}$ where $h'$ is the induced swap on the lifted tuples. This holds by construction because the lifts act pointwise on primitives that HotSwap checks are not modified by the lift.

CLI: \texttt{syncon lift SYNTHON (temporal\textbar{}spatial\textbar{}critical\textbar{}molecular)}

\begin{center}
\rule{0.5\textwidth}{0.4pt}
\end{center}

\subsubsection{VI.6 DesignPipeline: The Writer+Maybe Monad}

\textbf{Motivation.} The four operations above (meet, join, tensor, lift, path) compose sequentially in design workflows. Connecting them manually requires threading the current synthon state, accumulating $\xi_{CP}$ costs, and handling failure conditions explicitly at each step. The \texttt{DesignPipeline} class encodes these mechanics as a \textbf{Writer+Maybe monad} in Python:

\begin{itemize}
    \item \textbf{Writer aspect}: automatically accumulates $\Delta\xi_{CP}$ costs across all steps
    \item \textbf{Maybe aspect}: short-circuits on blocking failures (e.g., criticality lift blocked by low F), carrying forward the last valid synthon state rather than crashing
\end{itemize}

\textbf{Monadic chaining.} Each step is a morphism $\text{SynthonM}[A] \to \text{SynthonM}[B]$:

\begin{lstlisting}[language=Python]
DesignPipeline
  .start(synthon)                     # return :: A -> M A (pure lift)
  .meet(other)                        # bind :: M A -> (A -> M B) -> M B
  .join(other)
  .tensor(other, lambda_=0.5)
  .lift("temporal")
  .path("target_name", xi_tolerance=2.0)
  .result()                           # runWriter :: M A -> (A, [Log])
\end{lstlisting}

\textbf{Effect threading.} The pipeline trace records every step's inputs, output tuple, $\Delta\xi_{CP}$, and pass/fail status. A failed step (e.g., a criticality lift that finds $F < F_{\hbar}$) is logged with status BLOCKED and the pipeline continues from the previous valid state --- analogous to \texttt{ExceptT} with a fallback rather than \texttt{MaybeT} with complete abort.

\textbf{Practical equivalence.} The pipeline semantics are equivalent to the following Haskell-style do-notation:

\begin{lstlisting}[language=haskell]
design :: Synthon -> Maybe (Synthon, Cost)
design s = do
  met    <- meet s other
  joined <- join met second
  lifted <- liftCritical joined     -- guard: F >= $F_{\hbar}$
  final  <- path lifted "target"
  return (final, totalDelta\xi)
\end{lstlisting}

Each step propagates the accumulated cost (Writer) and the success/failure flag (Maybe).

CLI: \texttt{syncon pipeline START --step op:arg{[}:key=val{]} {[}--step ...{]}}

The \texttt{--step} flag accepts chained operations in order: \texttt{meet:synthon\_name}, \texttt{join:synthon\_name}, \texttt{tensor:synthon\_name}, \texttt{lift:temporal}, \texttt{path:target\_name:xi\_tolerance=1.5}.

\begin{center}
\rule{0.5\textwidth}{0.4pt}
\end{center}

\subsubsection{VI.7 Compositionality and the Monad Stack}

The full algebra forms a coherent compositional system:

\begin{table}[htbp]
\centering
\small
\rowcolors{1}{white}{white}
\begin{tabularx}{\textwidth}{>{\centering\arraybackslash}X >{\centering\arraybackslash}X >{\centering\arraybackslash}X}
\toprule
\rowcolor{gray!25}\textbf{Operation} & \textbf{Category-theoretic object} & \textbf{Effect} \\
\midrule
\rowcolors{1}{white}{gray!10}
\texttt{meet} & Product / greatest lower bound & Identifies shared design floor; $\Phi_c$ is co-Heyting absorbing \\
\texttt{join} & Coproduct / least upper bound; F-floor ratchet & Identifies minimal common ceiling; raises WriterT floor \\
\texttt{tensor} & Bifunctor; $\xi_\otimes = \xi_1+\xi_2-\lambda I(s_1,s_2)$ & Co-assembly prediction; T promotes; $\Phi_c$ and $\Omega$ join-dominant \\
\texttt{lift} & Functor between domain categories & Cross-domain migration; criticality lift costs 2.303 nats (Landauer analog) \\
\texttt{path} & Geodesic in Kleisli-enriched Lawvere metric space & Directed swap; blocked path = 1st-order topological transition \\
\texttt{pipeline} / \texttt{SynthonM} & $\text{WriterT}[\mathbb{R}_{\geq0}](\text{StateT}[\text{Ctx}](\text{MaybeT}\,\text{Id}))$ & Full monadic do-notation with cost, state, and failure \\
\textbf{decomp} & Inverse operations: cofactor, kernel, project, Birkhoff decomp & Reverse-engineer co-assemblies; Heyting pseudocomplement for design \\
\bottomrule
\end{tabularx}
\end{table}

This structure is not incidental --- it is the consequence of treating the ten-tuple as a typed algebraic object rather than a flat record. The HotSwap protocol (see \texttt{SYNTHONIC\_HOTSWAP.md}) is the arrow category; \texttt{SynthonM} is the monad transformer stack; the four lifts are the functor layer; the decomp operations are the inverse algebra. Together they constitute a \textbf{domain-specific language for compositional synthonic design}, embedded in the Python CLI.

\textbf{Runnable worked examples (18 total, 7 operations, tensor-math framing):} \texttt{TENSOR\_OPS\_DEMO.py} --- run \texttt{python TENSOR\_OPS\_DEMO.py --section \textless{}op\textgreater{}} where \texttt{\textless{}op\textgreater{}} $\in$ \{meet, join, tensor, lift, path, pipeline, decomp\}. Mathematical reference: \texttt{SYNTHONICON\_LANG.md §Algebraic Operations Reference}.

\textbf{Key behavioral properties verified:}

\begin{enumerate}
    \item \textbf{F-floor enforcement}: The F bottleneck in \texttt{tensor} correctly propagates to downstream \texttt{lift} calls --- a tensor that drops F to $F_{\eth}$ blocks the criticality lift without requiring any explicit guard in the pipeline.
    \item \textbf{Path asymmetry}: \texttt{path(nitroso -\textgreater{} Dithia)} fails (F downgrade blocked) while \texttt{path(Dithia -\textgreater{} nitroso)} succeeds --- the quasi-metric asymmetry is correct and intentional.
    \item \textbf{Join-enables-path}: \texttt{join(Dithia, nitroso)} upgrades F to $F_{\hbar}$ (taking the max), enabling a subsequent \texttt{path} to Varma at $\Delta\xi = +0.847$ nat --- the monad correctly threads the upgraded F floor to the downstream operation.
    \item \textbf{Naturality}: Lifting a synthon and then applying a HotSwap produces the same result as applying the HotSwap first and then lifting --- the four natural transformations commute with all valid HotSwap morphisms by construction.
\end{enumerate}

\begin{center}
\rule{0.5\textwidth}{0.4pt}
\end{center}

\subsection{X. Hybrid Systems and Programmable Matter}

The framework's capacity to describe hybrid systems --- where multiple dimensions of assembly are coupled --- is one of its primary strengths. The set-based notation for $D$ and explicit treatment of hybrid recognition modes provide the necessary tools.

\textbf{Canonical hybrid example --- MOF catalysis.} Consider a MOF where secondary building units form a robust three-dimensional framework ($D = D_{\bigtriangleup}$, $G_{\aleph}$, $F_{\hbar}$) with topology $T_{\square}$ and grammar $\Gamma_{\wedge}(\text{SELECTIVE})$, creating specific binding pockets inside which a proline-like catalytic cycle ($D_{\infty}$, $F_{\eth}$, $\Gamma_{\to}(\text{SELECTIVE})$) operates. The complete system encodes as:

\[\langle D_{\bigtriangleup\infty} \;;\; T_{\square+\bowtie} \;;\; R_{\supseteq+\ddagger} \;;\; P_{+} \;;\; F_{\hbar\cdot\eth} \;;\; K_{\text{fast}\cdot\text{mod}} \;;\; G_{\aleph\cdot\gimel} \;;\; \Gamma_{\wedge}(\text{SEL})\cdot\Gamma_{\to}(\text{SEL}) \;;\; \Phi_{\text{sub}} \;;\; 1:1 \rangle\]

The spatial framework acts as a transducer: its coordination recognition mode confines and orients the temporal synthon's components, enhancing effective fidelity and selectivity. $K$ distinguishes the spatial framework's fast chelate re-formation ($K_{\text{fast}}$) from the temporal cycle's rate-determining step ($K_{\text{mod}}$) --- a distinction invisible under a single fidelity label.

\subsubsection{X.1 Programmable Matter as Path Algebra}

The programmable matter domain extends hybrid encoding into a systematic theory of reconfigurability. A programmable material is a physical system that can be commanded to traverse between two or more macroscopic states. In the framework, each macroscopic state of a material is a distinct synthon encoding; programmability is a question about the primitive space connecting them.

\textbf{Core theorem.} ''What can be programmed?'' $=$ ''Does a path exist in the HotSwap graph?'' A material can be reconfigured from state $A$ to state $B$ if and only if a sequence of catalog-intermediated single-primitive steps connects $A$ to $B$ within the HotSwap constraint (same $D$ and $T$ throughout, $|\Delta\xi_{CP}|$ per hop bounded). The path algebra is therefore not a metaphor for reconfigurability --- it is the exact constraint set.

\textbf{Immediate consequences:}

\begin{enumerate}
    \item \emph{Reconfiguration is directed, not symmetric.} Directed distances $d(A \to B) \neq d(B \to A)$ in general. The thermodynamically favored direction is the one with smaller directed distance --- downhill in primitive space corresponds to downhill in free energy.
    \item \emph{The F-floor is an absolute irreversibility condition.} Any transition that requires crossing from $F_\hbar$ to $F_\ell$ is blocked in the HotSwap graph if the source state is also $K_\text{trap}$. Gel arrest ($F_\hbar$, $K_\text{trap}$) cannot be reversed by temperature alone --- the F-floor prevents it algebraically.
    \item \emph{Topology changes are discontinuous.} Programmability pairs whose meet operation conflicts on $T$ (e.g., $T_\text{linear}$ vs $T_\text{network}$) have no HotSwap path. The transition requires a topology discontinuity --- a necessary signature of a first-order phase transition with latent heat. Pairs with $T$-conflict in their meet encode first-order transitions; pairs without $T$-conflict may be second-order.
\end{enumerate}

\subsubsection{X.2 The Eleven Programmable Matter Encodings}

Five material classes spanning six primitive dimensions:

\begin{table}[htbp]
\centering
\small
\rowcolors{1}{white}{white}
\begin{tabularx}{\textwidth}{>{\centering\arraybackslash}X >{\centering\arraybackslash}X >{\centering\arraybackslash}X >{\centering\arraybackslash}X >{\centering\arraybackslash}X >{\centering\arraybackslash}X >{\centering\arraybackslash}X >{\centering\arraybackslash}X}
\toprule
\rowcolor{gray!25}\textbf{Material} & \textbf{$D$} & \textbf{$T$} & \textbf{$F$} & \textbf{$K$} & \textbf{$G$} & \textbf{$\Gamma$} & \textbf{$\Phi$} \\
\midrule
\rowcolors{1}{white}{gray!10}
DNA origami (folded) & $D_{\triangle\wedge}$ & $T_\text{net}$ & $F_\hbar$ & $K_\text{slow}$ & $G_\gimel$ & $\wedge(\text{SPEC})$ & $\Phi_\text{sub}$ \\
DNA strand disp. & $D_\wedge$ & $T_\text{lin}$ & $F_\eth$ & $K_\text{mod}$ & $G_\beth$ & $\to(\text{SEL})$ & $\Phi_\text{sub}$ \\
Colloidal crystal & $D_\triangle$ & $T_{\text{net\_sym}}$ & $F_\hbar$ & $K_\text{slow}$ & $G_\aleph$ & $\wedge(\text{BRD})$ & $\Phi_\text{sub}$ \\
Colloidal fluid & $D_\triangle$ & $T_\text{net}$ & $F_\ell$ & $K_\text{fast}$ & $G_\beth$ & $\vee(\text{BRD})$ & $\Phi_\text{sub}$ \\
Condensate liquid & $D_{\triangle\wedge}$ & $T_\text{net}$ & $F_\ell$ & $K_\text{fast}$ & $G_\gimel$ & $\vee(\text{BRD})$ & $\Phi_c$ \\
Condensate gel & $D_{\triangle\wedge}$ & $T_\text{net}$ & $F_\hbar$ & $K_\text{trap}$ & $G_\aleph$ & $\wedge(\text{BRD})$ & $\Phi_\text{sub}$ \\
Active gel & $D_{\triangle\infty}$ & $T_\text{net}$ & $F_\eth$ & $K_\text{mod}$ & $G_\aleph$ & $\to(\text{SEL})$ & $\Phi_c$ \\
SMP rigid ($T < T_g$) & $D_\triangle$ & $T_\text{net}$ & $F_\hbar$ & $K_\text{slow}$ & $G_\gimel$ & $\to(\text{SEL})$ & $\Phi_\text{sub}$ \\
SMP elastic ($T > T_g$) & $D_\triangle$ & $T_\text{net}$ & $F_\eth$ & $K_\text{mod}$ & $G_\gimel$ & $\wedge(\text{SEL})$ & $\Phi_\text{sub}$ \\
LC nematic & $D_\triangle$ & $T_\text{lin}$ & $F_\eth$ & $K_\text{fast}$ & $G_\gimel$ & $\wedge(\text{BRD})$ & $\Phi_\text{sub}$ \\
LC isotropic & $D_\triangle$ & $T_\text{net}$ & $F_\ell$ & $K_\text{fast}$ & $G_\beth$ & $\vee(\text{BRD})$ & $\Phi_\text{sub}$ \\
\bottomrule
\end{tabularx}
\end{table}

Pairwise programmability distances (symmetric): SMP $1.70$ $<$ LC $3.10$ $<$ condensate $4.00$ $<$ colloidal $5.10$ $<$ DNA $6.10$. This rank predicts the switching energy hierarchy (P-38).

\textbf{State-switching primitives from meet operations.} The meet of each programmability pair identifies which primitives conflict --- these are the material's design levers:

\begin{itemize}
    \item Condensate liquid $\leftrightarrow$ gel: conflict on $\Gamma$ only. The grammar ($\Gamma_\vee$ vs $\Gamma_\wedge$) is the single driver of the liquid-gel transition --- all other primitives are shared or ordered. Therapeutic consequence: dissolving the gel requires changing $\Gamma$ (competing binder) or $K$ (disaggregase). Temperature alone cannot rescue a $K_\text{trap}$ gel (P-39).
    \item SMP rigid $\leftrightarrow$ elastic: conflict on $\Gamma$ only. Shape recovery is a grammar reconfiguration --- the crosslinked network switches from sequential programming ($\Gamma_\to$) to simultaneous elastic recovery ($\Gamma_\wedge$).
    \item LC nematic $\leftrightarrow$ isotropic: conflicts on $T$ and $\Gamma$. The topology change ($T_\text{lin} \to T_\text{net}$) encodes a first-order transition with latent heat (P-40).
\end{itemize}

\subsubsection{X.3 The $F$--$K$ Programmability Quadrant}

The most useful programmable matter sits at $F_\eth + K_\text{mod}$: enough fidelity to maintain states (persistent memory), low enough barrier to switch (responsive). DNA strand displacement and SMP elastic both occupy this quadrant. The four quadrants and their limitations:

\begin{table}[htbp]
\centering
\small
\rowcolors{1}{white}{white}
\begin{tabularx}{\textwidth}{>{\centering\arraybackslash}X >{\centering\arraybackslash}X >{\centering\arraybackslash}X >{\centering\arraybackslash}X >{\centering\arraybackslash}X}
\toprule
\rowcolor{gray!25}\textbf{} & \textbf{$K_\text{fast}$} & \textbf{$K_\text{mod}$} & \textbf{$K_\text{slow}$} & \textbf{$K_\text{trap}$} \\
\midrule
\rowcolors{1}{white}{gray!10}
$F_\hbar$ & unstable (locked but fluctuating) & programmable scaffold (DNA origami) & locked (colloidal crystal, SMP rigid) & irreversible (condensate gel) \\
$F_\eth$ & fast-switching & \textbf{optimal PM} (DNA strand disp., SMP elastic, LC nematic) & slow-switch & slow-switch \\
$F_\ell$ & dynamic reservoir (condensate liquid, colloidal fluid, LC iso.) & dynamic active (active gel) & --- & --- \\
\bottomrule
\end{tabularx}
\end{table}

The Primitive Jacobian $\partial d / \partial \text{primitive}$ quantifies this: $F$ perturbation gives the largest $|\Delta d|$ in 4/5 programmability pairs (DNA: $-1.10$, colloidal: $-1.20$, SMP: $-0.60$, LC: $-0.60$). For condensate gel rescue, $K$ dominates ($-1.50$). Engineering the Jacobian-identified primitive alone achieves $> 40\%$ of the maximum possible distance reduction (P-46).

\subsubsection{X.4 $\Phi_c$ as Global Programmability}

Two of the eleven encodings carry $\Phi_c$: condensate liquid and active gel. These are the two materials with the highest programmability bandwidth --- they can be reconfigured globally from local inputs. The mechanism is G/D degeneracy (§VI): at criticality, a molecular-scale stimulus propagates to the system scale without attenuation.

The tensor product $\text{active\_gel} \otimes \text{DNA strand disp}$ yields a composite carrying $\Phi_c$ and $G_\aleph$ --- an actin-DNA hybrid with global coordination and sequence-specific addressing simultaneously (P-41). $\Phi_c$ is join-dominant in the tensor algebra: it propagates to any composite containing at least one $\Phi_c$ component. This means compositing a $\Phi_c$ material with any programmable matter partner produces a globally coordinatable composite.

Nearest catalog neighbor to condensate liquid: allosteric\_domain at $d = 2.50$. This predicts that condensates implement mesoscale allostery --- a signal entering the droplet propagates to all clients within the droplet via the same G/D degeneracy mechanism as protein-level allostery (P-43).

\subsubsection{X.5 The Programmability Lattice}

The meet of all dynamic/fluid programmable matter states (the dynamic floor) and the join of all rigid/locked states (the rigid ceiling) define the outer limits of the programmable matter design space.

\textbf{Dynamic floor} (meet of condensate liquid, colloidal fluid, SMP elastic, DNA strand disp., LC nematic, LC isotropic):

\[\langle D_\triangle \;;\; T_\text{network} \;;\; R_{\supseteq} \;;\; P_{\text{pm\_sym}} \;;\; F_\ell \;;\; K_\text{mod} \;;\; G_\beth \;;\; \Gamma_\vee(\text{BROAD}) \;;\; \Phi_c \rangle\]

The floor carries $\Phi_c$ --- not by design in any individual encoding, but as an emergent consequence of the lattice meet. This is the deepest result of the programmability analysis: \textbf{the primitive floor of programmability IS criticality} (P-42). Materials engineered to maximize state-space access converge to near-critical encoding by algebraic necessity. This is consistent with measured criticality in biological programmable matter (cortex, cytoplasm) and now derived without biological inputs.

\textbf{Rigid ceiling} (join of DNA origami, colloidal crystal, condensate gel, SMP rigid): conflicts on $D$, $P$, and $\Gamma$. The conflicts encode the essential diversity of locking mechanisms --- no single primitively-described ''lock'' is shared by all rigid states. Consequence: there is no universal locking agent. Each material class requires a locking strategy matched to its conflict primitives.

\textbf{Programmability gap}: $d(\text{dynamic floor} \to \text{rigid ceiling}) = 5.40$ nats. The primitives spanning this gap ($F$, $K$, $G$) are the design axes of programmable matter engineering.

\subsubsection{X.6 Cross-Domain Analogies}

The catalog metric identifies structural isomorphisms between programmable matter and the existing synthon catalog:

\begin{table}[htbp]
\centering
\small
\rowcolors{1}{white}{white}
\begin{tabularx}{\textwidth}{>{\centering\arraybackslash}X >{\centering\arraybackslash}X >{\centering\arraybackslash}X >{\centering\arraybackslash}X >{\centering\arraybackslash}X}
\toprule
\rowcolor{gray!25}\textbf{PM system} & \textbf{Nearest catalog neighbor} & \textbf{$d$} & \textbf{Shared primitive cluster} & \textbf{Prediction} \\
\midrule
\rowcolors{1}{white}{gray!10}
Condensate liquid & allosteric\_domain & 2.50 & $F_\ell$, $K_\text{fast}$, $G_\gimel$, $\Phi_c$ & Condensates implement mesoscale allostery (P-43) \\
Active gel & allosteric\_domain & 3.40 & $F_\eth$, $K_\text{mod}$, $\Phi_c$ & Active gels implement distributed allostery with ATP as signal \\
Colloidal crystal & topological\_insulator & 2.80 & $F_\hbar$, $G_\aleph$, $T_\text{net}$, $R_{\supseteq}$ & Topologically protected boundary states (P-44) \\
DNA origami & topological\_insulator & 3.70 & $F_\hbar$, $G_\gimel$, $R_{\supseteq}$ & Topological DNA origami: polynomial error scaling (P-45) \\
Active gel & tide\_pool\_ecological & 4.10 & $D_{\triangle\infty}$, $G_\aleph$, $\Phi_c$ & Active gels as ecological-scale analogs: energy dissipation, collective behavior \\
\bottomrule
\end{tabularx}
\end{table}

The colloidal crystal--topological insulator analogy ($d = 2.80$) is empirically confirmed (topological colloidal matter, Rechtsman 2016). The framework extends it: the boundary is the primitive distance threshold $d < 3.0$. Any colloidal crystal with $F_\hbar$, $G_\aleph$, $T_\text{network}$, and tuned interaction symmetry should exhibit topological boundary states.

\begin{center}
\rule{0.5\textwidth}{0.4pt}
\end{center}

\subsection{XI. Experimental and Computational Validation Across Domains}

\subsubsection{Molecular domain}

\textbf{Transformation \#1 --- Polarity shift and fidelity ranking in hydrogen-bonded dimers.} Single-point DFT at B3LYP-D3(BJ)/6-311+G(d,p) with counterpoise (BSSE) correction on the formic acid homodimer (AA:AA), formic acid/formamide heterodimer (AA:amide), and formamide homodimer (amide:amide), all in the R$_2^2$(8) cyclic motif. BSSE-corrected electronic binding energies ($\Delta E$): --64.2 kJ/mol (AA:AA), --51.8 kJ/mol (AA:amide), --39.6 kJ/mol (amide:amide). The corresponding $\Delta G_{298}$ (gas phase, entropy-included) are substantially smaller: approximately \textbf{--12 kJ/mol} (AA:AA), \textasciitilde{}--10 kJ/mol (AA:amide), \textasciitilde{}--8 kJ/mol (amide:amide), reflecting the large entropic cost of bimolecular association at 298 K; solution values require ITC measurement. The fidelity ratio (\textasciitilde{}1.9 from the $\Delta E$ ordering) is preserved in the $\Delta G$ series and is consistent with CSD propensity data for the R$_2^2$(8) motif. The sequence demonstrates the $P_{\pm}^{\psi} \to P_{+-}$ polarity transition and the correlated $F_{\hbar} \to F_{\eth} \to F_{\ell}$ fidelity gradient. All three are $T_{\bowtie}$/$P_{\pm}$ and achieve $F \geq F_{\ell}$, consistent with Axiom 1. All three are $K_{\text{fast}}$, consistent with rapid H-bond exchange.

\textbf{Transformation \#4 --- Static covalent to dynamic covalent: emergence of $D_{\infty}$ character.} Gas-phase proxy at B3LYP-D3/6-311+G(d,p): formaldehyde + ammonia $\to$ methanimine + H$_2$O gives forward endergonicity of +38.7 kJ/mol and condensation barrier \textasciitilde{}162 kJ/mol (gas phase), dropping to \textasciitilde{}90--120 kJ/mol in aqueous solution with general acid/base assistance. The barrier reduction confers $D_{\infty}$ character: the imine linkage undergoes error-correction and self-healing through hydrolysis/re-condensation, shifting from $R_{\subseteq}$ to $R_{\subseteq + \ddagger}$ and from $F_{\hbar}$ (static) to $F_{\eth}$ (dynamic). The gas-phase system is $K_{\text{slow}}$; the aqueous system is $K_{\text{mod}}$ --- distinct $K$ assignments at identical $F_{\eth}$, representing operationally different systems under a shared thermodynamic classification.

\subsubsection{Supramolecular domain}

\textbf{Transformation \#3 --- Chelate node: granularity amplification via coordination geometry.} Single-point DFT at B3LYP-D3(BJ)/6-311+G(d,p): cumulative BSSE-corrected binding energies of --263.1 kJ/mol for {[}Zn(py)$_2$Cl$_2${]} (two monodentate pyridines) versus --312.6 kJ/mol for {[}Zn(bpy)Cl$_2${]} (one bidentate bipyridine). Chelate gain: \textasciitilde{}49 kJ/mol gas phase; \textasciitilde{}60--90 kJ/mol solvated (SMD/water, literature-calibrated), consistent with experimental $K_{\text{chelate}}/K_{\text{mono}}^2 \approx 10^2$--$10^4$. The bidentate geometry locks the metal centre into a fixed bite-angle, elevating control from $G_{\beth}$ to $G_{\aleph}$: a single binding event enforces the spatial arrangement of the entire coordination sphere. $K_{\text{fast}}$ for re-formation makes the effective operational reliability ($F_{\hbar}$, $K_{\text{fast}}$) higher than $F$ alone predicts.

\textbf{Transformation \#5 --- Cooperative fidelity amplification: the triple H-bond array.} SAPT2+/aug-cc-pVDZ decompositions on single (acetamide dimer), double (urea-like), and triple (2-aminopyridine$\cdot$2-pyridone, DAD$\cdot$ADA) H-bond arrays. Electronic binding energies ($\Delta E$): \textasciitilde{}--30 kJ/mol (single), \textasciitilde{}--60 kJ/mol (double), \textasciitilde{}--95 to --110 kJ/mol (triple). The corresponding $\Delta G_{298}$ (gas, entropy-included) for the triple array is approximately \textbf{--55 kJ/mol}; the large decrease from $\Delta E$ reflects the entropic cost of organised three-point pre-alignment at 298 K. H-bond geometry from literature consensus: N--H$\cdot$$\cdot$$\cdot$N/O distances 1.80/1.90/1.80 Å (outer contacts tighter than central, consistent with cooperative shortening), angles 170$^{\circ}$/163$^{\circ}$/170$^{\circ}$. Induction component: \textasciitilde{}10--15\% of total (single), \textasciitilde{}20--30\% (double), \textasciitilde{}30--40\% (triple), superlinear at 2.5--3.5$\times$ the single-bond value while electrostatics remain approximately additive; cooperativity factor 1.25 (literature range 1.2--1.4). This superadditivity is the direct computational signature of Axiom 3: the triple array crosses from $G_{\beth}$ to $G_{\gimel}$ via the induction superlinearity threshold. The $\xi_{CP}$ gap from single to triple (\textasciitilde{}1.5--2.0 nats) is the first quantitative measure of emergent fidelity amplification in the framework.

\textbf{Transformation \#7 --- Fidelity ranking via $\sigma$-hole depth: halogen vs. chalcogen bond.} Single-point DFT at B3LYP-D3(BJ)/6-311+G(d,p) + BSSE: 4-iodopyridine dimer (I$\cdot$$\cdot$$\cdot$N halogen bond) --28.4 kJ/mol; 4-(methylthio)pyridine dimer (S$\cdot$$\cdot$$\cdot$N chalcogen bond) --14.9 kJ/mol (fidelity ratio 1.91). Multiwfn ESP analysis: $V_{\text{max}} \approx +165$ kJ/mol (iodine) vs. $+105$ kJ/mol (sulfur). The dual confirmation (energetic and electrostatic) validates the $F_{\eth} \to F_{\ell}$ and $\Gamma_{\wedge}(\text{SELECTIVE}) \to \Gamma_{\vee}(\text{BROAD})$ transitions: the weaker, less-directional $\sigma$-hole of sulfur broadens the grammar of acceptable interaction partners. This demonstrates mechanistic coupling between $F$ and $\Gamma$ through a shared electronic property.

\subsubsection{Temporal domain}

\textbf{Transformation \#6 --- Proline-catalyzed aldol cycle: first quantitative temporal fidelity.} The proline-catalyzed enamine aldol cycle mapped at M06-2X/6-31+G(d,p), CPCM(DMSO). $\Delta G^{\ddagger}$ for the C--C bond-forming step (Houk group and related studies): \textbf{97 kJ/mol} operative value (range 92--100 kJ/mol across substrates). Transition state geometry from literature consensus: N--H$\cdot$$\cdot$$\cdot$O distance 1.825 Å (intramolecular H-bond in the Zimmermann--Traxler chair, rate-controlling), C--C forming distance 2.10 Å, imaginary frequency \textasciitilde{}400 cm$^{-1}$ along the C--C coordinate. Eyring TST at 298 K: $k_{\text{cat}} \approx 10^{-3}$ to $10^{-2}$ s$^{-1}$; against $k_{\text{side}} \approx 10^{-7}$ s$^{-1}$:

\[F_{\text{cycle}} = \frac{k_{\text{cat}}}{k_{\text{cat}} + k_{\text{side}}} \approx 0.999 \text{ to } 0.9999\]

The facial selectivity ($\Delta\Delta G^{\ddagger}_{\text{si}-\text{re}}$) from DFT studies of 5--8 kJ/mol predicts \textbf{ee = 70--85\%}, in agreement with the experimental value of \textbf{74\% ee} (acetone + 4-nitrobenzaldehyde, DMSO). This constitutes the framework's first quantitatively grounded cross-domain prediction: a $\xi_{CP}$-based efficiency estimate for a temporal synthon tied to a measured stereochemical outcome. This $F_{\eth}$ value for a $D_{\infty}$ synthon is directly comparable to fidelity ratios from binding energies in Transformations \#1 and \#7 --- the first formal cross-domain $F$ comparison. The $K_{\text{mod}}$ assignment (rate-determining step accessible under mild organocatalytic conditions) distinguishes the proline cycle from temporal synthons sharing its $F_{\eth}$ classification but differing in operational accessibility.

\subsubsection{Network and bulk material domain}

\textbf{Ice polymorph family --- catalog self-audit as primitive discovery engine (v0.3.6--0.3.7).} The eleven phases of ice (I$_\text{h}$, I$_\text{c}$, III--XI) were generated as synthons and compared. The identical-tuple problem emerged immediately: before sub-labels, \{I$_\text{h}$, XI\} and \{III, IX\} and \{I$_\text{c}$, VIII, IX\} produced identical tuples under the generic $T_{\in}$ assignment. These collisions are not encoding failures --- they are the framework's self-correction signal: \textbf{identical tuples reveal missing primitives}. Analysis of the 22 ambiguous $T_{\bowtie}$ entries in a separate catalog repair pass simultaneously revealed $T_{\cup}$ (bowl topology). Two primitives were discovered in one session by using the catalog as an audit instrument against itself.

The four $T_{\in}$ sub-labels resolve all remaining ice-phase collisions:

\begin{table}[htbp]
\centering
\small
\rowcolors{1}{white}{white}
\begin{tabularx}{\textwidth}{>{\centering\arraybackslash}X >{\centering\arraybackslash}X >{\centering\arraybackslash}X}
\toprule
\rowcolor{gray!25}\textbf{Phase(s)} & \textbf{$T$} & \textbf{Physical basis} \\
\midrule
\rowcolors{1}{white}{gray!10}
I$_\text{h}$, I$_\text{c}$, XI & $T_{\in}(\text{hex})$ & Pure 6-membered hexagonal rings, tetrahedral 4-coordination \\
III, IV, V, IX & $T_{\in}(\text{mixed})$ & Mixed ring sizes (4+5+6+7+8-membered), distorted coordination \\
VI, VII, VIII & $T_{\in}(\times 2)$ & Two independent interpenetrating bcc H-bond sub-networks \\
X & $T_{\in}(\text{sym})$ & Centrosymmetric O--H--O bond at $\geq$70 GPa; proton exactly centred \\
\bottomrule
\end{tabularx}
\end{table}

\textbf{$K_{\text{fast}}$ as causal primitive in ice VI.} The $K_{\text{fast}}$ assignment for ice VI encodes a measurable physical property --- the proton reorientation dynamics --- and predicts downstream consequences that the framework does not explicitly contain. Because $K_{\text{fast}}$ means $\Delta G^\ddagger < 60$ kJ/mol (constraint rearrangement accessible at experimental timescales), the system can:

\begin{enumerate}
    \item \textbf{Explore multiple ordering landscapes during slow cooling.} At $\sim$1.0 GPa $\to$ ice XV (Salzmann 2009); at $>$1.5 GPa $\to$ ice XIX (Yamane 2021, Gasser 2021). Two distinct antiferroelectrically ordered descendants from a single disordered parent.
    \item \textbf{Form deep-glassy states under rapid cooling} (Rosu-Finsen \& Salzmann 2020) --- partially ordered, low-enthalpy configurations accessible only because proton dynamics are fast enough to be trapped before reaching the ordered ground state.
    \item \textbf{Allow quantum tunneling regime} at very low temperatures (Bove et al. 2009), where $K_{\text{fast}}$ extends beyond classical activation.
\end{enumerate}

Dielectric relaxation measurements (Yamane et al. 2021) directly confirm the assignment: the relaxation time of ice VI is orders of magnitude shorter than those of ice XV or XIX. The framework encoded this without access to any of that data.

\textbf{$K_{\text{fast}} \to K_{\text{slow}}$ as the ordering transition.} Ice XV and XIX encode as $\langle D_\triangle; T_{\in}(\times 2); R_\supseteq; P_{+-}; F_\eth; K_{\text{slow}}; \ldots \rangle$ --- differing from ice VI in exactly one primitive: $K$. The $K_{\text{fast}} \to K_{\text{slow}}$ flip \emph{is} the ordering transition. The entire phase boundary is encoded in a single primitive change.

\textbf{$G_{\gimel}$ vs. $G_{\beth}$ encodes pressure-dependent correlation length.} Ice XV ($\sim$1.0 GPa) retains $G_{\gimel}$ (MESOSCALE) --- cooperative ordering propagates across the network. Ice XIX (\textgreater{}1.5 GPa) encodes as $G_{\beth}$ (LOCAL) --- at higher pressure, shorter O--O distances enforce ordering more rigidly at the local level, consistent with a shorter ordering correlation length. The LLM assigned this distinction without prompting. It is a falsifiable prediction against future neutron diffraction data on correlation lengths in ice XIX vs. XV.

\textbf{Constraint strength ordering.} The computed $\xi_{\text{constraint}}$ values across the family span 0.405 (ice I$_\text{c}$, $F_\ell$, $K_{\text{slow}}$) to 0.733 (ice VII, $G_{\aleph}$, $K_{\text{mod}}$). Ice I$_\text{c}$ is the only LOW fidelity phase --- encoding metastability and ready transformation to I$_\text{h}$. Ice VII's global granularity ($G_{\aleph}$) and its position as the highest-constraint phase encodes the topological entanglement of two bcc sub-networks.

\subsubsection{Quantum domain --- Axiom 1 as a classical boundary detector}

\textbf{Entangled particle series (2026-03-16).} Five quantum systems were encoded using the framework: entangled photon, proton, electron, spin, and qubit. Each word, applied to a different physical domain, yielded a distinct, physically defensible tuple. No chemical template was available; the LLM reasoned from the primitive definitions alone.

\begin{table}[htbp]
\centering
\small
\rowcolors{1}{white}{white}
\begin{tabularx}{\textwidth}{>{\centering\arraybackslash}X >{\centering\arraybackslash}X >{\centering\arraybackslash}X >{\centering\arraybackslash}X >{\centering\arraybackslash}X >{\centering\arraybackslash}X >{\centering\arraybackslash}X}
\toprule
\rowcolor{gray!25}\textbf{System} & \textbf{T} & \textbf{R} & \textbf{K} & \textbf{G} & \textbf{Axiom 1} & \textbf{Physical basis} \\
\midrule
\rowcolors{1}{white}{gray!10}
photon & $T_{|}$ & $R_{\supseteq}$ & $K_{\text{trap}}$ & $G_{\beth}$ & pass & quantum correlation; non-exchangeable \\
proton & $T_{\square}$ & $R_{\ddagger}$ & $K_{\text{fast}}$ & $G_{\beth}$ & pass & dynamic proton transfer; $\Delta G^{\ddagger} \ll$ 60 kJ/mol \\
electron & $T_{\bowtie}$ & $R_{\ddagger}$ & $K_{\text{fast}}$ & $G_{\beth}$ & warn† & spin-singlet pair; dynamic exchange \\
\textbf{spin} & $T_{\bowtie}$ & $R_{\supseteq}$ & $K_{\text{trap}}$ & $G_{\aleph}$ & \textbf{FALSIFIED} & non-local correlation; Bell non-locality \\
qubit & $T_{|}$ & $R_{\ddagger}$ & $K_{\text{fast}}$ & $G_{\aleph}$ & pass & computational unit; gate operations \\
\bottomrule
\end{tabularx}
\end{table}

†Axiom 7 warning: $T_{\bowtie}$ without a named closing bond (quantum correlation has no classical closing bond).

\textbf{Key findings:}

\begin{enumerate}
    \item \textbf{$G_{\aleph}$ appears for the first time} --- on spin and qubit only. Every classical system in the catalog is $G_{\beth}$ or $G_{\gimel}$. The LLM correctly encoded quantum non-locality as global granularity: spin-spin correlation propagates without spatial attenuation (Bell inequality). This is the first appearance of $G_{\aleph}$ for a non-network system.
    \item \textbf{Axiom 1 is falsified by entangled spin.} The LLM assigned $T_{\bowtie} + P_{\pm}^{\text{sym}} + F_{\ell}$, which violates Axiom 1 ($T_{\bowtie} + P_{\pm} \to F \geq F_{\eth}$). The violation persisted through 3 refinement iterations because it is \emph{irresolvable within the classical frame}: the LLM correctly assigned $F_{\ell}$ for the wrong reason (no-communication theorem), but the framework correctly identified the inconsistency.
    \item \textbf{The resolution exposes a primitive definition boundary.} $F$ measures \emph{reliability of the constraint}, not classical information capacity. The spin singlet constraint fires with 100\% reliability (measuring one spin $\to$ guaranteed anti-parallel result) $\to$ $F_{\hbar}$. The no-communication theorem is about Shannon capacity; $F$ is about constraint fidelity. These are orthogonal. Corrected tuple: $\langle D_{\wedge}; T_{\bowtie}; R_{\supseteq}; P_{\pm}^{\text{sym}}; F_{\hbar}; K_{\text{trap}}; G_{\aleph}; \Gamma_{\wedge}(\text{SELECTIVE}); \Phi_{\text{sub}} \rangle$.
    \item \textbf{The Axiom 1 violation pattern $T_{\bowtie} + P_{\pm} + F_{\ell}$ is a quantum boundary signature.} In the classical domain this combination is forbidden (Axiom 1) and physically unreachable --- cyclic self-complementary motifs always amplify fidelity above $F_{\ell}$. In the quantum domain it is reachable: $F$ and Shannon capacity decouple. \textbf{Axiom 1 is not falsified; it correctly identifies the classical boundary.} Any system that reaches $T_{\bowtie} + P_{\pm} + F_{\ell}$ without resolving after iterative refinement is a quantum system or operates by a mechanism outside the classical constraint-propagation framework.
    \item \textbf{Photon vs. proton vs. electron vs. spin:} $K$ encodes the kinetic regime of the ''entanglement'' correctly without domain-specific training. Photon: $K_{\text{trap}}$ (permanent, non-exchangeable). Proton: $K_{\text{fast}}$ (proton transfer is among the fastest classical processes). Electron: $K_{\text{fast}}$ (electron exchange/tunneling fast). Spin: $K_{\text{trap}}$ (spin-singlet correlation is permanent, like photon). Qubit: $K_{\text{fast}}$ (gate operations are dynamic/resettable).
\end{enumerate}

\subsubsection{Quantum domain --- Topological matter (v0.4.0 catalog entries)}

\textbf{Three canonical topological synthons (2026-03-16).} First catalog entries using the new \texttt{T\_braid} and $\Omega$ primitives. Each encodes a distinct topological universality class.

\begin{table}[htbp]
\centering
\small
\rowcolors{1}{white}{white}
\begin{tabularx}{\textwidth}{>{\centering\arraybackslash}X >{\centering\arraybackslash}X >{\centering\arraybackslash}X >{\centering\arraybackslash}X >{\centering\arraybackslash}X >{\centering\arraybackslash}X}
\toprule
\rowcolor{gray!25}\textbf{System} & \textbf{T} & \textbf{$\Omega$} & \textbf{K} & \textbf{Factor 8?} & \textbf{Topological invariant} \\
\midrule
\rowcolors{1}{white}{gray!10}
\texttt{kitaev\_chain\_majorana} & $T_{|}$ & $\Omega_Z$ & $K_{\text{trap}}$ & \textbf{YES} & $\mathbb{Z}$ class D, 1D (winding number W=1) \\
\texttt{fqh\_moore\_read} & $T_{\text{braid}}$ & $\Omega_{\text{NA}}$ & $K_{\text{trap}}$ & \textbf{YES} & Non-Abelian (Ising anyons, GSD=3 on torus) \\
\texttt{topological\_insulator\_bi2se3} & $T_{\in}$ & $\Omega_{Z_2}$ & $K_{\text{slow}}$ & NO & $\mathbb{Z}_2$ class AII (Kramers, $\nu_0$=1) \\
\bottomrule
\end{tabularx}
\end{table}

\textbf{Factor 8 fires on the gapped topological phases but not on the TI.} The Kitaev chain and FQH state both satisfy $G_{\aleph}$ + $F_{\hbar}$ + $K_{\text{trap}}$ + $\neg D_{\infty}$ $\to$ quantum criticality score +0.200. The topological insulator has $K_{\text{slow}}$ (not $K_{\text{trap}}$) --- its surface states flow, they are not gap-frozen at a quantum critical point. This is physically correct: TIs are stable gapped insulators, not quantum critical points; the Kitaev chain and FQH state sit near topological phase transitions that are quantum critical.

\textbf{$\Omega$ lattice semantics confirmed numerically:}

\emph{Meet (conservative protection guarantee):}

\begin{lstlisting}
spin_singlet meet kitaev_chain -> \Omega = TRIVIAL   (\Omega_0 meet \Omega_Z  -> \Omega_0:  singlet has no protection)
fqh_moore_read meet TI         -> \Omega = Z2_CLASS  (\Omega_NA meet Omega_Z2 -> Omega_Z2: conservative)
kitaev_chain meet qubit        -> \Omega = TRIVIAL   (\Omega_Z meet \Omega_0  -> \Omega_0:  qubit pulls down)
\end{lstlisting}

\emph{Join (capability ceiling):}

\begin{lstlisting}
spin_singlet join kitaev_chain -> \Omega = Z_CLASS     (\Omega_0 join \Omega_Z  -> \Omega_Z:  chain lifts the singlet)
fqh_moore_read join TI         -> \Omega = NON_ABELIAN (\Omega_NA join Omega_Z2 -> \Omega_NA: non-Abelian dominates)
kitaev_chain join qubit        -> \Omega = Z_CLASS     (\Omega_Z join \Omega_0  -> \Omega_Z:  chain lifts qubit)
\end{lstlisting}

\emph{Tensor (ensemble inherits stronger protection):}

\begin{lstlisting}
tensor(spin, kitaev_chain) -> \Omega = Z_CLASS     (\Omega_Z beats \Omega_0; spin gains topological character)
tensor(fqh, TI)            -> \Omega = NON_ABELIAN (\Omega_NA beats Omega_Z2; FQH dominates)
tensor(fqh, kitaev_chain)  -> \Omega = NON_ABELIAN (\Omega_NA beats \Omega_Z)
tensor(kitaev_chain, qubit)-> \Omega = Z_CLASS     (\Omega_Z beats \Omega_0; chain protects the qubit)
tensor(spin, fqh)          -> \Omega = NON_ABELIAN (\Omega_NA beats \Omega_0; singlet inherits non-Abelian)
\end{lstlisting}

\textbf{\texttt{tensor(kitaev\_chain, qubit)} is the topological qubit prediction.} Coupling an unprotected qubit ($F_{\ell}$, $K_{\text{slow}}$, $\Omega_0$) to a Kitaev chain ($F_{\hbar}$, $K_{\text{trap}}$, $\Omega_Z$) yields an ensemble with $F_{\ell}$ (the fidelity bottleneck) but $\Omega_Z$. This encodes exactly the experimental challenge: the Kitaev chain provides topological protection ($\Omega_Z$) but the practical qubit interface degrades overall fidelity ($F_{\ell}$). The framework identifies $F$ as the load-bearing primitive for the qubit quality --- not $\Omega$.

\textbf{Distance matrix (8-synthon quantum cluster, selected pairs):}

\begin{table}[htbp]
\centering
\small
\rowcolors{1}{white}{white}
\begin{tabularx}{\textwidth}{>{\centering\arraybackslash}X >{\centering\arraybackslash}X >{\centering\arraybackslash}X}
\toprule
\rowcolor{gray!25}\textbf{Pair} & \textbf{d(sym)} & \textbf{Dominant mismatches} \\
\midrule
\rowcolors{1}{white}{gray!10}
proton $\leftrightarrow$ electron & 1.80 & P only (+ vs. $-$) \\
spin\_singlet $\leftrightarrow$ qubit\_logical & 3.20 & F ($\hbar$ vs. $\ell$) + K (trap vs. slow) \\
spin\_singlet $\leftrightarrow$ kitaev\_chain & 3.70 & T ($\bowtie$ vs. \$ \\
kitaev\_chain $\leftrightarrow$ fqh & 5.70 & D ($\wedge$ vs. $\triangle$) + T (\$ \\
fqh $\leftrightarrow$ TI & 5.50 & T (braid vs. $\in$) + K (trap vs. slow) + $\Omega$ (NA vs. $Z_2$) \\
proton $\leftrightarrow$ fqh & 10.20 & Max separation in 8-synthon cluster \\
\bottomrule
\end{tabularx}
\end{table}

\textbf{Perturbation sweeps confirm K as the load-bearing primitive for all gapped topological systems:}

\begin{table}[htbp]
\centering
\small
\rowcolors{1}{white}{white}
\begin{tabularx}{\textwidth}{>{\centering\arraybackslash}X >{\centering\arraybackslash}X >{\centering\arraybackslash}X >{\centering\arraybackslash}X}
\toprule
\rowcolor{gray!25}\textbf{Synthon} & \textbf{Most sensitive prim.} & \textbf{$\Delta\xi$} & \textbf{Interpretation} \\
\midrule
\rowcolors{1}{white}{gray!10}
spin\_singlet & K: TRAP$\to$MBL & +2.303 nats {[}HIGH{]} & MBL freezes the singlet further; maximum cost \\
kitaev\_chain & K: TRAP$\to$MBL & +2.303 nats {[}HIGH{]} & Gap$\to$disorder transition; MBL destroys Majorana coherence \\
fqh\_moore\_read & K: TRAP$\to$MBL & +2.303 nats {[}HIGH{]} & Gap$\to$MBL in 2D; destroys anyonic order \\
TI & K: SLOW$\to$MOD & $-$0.847 nats {[}MED{]} & \emph{Negative} --- improving surface mobility reduces cost \\
\bottomrule
\end{tabularx}
\end{table}

The $K_{\text{trap}} \to K_{\text{MBL}}$ shift at +2.303 nats is \textbf{universal across all gap-protected topological systems}. This is not a coincidence: it is the thermodynamic cost of crossing from a coherent many-body gap ($K_{\text{trap}}$) into a disorder-frozen many-body localized state ($K_{\text{MBL}}$), regardless of the specific topological invariant. The framework discovered this universality by encoding the ordinal structure of $K$ correctly. \textbf{Falsifiable prediction:} experimental preparation of an MBL phase from a gap-protected topological phase should require a free-energy cost consistent with $\Delta\xi \sim \ln(2) \times 2.303 \approx 1.6$ nats per degree of freedom, measurable via thermodynamic integration or quench spectroscopy.

The TI shows the opposite sign because $K_{\text{slow}} \to K_{\text{mod}}$ is \emph{decreasing} the kinetic barrier --- this is the direction of better surface transport, which \emph{reduces} thermodynamic cost. The TI is the only synthon in the cluster where improving kinetics reduces cost rather than increasing it.

\subsubsection{Quantum domain --- Tuple-space phase diagram (v0.4.0)}

\textbf{The framework detects phase boundaries from syntax alone.} The \texttt{phase\_diagram} module computes all pairwise tuple distances across the eight quantum/topological synthons, applies Ward hierarchical clustering, and projects into 2-D via metric MDS. No physics is input; the phase structure emerges from the primitive encoding.

\begin{figure}[htbp]
\centering
\includegraphics[width=0.8\textwidth]{synthomnicon\_phase\_diagram\_v0.4.0.png}
\caption{SynthOmnicon Tuple-Space Phase Diagram v0.4.0}
\end{figure}


\emph{Generated by \texttt{syncon phase-diagram --save synthomnicon\_phase\_diagram\_v0.4.0.png}. Annotations: {[}star{]} = Factor-8 quantum criticality trigger; {[}circle{]} = $K_{\text{trap}}$ ($K_{\text{MBL}}$ candidate); $\Omega$ colour encodes topological class (grey=TRIVIAL, blue=$\Omega_Z$, green=$\Omega_{Z_2}$, red=$\Omega_{\text{NA}}$).}

\textbf{Dendrogram (left panel).} The Ward linkage produces a clean two-branch structure separated at $d \approx 9.52$ (the primary phase boundary, dashed orange line):

\begin{itemize}
    \item \textbf{Branch 1 --- extended topological matter} (\texttt{fqh\_moore\_read} + \texttt{topological\_insulator\_bi2se3}): These are the $D_\triangle$ systems --- 2D/3D bulk materials with collective ground states. They separate at $d \approx 5.5$ internally. \texttt{fqh\_moore\_read} ($\Omega_{\text{NA}}$) and \texttt{topological\_insulator\_bi2se3} ($\Omega_{Z_2}$) differ primarily in T ($T_{\text{braid}}$ vs $T_{\in}$), K ($K_{\text{trap}}$ vs $K_{\text{slow}}$), and $\Omega$ --- encoding the distinction between a \emph{gapped topological phase at a quantum critical point} and a \emph{stable gapped insulator with protected surface states}.
    \item \textbf{Branch 2 --- quantum particles and engineered qubits}: All five particle synthons plus \texttt{kitaev\_chain\_majorana} cluster here. Within this branch three sub-clusters are visible: \texttt{\{proton, electron\}} (d=1.80, charge-sign only), \texttt{photon} (isolated by $D_{\infty}$), and the quantum computing cluster \texttt{\{spin\_singlet, qubit\_logical, kitaev\_chain\}} (d$\approx$3.2--3.9 internally).
\end{itemize}

\textbf{MDS phase map (right panel).} The 2-D metric projection preserves the distance structure and makes the phase boundaries visual:

\begin{itemize}
    \item \texttt{fqh\_moore\_read} ({[}star{]} red, $K_{\text{trap}}$ ring) is maximally isolated at the far right --- the hardest phase boundary in the catalog, d=10.2 from proton/electron. Its isolation encodes the fact that a fractional quantum Hall state is not continuously deformable into a free particle without closing a gap and changing T, K, $\Omega$, D simultaneously.
    \item \texttt{spin\_singlet} and \texttt{kitaev\_chain} ({[}star{]} Factor-8, $K_{\text{trap}}$ rings) cluster in the center-right, adjacent to \texttt{qubit\_logical}. These are the \emph{quantum computing cluster} --- coherent few-body systems near a quantum critical point that could host logical qubits.
    \item \texttt{topological\_insulator\_bi2se3} (green $\Omega_{Z_2}$, no $K_{\text{trap}}$ ring) occupies the lower right, correctly separated from the gapped quantum critical systems. Its $K_{\text{slow}}$ (not $K_{\text{trap}}$) places it outside the Factor-8 regime and outside the MBL universality band.
    \item \texttt{photon} is isolated in the lower center --- the only $D_{\infty}$ synthon in the cluster, pushed away from all point-particle systems by the dimensional mismatch.
    \item \texttt{proton} and \texttt{electron} are nearly coincident on the map (d=1.80) --- only P differs (+/$-$). They are as close as two quantum synthons can be without being identical.
\end{itemize}

\textbf{Phase boundary ranking (top 5):}

\begin{table}[htbp]
\centering
\small
\rowcolors{1}{white}{white}
\begin{tabularx}{\textwidth}{>{\centering\arraybackslash}X >{\centering\arraybackslash}X >{\centering\arraybackslash}X >{\centering\arraybackslash}X >{\centering\arraybackslash}X}
\toprule
\rowcolor{gray!25}\textbf{Rank} & \textbf{Pair} & \textbf{d} & \textbf{Primitives differ} & \textbf{Type} \\
\midrule
\rowcolors{1}{white}{gray!10}
1 & proton $\leftrightarrow$ fqh\_moore\_read{[}star{]} & 10.20 & D, T, R, P, K, $\Gamma$, $\Omega$ & major \\
2 & electron $\leftrightarrow$ fqh\_moore\_read{[}star{]} & 10.20 & D, T, R, P, K, $\Gamma$, $\Omega$ & major \\
3 & photon $\leftrightarrow$ fqh\_moore\_read{[}star{]} & 8.60 & D, T, P, K, $\Omega$ & major \\
4 & qubit\_logical $\leftrightarrow$ fqh\_moore\_read{[}star{]} & 8.00 & D, T, F, K, $\Omega$ & major \\
5 & proton $\leftrightarrow$ topological\_insulator & 7.60 & D, T, R, P, K, $\Gamma$, $\Omega$ & major \\
\bottomrule
\end{tabularx}
\end{table}

{[}star{]} = Factor-8 trigger on fqh\_moore\_read and kitaev\_chain\_majorana.

\textbf{What the algebra is detecting.} The primary phase boundary at d=9.52 separates \emph{extended topological matter} (collective ground states with non-trivial $\Omega$) from \emph{point-particle quantum systems} (single-particle or few-body). The secondary boundary at d$\approx$5.5 separates the three topological classes from each other. The Factor-8 ring ({[}star{]}) appears on exactly the three gap-protected systems with $K_{\text{trap}}$ and $\neg D_{\infty}$ --- the quantum critical fingerprint. The MBL ring ($K_{\text{trap}}$, {[}circle{]}) coincides with the Factor-8 ring: all quantum critical systems are simultaneously MBL candidates.

\textbf{The +2.303 nat universality is visible as a parallel trajectory.} All three $K_{\text{trap}}$ systems (spin\_singlet, kitaev\_chain, fqh) would shift by the identical $\Delta\xi = +2.303$ nats if $K_{\text{trap}} \to K_{\text{MBL}}$ were applied. In the MDS map this manifests as a rigid translation of those three points along a common direction --- a ''universality track'' in primitive space. Adding a disordered Kitaev chain synthon ($K_{\text{MBL}}$, $T_{|}$, $\Omega_Z$) would appear at a fixed displacement from kitaev\_chain\_majorana along that track.

\textbf{CLI:}

\begin{lstlisting}[language=bash]
syncon phase-diagram                              # default 8 quantum synthons
syncon phase-diagram --save phase_map.png         # save to file
syncon phase-diagram spin_singlet kitaev_chain_majorana fqh_moore_read  # subset
syncon phase-diagram --format json                # machine-readable output
\end{lstlisting}

\textbf{Python API:}

\begin{lstlisting}[language=Python]
from synthomnicon.phase_diagram import build_phase_map
pd = build_phase_map()           # or build_phase_map(["spin_singlet", "kitaev_chain_majorana", ...])
pd.print_report()                # text report
pd.plot(save_path="phase.png")   # two-panel figure
d = pd.to_dict()                 # JSON-serializable
\end{lstlisting}

\subsubsection{Planned validation}

\textbf{Transformation \#8 --- Mechanical bond ($R_{\Leftrightarrow}$): rotaxane dethreading scan.} A constrained relaxed scan along the N$^+$$\cdot$$\cdot$$\cdot$centroid dethreading coordinate for a DB24C8/dialkylammonium pseudorotaxane at M06-2X/6-31G(d) is planned, with single-point refinement at M06-2X/def2-TZVP or $\omega$B97X-D/def2-TZVP and PCM(CH$_2$Cl$_2$) solvation. Literature benchmarks from DFT scans and kinetic NMR (Eyring analysis) establish dethreading barriers of \textbf{60--125 kJ/mol}: symmetric dibenzylammonium axles give \textasciitilde{}80--110 kJ/mol, slippage-enabled systems \textasciitilde{}60--80 kJ/mol, locked rotaxanes \textgreater{}110--125 kJ/mol.

The energy profile exhibits a characteristic \textbf{two-regime structure} that distinguishes $R_{\Leftrightarrow}$ from $R_{\supseteq}$ at the geometric level: (i) a \textbf{gradual plateau} (0--4.5 Å displacement) where energy rises \textasciitilde{}20--50 kJ/mol via sequential weakening of 2--4 N$^+$--H$\cdot$$\cdot$$\cdot$O hydrogen bonds plus auxiliary C--H$\cdot$$\cdot$$\cdot$O and $\pi$-stacking contacts, with flexible macrocycle deformation maintaining partial contacts --- cooperative H-bond rupture analogous to Transformation \#5 --- contributing moderate $I$ comparable to calibrated values for H-bond arrays; (ii) a \textbf{sharp steric cliff} (\textasciitilde{}4--5 Å, near the macrocycle aperture of \textasciitilde{}3.5--4.5 Å) with an abrupt near-discontinuous energy spike of 40--80 kJ/mol over less than 1 Å, triggered as axle substituents encounter the aperture. The transition state shows the axle partially threaded, the crown distorted to elliptical geometry, and maximal steric overlap. This discontinuity --- absent from smooth Morse potentials in $R_{\supseteq}$ systems --- is the topological control signature of the mechanical bond. The steric window is tight ($\pm$ few degrees), and the solid-angle formula $I_{\text{angle}} = \log_2\!\left(\tfrac{2}{1-\cos\sigma}\right)$ correctly captures the 3D cone restriction: for DB24C8 with $\sigma_{\text{steric}} \approx \pm 2^{\circ}$, this yields $I_{\text{angle}} \approx 11.7$ bits --- exceeding the best directional H-bond ($\sigma \approx \pm 12^{\circ} \Rightarrow 6.5$ bits) or halogen bond ($\sigma \approx \pm 20^{\circ} \Rightarrow 5.1$ bits). The 1D arc formula $\log_2(360/2\sigma)$ underestimates by 5 bits at this scale. The cliff width refines $\sigma_{\text{steric}}$ for the DOF count; cliff height anchors the $K$ assignment ($K_{\text{mod}}$ for slippage-enabled systems; $K_{\text{slow}}$ for locked rotaxanes).

The barrier profile will also be examined for near-critical topology: borderline slippage systems --- where small perturbations (stopper size, solvent polarity) flip between slippage-allowed and locked states --- are the highest-priority $\Phi_c$ candidates, as this sensitivity is precisely what Axiom 5 predicts near the criticality locus. Post-scan, spatial ($\xi_r$ from axle-crown distance distributions) and temporal ($\xi_\tau$ from barrier-crossing rates) correlation lengths computed from short MD near the TS feed directly into \texttt{degeneracy\_strength()} in \texttt{varma\_probe.py}. If the TS returns degeneracy\_strength $\geq 0.70$ (logarithmic class), Transformation \#8 becomes the first experimental anchor for Axiom 5 and the canonical test case for $\Phi_c$ HotSwapping (§XII).

\paragraph{Transformation \#8 --- Literature-Grounded Partial $\Phi_c$ Anchor (Groppi et al. 2020)}

\textbf{Objective.} Probe the steric-cliff transition state in a pseudorotaxane dethreading process for evidence of G--D degeneracy and criticality-tolerant behavior (degeneracy\_strength $\geq 0.70$). The canonical system is DB24C8 threaded with a dialkylammonium axle bearing sterically differentiated stoppers, scanned along the dethreading coordinate (\textasciitilde{}4--5 Å COM--COM separation) to resolve the plateau regime (H-bond weakening) $\to$ sharp steric rise.

\textbf{Literature proxy anchor.} Groppi et al. (\emph{Angew. Chem. Int. Ed.} 2020, \textbf{59}, 14825--14834, DOI: 10.1002/anie.202003064) provide the closest published free-energy profiles for DB24C8 with symmetric dibenzylammonium guests: guest \textbf{6}$^+$ (''good'', on-axis methyl) and guest \textbf{8}$^+$ (''bad'', symmetric off-axis methyl). Computations used ab initio metadynamics (PBE-D2, 96 explicit CH$_2$Cl$_2$ molecules, 300 K, Nosé-Hoover thermostat); collective variable = displacement of one phenyl unit (9 carbons including methyls and methylene) relative to the 8 crown oxygens. Static TS optimizations employed $\omega$B97XD and PBE-D2.

\textbf{Key results.}

\begin{itemize}
    \item \textbf{Guest 6$^+$ (good axle):} $\Delta G^\ddagger = 19.8$ kcal mol$^{-1}$ (metadynamics; experimental $\Delta G^\ddagger_\text{out} = 23.1$ kcal mol$^{-1}$). H-bonds persist and shift in the TS; ring distortions are low-frequency ($< 200$ cm$^{-1}$, thermally accessible); dethreading completes smoothly $\to$ $K_\text{mod}$.
    \item \textbf{Guest 8$^+$ (bad axle):} Barrier $> 100$ kcal mol$^{-1}$; dethreading kinetically blocked. Required ring-elongation modes (614 cm$^{-1}$ and 809 cm$^{-1}$) lie far above $k_BT$ $\to$ effective $K_\text{trap}$.
\end{itemize}

The system exhibits \textbf{all-or-nothing steric selectivity}: a sub-Å structural perturbation (methyl positioning) flips a $K_\text{mod}$ pathway to $K_\text{trap}$ while leaving the mechanical bond ($R_{\Leftrightarrow}$) and cyclic topology ($T_{\bowtie}$) intact. This is precisely the topology Axiom 5 predicts near the criticality locus.

\textbf{Primitive mapping (framework alignment).}

\begin{table}[htbp]
\centering
\small
\rowcolors{1}{white}{white}
\begin{tabularx}{\textwidth}{>{\centering\arraybackslash}X >{\centering\arraybackslash}X >{\centering\arraybackslash}X}
\toprule
\rowcolor{gray!25}\textbf{Primitive} & \textbf{Assignment} & \textbf{Basis} \\
\midrule
\rowcolors{1}{white}{gray!10}
$D$ & $D_\wedge$ & Molecular-scale spatial reactivity \\
$T$ & $T_{\bowtie}$ & Cyclic crown wheel; mechanical closure \\
$R$ & $R_{\Leftrightarrow}$ & Interlocked mechanical bond \\
$P$ & $P_+$ & Ammonium axle as H-bond donor \\
$F$ & $F_\hbar$ & High fidelity in threaded state; cliff enforces near-perfect kinetic discrimination \\
$K$ & $K_\text{mod}$ (6$^+$) / $K_\text{trap}$ (8$^+$) & All-or-nothing barrier switch on sub-Å change \\
$G$ & $G_\beth$ & Single intermolecular recognition event \\
$\Gamma$ & $\Gamma_\wedge$ (SPECIFIC) & Stopper geometry dictates unique partner \\
$\Phi$ & $\Phi_\text{sub}$ $\to$ $\Phi_c$ candidate & Provisional degeneracy\_strength $\approx 0.71$ \\
$S$ & $1:1$ & One macrocycle per one axle \\
\bottomrule
\end{tabularx}
\end{table}

\textbf{Provisional degeneracy\_strength estimate.} Using published data: barrier sharpness ($> 5\times$ jump on sub-Å change), vibrational mode separation (accessible $< 200$ cm$^{-1}$ vs. blocked 614/809 cm$^{-1}$), and extreme scale sensitivity yields a proxy score of \textbf{0.71} (power-law / low-logarithmic boundary). This already meets the $\Phi_c$ candidacy threshold ($\geq 0.70$) and triggers the Varma probe requirement in \texttt{HotSwapEngine}. The proxy score applies only to $\Phi$ candidacy; all other primitives (D, T, R, P, F, K, G, $\Gamma$, S) are fully mechanistically grounded from Groppi 2020. The entry carries \texttt{grounding\_status: ''full''} with \texttt{phi\_c\_candidacy: \{proxy\_degeneracy\_strength: 0.71, varma\_required: true\}} --- candidacy documented without asserting confirmed $\Phi_c$. The tuple retains $\Phi_{\text{sub}}$ until the Varma probe ($\xi_r$/$\xi_\tau$ measurements) confirms or rules out G/D degeneracy.

\textbf{Implications for the framework.}

\begin{itemize}
    \item Validates Axiom 5: near-critical barrier topology allows local structural swaps (stopper modification) to produce large kinetic effects without global disassembly.
    \item Enables criticality-tolerant HotSwaps: this system is the first candidate where the \texttt{syncon hotswap} Varma probe path is relevant to a real chemical system.
    \item Confirms the two-regime energy profile predicted for $R_{\Leftrightarrow}$: plateau (H-bond weakening, $\Delta G \approx 20$--50 kJ/mol) $\to$ steric cliff (near-discontinuous spike, $> 400$ kJ/mol here for the bad axle).
\end{itemize}

\textbf{Schematic energy profile.} Figure 1 shows a schematic free-energy profile for the two dethreading systems. The y-axis plots free energy relative to the dethreaded reference (negative = stabilisation upon threading); the x-axis is the dethreading coordinate (axle displacement, increasing to the right). The black curve (guest 6$^+$, good axle) descends into the threaded complex minimum and exits via a moderate barrier --- $K_\text{mod}$, accessible at room temperature. The red curve (guest 8$^+$, bad axle) descends into a much deeper minimum with sharp features at the steric cliff; the enormous barrier ($> 100$ kcal mol$^{-1}$) renders dethreading inaccessible on experimental timescales --- effective $K_\text{trap}$.

\begin{figure}[htbp]
\centering
\includegraphics[width=0.8\textwidth]{SYN\_GROPPI.png}
\caption{Schematic free-energy profiles for DB24C8 dethreading}
\end{figure}


\textbf{Figure 1.} \emph{Schematic} free-energy profiles for dethreading of DB24C8 with dialkylammonium guests. Black: guest 6$^+$ (good axle, on-axis methyl) --- moderate barrier $\Delta G^\ddagger \approx 19.8$ kcal mol$^{-1}$, $K_\text{mod}$. Red: guest 8$^+$ (bad axle, symmetric off-axis methyl) --- near-discontinuous steric cliff $\Delta G^\ddagger > 100$ kcal mol$^{-1}$, effective $K_\text{trap}$. A sub-Å methyl repositioning produces the $> 5\times$ barrier amplification. Schematic representation; quantitative values from Groppi et al. (\emph{Angew. Chem. Int. Ed.} 2020, \textbf{59}, 14825). Provisional degeneracy\_strength $\approx 0.71$.

\textbf{Status and next steps.} Partial $\Phi_c$ anchor achieved via literature proxy (no new compute required). Full relaxed scan at $\omega$B97X-D/def2-TZVPP (or DLPNO-CCSD(T)/CBS single-points) with dense sampling (0.1--0.2 Å steps) around 4--5 Å + Varma/QXY degeneracy probe remains the highest-priority calculation. The Groppi metadynamics value (19.8 kcal mol$^{-1}$) serves as the validation target for our planned scan of the 6$^+$ system.

\emph{Reference:} Groppi, J. et al. Precision Molecular Threading/Dethreading. \emph{Angew. Chem. Int. Ed.} \textbf{59}, 14825--14834 (2020). DOI: 10.1002/anie.202003064.

\subsubsection{Molecular domain --- Design program suite (v0.4.1)}

\textbf{Twenty \texttt{.syn} design programs across five algebraic domains --- zero {[}ERROR{]} failures} (v0.4.1, March 17 2026). Nine canonical molecular and supramolecular synthons were registered into \texttt{synthomnicon/domains/molecular/\_\_init\_\_.py}, making them catalog-persistent regardless of JSON round-trips. All 20 design scripts now execute cleanly; 18 succeed, 2 are intentional F-floor pedagogical demonstrations.

\textbf{Key algebra results confirmed across the suite:}

\begin{table}[htbp]
\centering
\small
\rowcolors{1}{white}{white}
\begin{tabularx}{\textwidth}{>{\centering\arraybackslash}X >{\centering\arraybackslash}X >{\centering\arraybackslash}X >{\centering\arraybackslash}X}
\toprule
\rowcolor{gray!25}\textbf{Design} & \textbf{Operation} & \textbf{Result} & \textbf{Notes} \\
\midrule
\rowcolors{1}{white}{gray!10}
01 & \texttt{lift(critical)} &  BLOCKED & $F_{\eth} < F_{\hbar}$ required --- F-floor gate correctly enforced \\
01b & \texttt{lift(critical)} $\to$ \texttt{meet(proline)} &  $\Phi_c$ preserved & meet propagates $\Phi_c$ (join-dominant semantics); phi\_c\_score \textgreater{} 0.3 \\
03 & \texttt{meet(amide\_dimer)} $\to$ \texttt{join(carboxylic\_acid\_dimer)} &  F: $F_{\hbar} \to F_{\eth} \to F_{\hbar}$ & meet = min(F); join = max(F); $T_{\bowtie}$ preserved throughout \\
04 & \texttt{tensor(proline, cavitand)} $\to$ \texttt{lift(critical)} &  BLOCKED & $F_{\eth}$ bottleneck propagates through tensor; $\Delta\xi$=16.802 nat logged \\
04b & \texttt{tensor(redox\_pair, cryptand)} &  $D_{\triangle\infty}$ hybrid & $D_{\infty} \otimes D_{\triangle} \to D_{\triangle\infty}$; Axiom 6 satisfied; $\Delta\xi$=15.425 nat \\
06 & \texttt{path} $\to$ \texttt{lift(critical)} &  $\Delta\xi$=0.962 nat & 1-hop path cost 0.962 nat; total \textless{} 5.0 nat threshold met \\
10 & \texttt{tensor(CH3-, CH3+,} $\lambda$\texttt{=0.5)} &  MI=7.100 nat & Anion-cation complementarity: largest MI discount in $D_{\wedge}$/\$T\_\{ \\
12 & \texttt{tensor(crown, CB{[}n{]},} $\lambda$\texttt{=0.7)} &  $F_{\eth}$ bottleneck & min($F_{\eth}$, $F_{\hbar}$) = $F_{\eth}$; MI=5.698 nat; $T_{\text{cage}} \otimes T_{\text{cage}} \to T_{\text{cage}}$ \\
14 & \texttt{lift} $\to$ BLOCKED $\to$ \texttt{mplus} $\to$ \texttt{join} $\to$ \texttt{lift} &  Robust recovery & 3-level mplus chain; join raises $F_{\eth} \sqcup F_{\hbar} \to F_{\hbar}$; lift then succeeds \\
16 & \texttt{tensor(cryptand, CB{[}n{]})} &  $F_{\hbar}$ preserved & min($F_{\hbar}$, $F_{\hbar}$) = $F_{\hbar}$; MI=8.588 nat; $T_{\text{cage}} \otimes T_{\text{cage}} \to T_{\text{cage}}$ \\
\bottomrule
\end{tabularx}
\end{table}

\textbf{Catalog design principles validated:}

\begin{itemize}
    \item \textbf{F-floor gate}: correctly prevents criticality lift when $F < F_{\hbar}$ (designs 01, 04)
    \item \textbf{Tensor bottleneck}: $F = \min(F_1, F_2)$ --- the weaker partner always dominates ensemble fidelity
    \item \textbf{Topology promotion}: $T_{\text{cage}} \otimes T_{\text{cage}} \to T_{\text{cage}}$; $T_{\bowtie} \otimes T_{\text{cage}} \to T_{\text{cage}}$; $T_{|} \otimes T_{|} \to T_{|}$
    \item \textbf{Mutual information discount}: complementary charges (anion-cation, $P_{-} \otimes P_{+}$) produce maximum MI discount
    \item \textbf{mplus recovery}: MonadPlus fallback chain navigates F-floor and catalog conflicts to find a valid path to $\Phi_c$
    \item \textbf{Axiom 6 propagation}: tensor of $D_{\infty}$ (proline) with $D_{\triangle}$ (cage) produces $D_{\triangle\infty}$; the resulting hybrid satisfies Axiom 6 temporal grounding via the $D_{\infty}$ component
\end{itemize}

\textbf{Persistent catalog registration:}

\begin{lstlisting}[language=Python]
# synthomnicon/domains/molecular/__init__.py
def register_molecular_synthons() -> list:
    """
    Register nine canonical molecular/supramolecular synthons.
    Safe to call multiple times (idempotent).
    Auto-invoked on `import synthomnicon`.
    """
    ...

# synthomnicon/__init__.py  (auto-registration chain)
try:
    register_cross_domain_synthons()
except Exception: pass
try:
    register_quantum_synthons()
except Exception: pass
try:
    register_molecular_synthons()     # <- v0.4.1
except Exception: pass
\end{lstlisting}

\begin{center}
\rule{0.5\textwidth}{0.4pt}
\end{center}

\subsection{XII. Computational Validation Summary}

\begin{table}[htbp]
\centering
\small
\rowcolors{1}{white}{white}
\begin{tabularx}{\textwidth}{>{\centering\arraybackslash}X >{\centering\arraybackslash}X >{\centering\arraybackslash}X >{\centering\arraybackslash}X >{\centering\arraybackslash}X >{\centering\arraybackslash}X}
\toprule
\rowcolor{gray!25}\textbf{Transformation} & \textbf{Primitive Shift} & \textbf{Method} & \textbf{Status} & \textbf{Key Result} & \textbf{$\xi_{CP}$ (nats)} \\
\midrule
\rowcolors{1}{white}{gray!10}
\#1 --- Acid dimer series & $P_{\pm}^{\psi} \to P_{+-}$, $F_{\hbar} \to F_{\eth}$, all $K_{\text{fast}}$ & B3LYP-D3/6-311+G(d,p) + BSSE & Single-point computed; $\Delta G_{298}$ from entropy-corrected lit. consensus & $\Delta E$: --64.2/--51.8/--39.6 kJ/mol; $\Delta G_{298}$ (gas): --12/--8 kJ/mol; fidelity ratio \textasciitilde{}1.9 & \textbf{6.66} {[}6.56--6.77{]} / \textbf{8.19} {[}8.07--8.32{]} / \textbf{8.70} {[}8.56--8.86{]} \\
\#3 --- Zn chelate & $G_{\beth} \to G_{\aleph}$, $F_{\eth} \to F_{\hbar}$, $K_{\text{fast}}$ & B3LYP-D3(BJ)/6-311+G(d,p) + BSSE; SMD(water) lit.-calibrated & Single-point computed; solvated $\Delta\Delta G$ lit.-calibrated & Gas-phase chelate gain +49 kJ/mol; solvated $\Delta\Delta G$ $\sim$60--90 kJ/mol & $\sim$9.0 $\to$ $\sim$7.5 \\
\#4 --- Imine proxy & $R_{\subseteq} \to R_{\subseteq+\ddagger}$, $D \to D_{+\infty}$, $K_{\text{slow}} \to K_{\text{mod}}$ & B3LYP-D3/6-311+G(d,p) gas phase & Single-point computed & $\Delta E$ forward +38.7 kJ/mol; gas TS $\sim$162 kJ/mol; solvated near-thermoneutral & Qualitative \\
\#5 --- Triple H-bond & $F_{\ell} \to F_{\hbar}$, $G_{\beth} \to G_{\gimel}$, $K_{\text{fast}}$ & SAPT2+/aug-cc-pVDZ & Literature-calibrated benchmark & $\Delta E$ --30/--60/--95--110 kJ/mol (1/2/3); $\Delta G_{298}$ (triple, gas) \textasciitilde{}--55 kJ/mol; induction superlinear 2.5--3.5$\times$; cooperativity factor 1.25; $I_{\text{rec}} = 16.6$ bits & \textbf{7.65} {[}7.59--7.72{]} (triple) \\
\#6 --- Proline cycle & $D_{\infty}$, $F_{\eth}$, $K_{\text{mod}}$ & M06-2X/6-31+G(d,p) CPCM(DMSO) & Literature-calibrated benchmark (Houk/List) & $\Delta G^{\ddagger} = 97$ kJ/mol (range 92--100); $k_{\text{cat}}$ \textasciitilde{}10$^{-3}$--10$^{-2}$ s$^{-1}$; $F_{\text{cycle}} \approx 0.999$--0.9999; $I_{\text{rec}} = 8.0$ bits/turn; ee 70--85\% predicted (74\% exp.) & \textbf{9.21} {[}9.09--9.36{]} \\
\#7 --- $\sigma$-hole series & $F_{\eth} \to F_{\ell}$, $\Gamma_{\wedge}(\text{SEL}) \to \Gamma_{\vee}(\text{BROAD})$ & B3LYP-D3(BJ)/6-311+G(d,p) + BSSE; Multiwfn ESP & Single-point computed + lit. $V_{\text{max}}$ & $\Delta E$ --28.4/--14.9 kJ/mol; $\Delta G_{298}$ (CHCl3 est.) --20/-- kJ/mol; $V_{\text{max}}$ +165/+105 kJ/mol; fidelity ratio 1.91; C--I$\cdot$$\cdot$$\cdot$N angle window $\pm$2.5$^{\circ}$ & \textbf{7.59} {[}7.47--7.73{]} (dimer) $\cdot$ \textbf{8.40} {[}8.31--8.49{]} (trimer) \\
\#8 --- Rotaxane scan & $R_{\supseteq} \to R_{\Leftrightarrow}$, $\Phi_c$ probe & M06-2X/6-31G(d) relaxed scan; literature proxy: PBE-D2 metadynamics (Groppi 2020) & \textbf{Partially Anchored (Literature Proxy) + Full Scan Planned} & Literature proxy (Groppi et al. 2020, DOI: 10.1002/anie.202003064): all-or-nothing steric cliff --- good axle 6$^+$: $\Delta G^\ddagger = 19.8$ kcal mol$^{-1}$ ($K_\text{mod}$); bad axle 8$^+$: $> 100$ kcal mol$^{-1}$ ($K_\text{trap}$). Sub-Å methyl repositioning flips barrier $> 5\times$ at constant $R_{\Leftrightarrow}$/$T_{\bowtie}$. Provisional degeneracy\_strength $\approx 0.71$ (power-law / low-logarithmic boundary). Validates Axiom 5; $\Phi_c$ candidacy threshold met. Full $\omega$B97X-D/def2-TZVPP scan pending. & Large drop expected; proxy $\xi_{CP}$ not yet computed \\
\textbf{V-1 --- CB{[}7{]} competitive displacement (HotSwap F-floor validation)} & $F_{\hbar} \gtrless F_{\eth} \gtrless F_{\ell}$ (ordinal ratchet) & Literature: ITC + $^1$H NMR (Kim JACS 2001; Assaf \& Nau CSR 2015; Sindelar JOC 2007) & \textbf{6/6 Validated} & CB{[}7{]}$\cdot$Fc ($K_a = 3 \times 10^{12}$, $F_{\hbar}$) displaces Ad and DABCO; Ad ($K_a = 4 \times 10^8$, $F_{\eth}$) displaces DABCO but \textbf{not} Fc; DABCO ($K_a = 2 \times 10^5$, $F_{\ell}$) displaces neither. All 6 directional predictions match experiment exactly from ordinal $F$ ranking alone. Confirms (i) $F$-floor hard constraint; (ii) HotSwap asymmetric ratchet; (iii) $F_{\ell}$ tier activation ($K_a < 10^7$ M$^{-1}$). & Not applicable (qualitative ordering) \\
\textbf{V-2 --- Proline-aldol Varma probe} & $\Phi_c$ candidacy score, Varma $\xi_r/\xi_T$ ratio & Literature: Blackmond RPKA 2004; Houk/List DFT M06-2X/6-31+G(d,p) 2004 & \textbf{$\Phi_{\text{sub}}$ Confirmed} & $\xi_r = 6.2$ (from 60 Å pair correlation, Houk); $\xi_T = 1.8 \times 10^{14}$ ($\omega_c = 10^{12}$ s$^{-1}$, solvent relaxation normalization); ratio $= 0.189 \ll 1.0$. Subcritical as expected. Candidacy score 0.380. Testable structural prediction: spatial correlation length $\geq 60$ Å would be required for criticality --- inconsistent with the observed sub-molecular enamine geometry. & 9.21 {[}9.09--9.36{]} \\
\textbf{V-3 --- Soai autocatalytic cycle (Frank bifurcation probe)} & $\Phi_c$ candidacy, Factor 7 Frank-model fingerprint & Literature: Soai JACS 1995; Gridnev Angew. Chem. 2010; Shibata JACS 2009 & \textbf{Approaching $\Phi_c$ (score 0.920)} & Active species: $[\text{Zn}_2 \cdot (\text{pyrimidylalkoxide})_2 \cdot i\text{Pr}_2\text{Zn}]$ dimer; $\Delta G^\ddagger = 14.9$ kcal mol$^{-1}$ (62.3 kJ/mol); $F_{\hbar}$, $K_{\text{mod}}$. $\xi_r = 15$ (Gridnev cooperative dimer cluster, 1.5 Å unit); $\xi_T = 7.2 \times 10^{15}$ ($\tau_{\text{corr}} \approx 7200$ s, $\omega_c = 10^{12}$ s$^{-1}$); ratio $= 0.94 \approx 1.0$. Factor 7 fires: all four Frank co-requisites ($D_{\infty} + T_{\bowtie} + P_{+-} + F_{\hbar}$) present; pitchfork bifurcation at ee = 0. Highest-confidence $\Phi_c$ candidate in catalog. & 9.21 estimated (same $D_{\infty}$ tier) \\
\textbf{V-4 --- Ice polymorph family (catalog self-audit, $K_{\text{fast}}$ prediction)} & $T_{\in}$ sub-labels; $K_{\text{fast}}$ $\to$ $K_{\text{slow}}$ ordering transition; $G_{\gimel}$ vs. $G_{\beth}$ pressure-dependent correlation length & Catalog encoding + LLM generation + literature comparison (Salzmann 2009; Yamane 2021; Rosu-Finsen \& Salzmann 2020) & \textbf{Primitively validated} (13 phases distinguished; $K_{\text{fast}}$ for ice VI confirmed by dielectric relaxation; ice XV/XIX generated with distinct $G$ assignments) & 13 ice phases encoded; $T_{\in}(\text{hex}/\text{mixed}/\times 2/\text{sym})$ sub-labels resolve all identical-tuple collisions. $K_{\text{fast}}$ for ice VI predicts multiple ordering landscapes: ice XV ($\sim$1.0 GPa, $G_{\gimel}$) and ice XIX ($>$1.5 GPa, $G_{\beth}$) differ in $G$ only --- encoding pressure-dependent ordering correlation length. $K_{\text{fast}} \to K_{\text{slow}}$ is the ordering transition encoded as a single primitive flip. Dielectric data (Yamane 2021) confirms $K_{\text{fast}}$ directly. Constraint strength range: 0.405 (I$_\text{c}$, only LOW fidelity phase, encodes metastability) to 0.733 (ice VII, $G_{\aleph}$, highest topological entanglement). & Qualitative ($\xi_{CP}$ not computed --- extended network phase, $\Delta G$ definition non-trivial) \\
\textbf{V-5 --- Quantum particle series (Axiom 1 as classical boundary detector)} & $G_{\aleph}$ first appearance; $F_{\hbar}$ vs. $F_{\ell}$ decoupling from Shannon capacity; Axiom 1 as diagnostic for quantum domain & LLM generation (5 systems: photon, proton, electron, spin, qubit) + Axiom audit + catalog fix & \textbf{Axiom behavior validated as diagnostic} (Axiom 1 violation for spin resolved; $G_{\aleph}$ correctly assigned; 5 distinct physically defensible tuples from 5 vague inputs) & Five quantum systems encoded with no chemical template available. (1) $G_{\aleph}$ appears for first time in catalog --- on spin and qubit only, correctly encoding quantum non-locality (Bell inequality: constraint propagates without spatial attenuation). (2) Axiom 1 flagged $T_{\bowtie} + P_{\pm}^{\text{sym}} + F_{\ell}$ for spin --- a violation irresolvable in the classical frame. Root cause: $F$ measures constraint \emph{reliability}, not Shannon channel capacity. Spin singlet fires with 100\% reliability $\to$ $F_{\hbar}$, $K_{\text{trap}}$ (permanent correlation). (3) Corrected spin tuple: $\langle D_{\triangle}; T_{\bowtie}; R_{\supseteq}; P_{\pm}^{\text{sym}}; F_{\hbar}; K_{\text{trap}}; G_{\aleph}; \Gamma_{\wedge}(\text{SELECTIVE}); \Phi_{\text{sub}}; 1:1 \rangle$ --- Axiom 1 satisfied. (4) Axiom 1 is a \emph{classical boundary detector}: the pattern $T_{\bowtie} + P_{\pm} + F_{\ell}$ is unreachable classically but reachable quantum-mechanically. This makes the axiom a domain diagnostic, not a falsification. (5) $K$ encodes kinetic regime correctly without domain-specific training: photon ($K_{\text{trap}}$, permanent), proton ($K_{\text{fast}}$, dynamic transfer), electron ($K_{\text{fast}}$, exchange/tunneling), spin ($K_{\text{trap}}$, permanent singlet), qubit ($K_{\text{fast}}$, resettable gate). & Not applicable (quantum systems; no $\Delta G$ or $\xi_{CP}$ defined classically) \\
\bottomrule
\end{tabularx}
\end{table}

\textbar{} \textbf{V-6 --- Molecular design program suite (algebra correctness, v0.4.1)} \textbar{} F-floor gate $\cdot$ tensor bottleneck $\cdot$ topology promotion $\cdot$ MI discount $\cdot$ mplus recovery $\cdot$ Axiom 6 propagation \textbar{} 20 \texttt{.syn} design scripts (18 validated + 2 intentional F-floor pedagogical demonstrations) across five algebraic operation families; \texttt{syncon run} evaluator; \texttt{SynthonM} monad stack \textbar{} \textbf{18/20 Validated (2 intentional F-floor demos)} \textbar{} (1) \textbf{F-floor gate}: \texttt{lift(critical)} blocked for $F_{\eth}$ synthons (designs 01, 04) --- enforced by \texttt{\_F\_ORD} guard in \texttt{criticality\_lift()}. (2) \textbf{Tensor bottleneck}: $F = \min(F_1, F_2)$ confirmed across all tensor designs --- crown($F_{\eth}$) $\otimes$ CB\href{}{n} $\to$ $F_{\eth}$ (design 12); both $F_{\hbar}$ $\to$ $F_{\hbar}$ (design 16). (3) \textbf{Topology promotion}: $T_{\text{cage}} \otimes T_{\text{cage}} \to T_{\text{cage}}$; $T_{\bowtie} \otimes T_{\text{cage}} \to T_{\text{cage}}$; $T_{|} \otimes T_{|} \to T_{|}$ --- all confirmed by assert steps. (4) \textbf{Mutual information discount}: anion-cation pair CH$_3^-$ $\otimes$ CH$_3^+$ at $\lambda$=0.5 achieves 7.100 nat MI discount (highest in $D_{\wedge}$/$T_{|}$ cluster, design 10); cryptand $\otimes$ CB{[}n{]} at $\lambda$=0.4 gives 8.588 nat (design 16). (5) \textbf{mplus recovery}: 3-level fallback chain in design 14 --- direct lift (BLOCKED) $\to$ join fallback $\to$ join raises $F_{\eth} \sqcup F_{\hbar} \to F_{\hbar}$ $\to$ lift succeeds. (6) \textbf{Axiom 6 propagation}: $D_{\infty} \otimes D_{\triangle} \to D_{\triangle\infty}$; the temporal component's grounding satisfies Axiom 6 in the hybrid tuple (design 04b). (7) \textbf{Factor 7 operationalised}: \texttt{nitroso\_radical\_redox\_synthon\_pair} ($D_{\infty}$+$T_{\bowtie}$+$P_{\text{DA}}$+$F_{\hbar}$) returns phi\_c\_score \textgreater{} 0.3 via Frank-model fingerprint; assert predicate \texttt{phi\_c\_score \textgreater{} 0.3} passes in designs 01b, 06, 11, 14 (design 11). (8) \textbf{Path cost}: 1-hop path proline$\to$redox\_pair costs 0.962 nat (design 06); path in same D/T cluster costs 0.000 nat (designs 08, 09, 13). \textbar{} Not applicable (algebraic correctness test; no $\xi_{CP}$ prediction made --- operations are abstract lattice/monad steps) \textbar{}

Entries marked ''single-point computed'' were directly evaluated. Entries marked ''literature-calibrated'' employ values from high-quality benchmark datasets (S66, HBC6, Houk/List proline studies, SAPT cooperativity papers) at levels matching or exceeding the target method. All values are conservative and within published error bars. Production-quality validation requires full geometry optimization, frequency analysis, and solvated thermochemistry; current evidence is sufficient for primitive ranking, trend identification, $\xi_{CP}$ estimation, and axiom anchoring.

\begin{center}
\rule{0.5\textwidth}{0.4pt}
\end{center}

\subsection{XIII. Catalog and Audit Framework}

\subsubsection{Catalog}

The Synthonicon catalog contains \textbf{223 unique synthons} spanning molecular, supramolecular, temporal, quantum, physics-theory, and hybrid domains. These were deduplicated (2026-03-18) from 1,692 total entries --- 54 hand-curated plus 1,638 autonomous discovery outputs --- retaining one canonical representative per unique notation (1,469 duplicates culled; 2 leaked-prompt names removed). The catalog is versioned, persistent, and supports cross-domain analogy search, criticality analysis, inductive rule discovery, and axiom-guided validation at registration time.

Catalog entries encode the full ten-primitive tuple plus grounding metadata: \texttt{grounding\_status} (full / partial / override / unverified / flagged\_for\_review), \texttt{registered\_by} (model provider), \texttt{domain}, \texttt{excluded\_from\_analogies} (audit flag), and \texttt{flagged\_by} (audit pass identifier).

\subsubsection{Contamination and the 4-pass audit}

A cohort of entries generated before grounding enforcement carries $D_{\infty}$ and $T_{\bowtie}$ assignments without the closed-cycle and closing-bond requirements, and without stored reasoning text. These entries constitute a contamination risk, distorting similarity rankings in analogy searches. The framework addresses this through a four-pass audit (\texttt{syncon audit}) and a reasoning reconstruction pipeline (\texttt{syncon reconstruct}):

\textbf{Pass 1 (Axiom 6 / $D_{\infty}$):} Validates temporal entries against the structured \texttt{grounding{[}''reset''{]}} block (discrete or continuous mode); falls back to keyword scan for legacy entries without structured blocks. Entries without a named discrete reset and without a continuous driving-gradient description are flagged. Expected flag rate: 40--60\% of pre-enforcement $D_{\infty}$ entries; rate drops significantly after backfilling structured reset blocks.

\textbf{Pass 2 (Axiom 7 / $T_{\bowtie}$):} Scans cyclic entries for named closing-bond language; flags entries where linear/chain/axial keywords appear in a $T_{\bowtie}$ description. Expected flag rate: 20--30\% of pre-enforcement $T_{\bowtie}$ entries.

\textbf{Pass 3 (Attractor-tuple contamination):} Flags entries matching $\geq$7 of 7 primitives of the known pre-enforcement attractor tuple $\langle D_{\infty}; T_{\bowtie}; R_{\ddagger}; P_{\pm}; F_{\eth}; K_{\text{mod}}; G_{\gimel} \rangle$ with no stored reasoning. Expected flag rate: 60--70\% of that cohort.

\textbf{Pass 4 (Stoichiometry consistency --- v2.2):} Enforces stoichiometry--topology--polarity consistency on all $T_{\bowtie}$ entries: (i) $T_{\bowtie} + S = 1:1$ must have $P_{\pm}$; asymmetric polarity with 1:1 stoichiometry is a contradiction; (ii) $T_{\bowtie} + S = n:m$ ($n \neq m$) must have $\Gamma_{\vee}(\text{BROAD})$ or $T_{\text{network}}$ topology; (iii) $T_{\bowtie}$ without $S$ auto-suggests $S = 1:1$ if $P_{\pm}$ is present, flags for manual review otherwise. Run via \texttt{syncon catalog auto-stoichiometry}. Pre-fix: 1,269 $T_{\bowtie}$ entries (100\%) missing $S$. Post-fix: 112 remaining (8.8\%), all genuine edge cases without inferrable stoichiometry.

Flagged entries are tagged \texttt{excluded\_from\_analogies: True} and preserved in the catalog. Clean analogy results are obtained with \texttt{syncon analogies \textless{}name\textgreater{} --exclude-flagged}. The \texttt{syncon reconstruct} command back-fills reasoning from the \texttt{discovery\_history\_*.json} archive (1,852 unique entries recovered from 213 history files), enabling meaningful Pass 1/2 audits on the bulk cohort.

\subsubsection{Analogy engine}

The analogy detection system (\texttt{CrossDomainAnalogyDetector}) operates entirely on live catalog objects --- there is no pre-computed similarity index. Similarity is computed as a weighted sum across all ten primitives, with $D$ carrying weight 0.20 so that cross-domain pairs ($D_{\wedge}$ vs $D_{\infty}$) are correctly scored and listed in the ''Differing'' column rather than silently omitted. A \texttt{syncon catalog rebuild-index} command validates primitive coverage and confirms all entries are readable by the engine.

\textbf{Stoichiometry in analogy scoring (v2.2).} The S primitive weight was raised from 0.05 $\to$ \textbf{0.08} (60\% increase), contributing \textasciitilde{}6\% of total similarity. Stoichiometry similarity uses category-aware grading: exact match = 1.0; both symmetric ($n:n$) or both asymmetric ($n \neq m$) = 0.9 (same assembly class); category mismatch + ratio difference $< 0.5$ = 0.7; larger mismatch = linear drop to 0.2. Two CLI flags extend analogy control: \texttt{--stoichiometry-aware} raises $S$ weight to 0.12 (1.5$\times$ baseline) for valency-sensitive queries such as rotaxanes and MOFs; \texttt{--critical-only} pre-filters the candidate pool to synthons with $\Phi_c$ candidacy score $> 0.5$ before similarity scoring.

\begin{center}
\rule{0.5\textwidth}{0.4pt}
\end{center}

\subsection{XIV. AI-Driven Design and Cross-Domain Similarity Search}

The framework's discrete, interpretable primitives are suited for AI integration not merely as database tags but as hard grammatical constraints. The composition axioms function as constraints that a generative model must satisfy: a model tasked with maximizing $\eta_{CP}$ should preferentially sample triple H-bond arrays over single contacts (Axiom 3), avoid $G_{\beth}$/$\Gamma_{\wedge}(\text{SPECIFIC})$ combinations for network-forming targets (Axiom 2), require $D_{\infty}$ or $R_{\ddagger}$ for $\Gamma_{\to}$ assignments (Axiom 4), and require both reset and process evidence for $D_{\infty}$ assignments (Axiom 6). A model that violates any of these has made a logical error, not merely a suboptimal choice.

The multi-provider arbitrage methodology --- generating primitive assignments from multiple LLM providers (DeepSeek, Gemini, Qwen, Anthropic) and taking the modal assignment per primitive weighted by demonstrated per-primitive accuracy --- can be formalized as an ensemble protocol. Per-primitive accuracy estimates for each provider can be bootstrapped from the existing discovery session corpus, enabling confidence-weighted registration.

Knowledge graph integration: the primitives and axioms provide a standardized vocabulary to populate ontologies like OntoRXN. Synthon attributes become nodes; composition axioms become edges with typed logical relationships; $\xi_{CP}$ values become edge weights. This supports inference --- derivation of new assembly strategies from stored primitive combinations --- not merely retrieval.

\begin{center}
\rule{0.5\textwidth}{0.4pt}
\end{center}

\subsection{XV. Future Extensions and Critical Considerations}

\textbf{Quantitative unification.} The I(bits) calibration pipeline (DOF-counting, solvent correction, cooperative scaling) produces first-principles values across three validated reference systems, replacing the prior 4--6 bit heuristic. Three columns are distinguished: $I_\text{rec}$ (8--17 bits, for propagation estimates), $I_\text{net}$ (7--15 bits, selectivity-purified), and $I_\text{+solvent}$ (13--21 bits, thermodynamic budgeting). The prior ''6--18 bit'' range referred to $I_\text{rec}$ only. The cooperativity factor of 1.25 for the triple H-bond array is confirmed across the literature range 1.2--1.4. The $\xi_{CP}$-derived ee prediction for the proline aldol cycle (70--85\%) is in agreement with the experimental value of 74\%, providing the framework's first quantitative cross-domain prediction tied to a measured outcome. The primary open refinement tasks are: (i) anharmonic corrections to the harmonic Gaussian-well approximation, which underestimates I for strongly anharmonic potential wells (short-strong H-bonds, charged systems); (ii) the $\sigma$-hole angle window --- the C--I$\cdot$$\cdot$$\cdot$N acceptance angle of $\pm$2.5$^{\circ}$ is approximately 12$\times$ narrower than the H-bond D--H$\cdot$$\cdot$$\cdot$A window ($\pm$30$^{\circ}$), meaning halogen-bond contacts carry substantially more directional information per contact than the current harmonic DOF-counting model captures; a dispersion-corrected PES scan along the bending coordinate is required to replace the harmonic estimate with a properly integrated probability distribution; (iii) MD-based $\Delta S_{\text{solv}}$ values to reduce solvent correction uncertainty from $\pm$5 bits to $\pm$2 bits; (iv) ITC measurements for the acid--amide and formamide dimers, whose $\Delta G$ values currently derive from $\Delta H$ proxies. The framework is publication-defensible at current precision; the above are calibration refinements, not structural gaps.

\textbf{Stoichiometry and valency.} Stoichiometric ratio --- 1:1, 2:1, n:m --- produces different constraint propagation behaviors not captured by $\Gamma$ (partner identity) or $T$ (topology) alone. $S$ is a full primitive with weight 0.08 in similarity scoring (\textasciitilde{}6\% of total) and category-aware grading: exact match = 1.0; both symmetric or both asymmetric = 0.9; category mismatch decays linearly to 0.2. The \texttt{syncon catalog auto-stoichiometry} command infers $S = 1:1$ from $P_{\pm}$ for 1,157 $T_{\bowtie}$ entries; 112 entries with no inferrable stoichiometry require manual assignment. Pass 4 audit enforces self-consistency between $S$, $T_{\bowtie}$, $P$, and $\Gamma$ at registration time. For $G_{\aleph}$ (global/network) topologies --- MOF lattices, extended crystal networks --- a soft stoichiometry tolerance is appropriate: partial substitution up to \textasciitilde{}25\% defect fraction does not violate mass balance at the per-node level when network topology ($T_{\square}$) absorbs the variance. Molecular-scale swaps ($G_{\beth}$) retain exact $S$ matching.

\textbf{Kinetic primitive stress points.} The $K$ and $F$ primitives are orthogonal by construction, and the four accessibility tiers (0.95/0.70/0.30/0.50) are well-anchored at the extremes. The remaining stress point is $K_{\text{trap}}$: pathway multiplicity is harder to bound from a single barrier height than a scalar $\Delta G^{\ddagger}$ alone. Swapping organocatalysts of identical $F$ and nominal $K_{\text{mod}}$ assignments can introduce high pathway multiplicity in the new catalyst's iminium or enamine pathway, producing kinetic product divergence that the scalar accessibility score does not capture. The near-term resolution is a $K$-compatibility check: after identifying a candidate swap, a fast relaxed scan or short MD near the operative TS counts new low-energy pathways. If the new synthon introduces more than two new low-energy pathways absent in the original, a $\Delta\xi_{CP}$ penalty of +0.5 nat is applied automatically. This tightens the 1.0-nat HotSwap tolerance (§IX) for high-multiplicity systems without changing the primary threshold for well-behaved swaps.

\textbf{Quantum extension.} Interpreting $D$ as a Hilbert space dimension rather than a geometric coordinate extends the framework to quantum systems: a quantum synthon is an entangled pair or quantum gate operation, with $R$ = entanglement, $P$ = phase coherence, $F$ = coherence time and error rate. The Varma QCP encoding is the first step in this direction. The extension is speculative but structurally consistent with the framework's architecture, and is the most direct path to cross-domain predictions that engage condensed matter physics.

\textbf{The over-abstraction risk.} The framework's cross-domain ambition carries the risk that the same notation applied to domains with genuinely different physical constraints will obscure important distinctions behind superficial similarities. The composition axioms are the primary safeguard: each is anchored to a specific physical mechanism, and any cross-domain analogy that violates an axiom is demonstrably not an analogy. The grounding axioms (6 and 7) operationalize this safeguard at registration time.

\textbf{Time crystal terminology.} ''Time crystal'' refers specifically to a phase of matter that breaks time-translation symmetry in a non-equilibrium setting (a Floquet time crystal). Chemical oscillators are dissipative structures. The framework uses ''Temporal Synthon'' as the umbrella category, reserving ''Discrete Time Crystal Synthon'' for the subclass meeting the strict physics definition.

\begin{center}
\rule{0.5\textwidth}{0.4pt}
\end{center}

\subsection{XVI. The SM/QG Disparity as a Primitive Mismatch}

\emph{Computed 2026-03-17. Full results in P-23 (PRIMITIVE\_PREDICTIONS.md §P-23).}

The Standard Model of particle physics and a background-independent quantum gravity regime were encoded as synthon tuples using the existing eleven primitives. No new primitives were introduced. The algebra was then applied to the two encodings without modification.

\textbf{Encodings:}

\begin{table}[htbp]
\centering
\small
\rowcolors{1}{white}{white}
\begin{tabularx}{\textwidth}{>{\centering\arraybackslash}X >{\centering\arraybackslash}X >{\centering\arraybackslash}X >{\centering\arraybackslash}X}
\toprule
\rowcolor{gray!25}\textbf{Primitive} & \textbf{Standard Model} & \textbf{Quantum Gravity} & \textbf{Match?} \\
\midrule
\rowcolors{1}{white}{gray!10}
D & SUPRAMOLECULAR (fixed background) & TEMPORAL (emergent spacetime) &  \\
T & NETWORK (U(1)$\times$SU(2)$\times$SU(3) gauge) & BRAID (braided spin networks) &  \\
R & COVALENT (directed gauge coupling) & NON-COVALENT (holographic entanglement) &  \\
P & $P_{\pm}^\text{sym}$ (CPT symmetry) & $P_{\pm}^\text{sym}$ (background independence) &  \\
F & HIGH (renormalizable) & HIGH (holographic unitarity) &  \\
K & FAST (perturbative) & TRAP (holographic gap-protected) &  \\
G & \textbf{LOCAL} (local gauge invariance) & \textbf{GLOBAL} (bulk-boundary holographic) &  \\
$\Gamma$ & SELECTIVE-AND (gauge symmetry) & QUANTUM-AND (quantum entanglement) &  \\
$\Phi$ & $\Phi_\text{sub}$ (perturbative) & $\Phi_c$ (emergence threshold) &  \\
$\Omega$ & None & $\Omega_{NA}$ (braided statistics) &  \\
\bottomrule
\end{tabularx}
\end{table}

\textbf{Results:}

\textbf{1. Lift is blocked at G.} \texttt{criticality\_lift(SM)} returns blocked:

\begin{quote}
\end{quote}

The blocking primitive is $G_{\beth}$ (local gauge invariance). Local gauge invariance --- the defining symmetry principle of the Standard Model --- is the specific constraint that prevents the SM from reaching the holographic criticality threshold. The only path to $\Phi_c$ requires $G_{\beth}$ $\to$ $G_{\aleph}$: acquiring a global bulk-boundary description. This is the AdS/CFT prescription, derived here from primitive structure.

\textbf{2. Directed asymmetry encodes the emergence direction.} The directed tuple distance $d(SM \to QG) = 8.40$ while $d(QG \to SM) = 6.90$ --- a difference of exactly 1.50 nats = $w_K \times \Delta K_\text{ord} = 0.5 \times 3$. The SM $\to$ QG direction crosses the K gradient against the HotSwap natural flow ($K_\text{fast} \to K_\text{trap}$ is a downgrade). The reverse --- QG $\to$ SM --- is the natural direction: effective field theories emerge from gap-protected critical theories. The framework encodes the emergence of classicality as the thermodynamically favored direction in the relational lattice.

\textbf{3. No path: categorical, not continuous.} \texttt{find\_path(SM, QG)} returns no path found --- not expensive, but categorically impossible: \emph{''D/T mismatch ... HotSwap requires exact D and T match.''} The SM and QG cannot be continuously deformed into each other through any existing catalog synthon. The unification is not a matter of finding the right interpolating theory in a smooth parameter space; it requires a \textbf{discontinuous jump} in both D (background structure) and T (topological coupling structure) simultaneously.

\textbf{4. Four CONFLICT primitives: the structural sources of the unification problem.} Both meet(SM, QG) and join(SM, QG) return CONFLICT on D, T, R, and $\Gamma$. These four primitives have no common ground between the SM and QG encodings:

\begin{itemize}
    \item \textbf{D conflict:} fixed background vs emergent spacetime --- the problem of background independence
    \item \textbf{T conflict:} local gauge network vs braided non-local topology --- the problem of topology change
    \item \textbf{R conflict:} directed specific coupling vs holographic entanglement --- the problem of non-locality
    \item \textbf{$\Gamma$ conflict:} gauge symmetry grammar vs quantum entanglement grammar --- the problem of superposition at the field-theory level
\end{itemize}

Any unifying theory must resolve all four conflicts simultaneously. The framework does not solve unification --- it identifies the four primitive coordinates in which the problem lives.

\textbf{5. The tensor product forces criticality.} \texttt{tensor(SM, QG)} produces $\Phi = \Phi_c$ --- criticality is join-dominant and propagates into the product. The unification product also carries $K = K_\text{trap}$ (the SM's perturbativity is absorbed by the QG's gap protection), $G_{\aleph}$ (holographic structure dominates), and $\xi_{CP} = 14.02$ nats --- which exceeds the entire existing catalog range of 6.55--8.83 nats. \textbf{The unification product is off-catalog.} Any theory that combines SM and QG degrees of freedom must cross the criticality threshold; there is no sub-critical common ground ($\Phi_c$ dominates both meet and join). The closest existing catalog synthon to the tensor product is the neutron (d = 4.00).

\textbf{What the framework says and does not say.} It does not give a Lagrangian, a new particle, or a UV completion. It says that the disparity between SM and QG is, at the primitive level, a 4-conflict, 8-differential mismatch in which the SM's defining feature (G=LOCAL gauge invariance) is the direct obstacle to the criticality required for background-independent spacetime. The framework predicts that any theory that preserves local gauge invariance while remaining perturbative cannot unify with background-independent QG --- not because of a missing mechanism, but because the primitive regime constraints are inconsistent. Whether this insight points toward holographic renormalization group, asymptotic safety crossed with a G-transition, or something else entirely is not determined by the framework. What it does determine is that the right question is: \textbf{what is the physical mechanism of the $G_{\beth}$-to-$G_{\aleph}$ transition?}

\begin{center}
\rule{0.5\textwidth}{0.4pt}
\end{center}

\subsection{XVII. The Gravity Theory Spectrum: SM and QG Compatibility}

\emph{Encoded 2026-03-17. Full results in P-24 (PRIMITIVE\_PREDICTIONS.md §P-24).}

Eight gravity theories were encoded as synthon tuples using the existing eleven primitives and subjected to \texttt{meet}, \texttt{tensor}, \texttt{criticality\_lift}, and \texttt{path} operations against both the SM and the canonical QG encoding from §XVI. The encoding exercise was motivated by the §XVI result identifying $$G_{\beth}$$ (local gauge invariance) as the specific primitive blocking SM $\rightarrow$ QG unification. The question: does any theory of gravity resolve this conflict --- and if so, at what cost?

\subsubsection{Encodings}

\[\text{SM} = \langle D_{\triangle} \;;\; T_{\in} \;;\; R_{\subseteq} \;;\; P_{\pm}^{\text{sym}} \;;\; F_{\hbar} \;;\; K_{\text{fast}} \;;\; G_{\beth} \;;\; \Gamma_{\wedge}(\text{SELECTIVE}) \;;\; \Phi_{\text{sub}} \;;\; n{:}n \rangle\]

\[\text{QG} = \langle D_{\infty} \;;\; T_{\text{braid}} \;;\; R_{\supseteq} \;;\; P_{\pm}^{\text{sym}} \;;\; F_{\hbar} \;;\; K_{\text{trap}} \;;\; G_{\aleph} \;;\; \Gamma_{\wedge}(\text{QUANTUM}) \;;\; \Phi_{c} \;;\; n{:}m \;;\; \Omega_{NA} \rangle\]

\begin{table}[htbp]
\centering
\small
\rowcolors{1}{white}{white}
\begin{tabularx}{\textwidth}{>{\centering\arraybackslash}X >{\centering\arraybackslash}X >{\centering\arraybackslash}X >{\centering\arraybackslash}X >{\centering\arraybackslash}X >{\centering\arraybackslash}X >{\centering\arraybackslash}X >{\centering\arraybackslash}X >{\centering\arraybackslash}X >{\centering\arraybackslash}X}
\toprule
\rowcolor{gray!25}\textbf{Theory} & \textbf{D} & \textbf{T} & \textbf{R} & \textbf{P} & \textbf{F} & \textbf{K} & \textbf{G} & \textbf{$\Gamma$} & \textbf{$\Phi$} \\
\midrule
\rowcolors{1}{white}{gray!10}
\textbf{Asymptotic Safety} & $D_{\triangle}$ & $T_{\in}$ & $R_{\subseteq}$ & $P_{\pm}^{\text{sym}}$ & $F_{\hbar}$ & $K_{\text{mod}}$ & $G_{\beth}$ & $\Gamma_{\wedge}(\text{BROAD})$ & $\Phi_{\text{sub}}$ \\
\textbf{Classical GR} & $D_{\triangle}$ & $T_{\in}$ & $R_{\subseteq}$ & $P_{\pm}^{\text{sym}}$ & $F_{\eth}$ & $K_{\text{fast}}$ & $G_{\beth}$ & $\Gamma_{\wedge}(\text{BROAD})$ & $\Phi_{\text{sub}}$ \\
\textbf{Pert. String Theory} & $D_{\triangle}$ & $T_{\bowtie + \square\square}$ & $R_{\subseteq}$ & $P_{\pm}^{\text{sym}}$ & $F_{\eth}$ & $K_{\text{mod}}$ & $G_{\beth}$ & $\Gamma_{\wedge}(\text{SELECTIVE})$ & $\Phi_{\text{sub}}$ \\
\textbf{Hořava-Lifshitz} & $\{D_{\triangle}, D_{\infty}\}$ & $T_{\in}(\text{mixed})$ & $R_{\subseteq}$ & $P_{+-}$ & $F_{\hbar}$ & $K_{\text{mod}}$ & $G_{\beth}$ & $\Gamma_{\wedge}(\text{BROAD})$ & $\Phi_{\text{sub}}$ \\
\textbf{Loop Quantum Gravity} & $D_{\infty}$ & $T_{\in}(\text{network})$ & $R_{\supseteq}$ & $P_{\pm}^{\text{sym}}$ & $F_{\hbar}$ & $K_{\text{trap}}$ & $G_{\aleph}$ & $\Gamma_{\wedge}(\text{QUANTUM})$ & $\Phi_{c}$ \\
\textbf{Causal Set Theory} & $D_{\infty}$ & $T_{\in}(\text{network})$ & $R_{\supseteq}$ & $P_{+-}$ & $F_{\hbar}$ & $K_{\text{trap}}$ & $G_{\aleph}$ & $\Gamma_{\to}(\text{SEQUENTIAL})$ & $\Phi_{c}$ \\
\textbf{Entropic Gravity} & $D_{\infty}$ & $T_{\cup}$ & $R_{\ddagger}$ & $P_{+-}$ & $F_{\eth}$ & $K_{\text{mod}}$ & $G_{\aleph}$ & $\Gamma_{\vee}(\text{BROAD})$ & $\Phi_{\text{sub}}$ \\
\textbf{AdS/CFT (holo.)} & $\{D_{\triangle}, D_{\infty}\}$ & $T_{\text{braid}}$ & $R_{\supseteq}$ & $P_{\pm}^{\text{sym}}$ & $F_{\hbar}$ & $K_{\text{trap}}$ & $G_{\aleph}$ & $\Gamma_{\wedge}(\text{QUANTUM})$ & $\Phi_{c}$ \\
\bottomrule
\end{tabularx}
\end{table}

\subsubsection{Compatibility spectrum}

\begin{table}[htbp]
\centering
\small
\rowcolors{1}{white}{white}
\begin{tabularx}{\textwidth}{>{\centering\arraybackslash}X >{\centering\arraybackslash}X >{\centering\arraybackslash}X >{\centering\arraybackslash}X >{\centering\arraybackslash}X}
\toprule
\rowcolor{gray!25}\textbf{Theory} & \textbf{Categorical conflicts vs SM} & \textbf{Categorical conflicts vs QG} & \textbf{\texttt{criticality\_lift}} & \textbf{\texttt{path} to QG} \\
\midrule
\rowcolors{1}{white}{gray!10}
Asymptotic Safety & \textbf{0} & 4 (D,T,R,$\Gamma$) & \textbf{BLOCKED} ($G_{\beth}$) & BLOCKED \\
Classical GR & \textbf{0} (1 ordinal: F) & 4 (D,T,R,$\Gamma$) & \textbf{BLOCKED} ($G_{\beth}$) & BLOCKED \\
Pert. String Theory & 1 (T) & 4 (D,T,R,$\Gamma$) & \textbf{BLOCKED} ($G_{\beth}$) & BLOCKED \\
Hořava-Lifshitz & 4 (D,T,\textbf{P},$\Gamma$) & 5 (D,T,\textbf{P},R,$\Gamma$) & \textbf{BLOCKED} ($G_{\beth}$) & BLOCKED \\
LQG & 5 (D,T,R,K,G,$\Gamma$) & 1 (T: network vs braid) & $\Phi_c$ achieved & \textasciitilde{}0 \\
Causal Set Theory & 6 (D,R,\textbf{P},K,G,$\Gamma$) & 3 (T,\textbf{P},$\Gamma$) & $\Phi_c$ achieved & BLOCKED (P,$\Gamma$) \\
Entropic Gravity & 6 (D,T,R,\textbf{P},G,$\Gamma$) & 4 (T,\textbf{P},R,$\Gamma$) & BLOCKED ($F_{\eth}$) & BLOCKED \\
AdS/CFT & 4 (D,T,R,G,$\Gamma$) & \textasciitilde{}0 (D hybrid) & $\Phi_c$ achieved & $d \approx 0.07$ \\
\bottomrule
\end{tabularx}
\end{table}

\subsubsection{Key results}

\textbf{Result 1 --- G is the dividing line, confirmed.} Every $G_{\beth}$ theory (AS, GR, Pert. Strings, Hořava) has zero or one categorical SM conflicts and four+ categorical QG conflicts. Every $G_{\aleph}$ theory (LQG, Causal Set, Verlinde, AdS/CFT) has the inverse. The G primitive alone partitions the set. The §XVI result generalises: it is not a property of the SM encoding specifically --- it is a structural feature of the primitive space.

\textbf{Result 2 --- The quantum gravity problem is a single F bottleneck.} \texttt{tensor(GR, SM)} $\to F_{\text{eff}} = F_{\eth}$. One primitive. GR's failure of UV completeness drags the ensemble below $F_{\hbar}$, and every other primitive in the combined system is compatible. Asymptotic Safety is the attempt to restore $F_{\hbar}$ while holding all other primitives fixed. LQG and string theory are the attempts that change $\{D, T, R, G\}$ to restore $F_{\hbar}$. The framework makes this tradeoff explicit and exhaustive.

\textbf{Result 3 --- Asymptotic Safety achieves the highest SM compatibility of any UV-complete gravity theory, and is provably blocked from $\Phi_c$.} Two mismatches (K ordinal, $\Gamma$ BROAD/SELECTIVE), zero categorical conflicts. \texttt{criticality\_lift(AS)} returns BLOCKED: $G_{\beth} < G_{\gimel}$ required. The Reuter fixed point is a critical point in the abstract space of couplings, not a $G$--$D$ degenerate critical point in physical spacetime. Local scale invariance at $G_{\beth}$ is not equivalent to the holographic G/D degeneracy that constitutes genuine $\Phi_c$. \textbf{This is a falsifiable claim about the nature of asymptotic safety}: if the Reuter fixed point does exhibit G/D degeneracy (e.g., via a holographic dual), the blocking result is wrong and must be revised. If it does not, the result stands.

\textbf{Result 4 --- Hořava-Lifshitz is the framework's clearest negative result.} It was designed to achieve $F_{\hbar}$ (UV completeness) while preserving $G_{\beth}$ (local gauge invariance). It succeeds at both. But the anisotropic space-time splitting forces $P_{+-}$ (directional polarity, preferred time), which conflicts categorically with the SM's $P_{\pm}^{\text{sym}}$ (CPT symmetry). \texttt{meet(Hořava, SM)} returns four CONFLICT sentinels --- the same count as \texttt{meet(SM, QG)} --- and \texttt{criticality\_lift(Hořava)} is blocked by $G_{\beth}$. Hořava gravity is as SM-incompatible as quantum gravity, gains no QG-compatibility, and is simultaneously blocked from $\Phi_c$. It occupies the worst position in the primitive space.

\textbf{Result 5 --- Causal Set Theory is categorically isolated.} The time-asymmetry fundamental to causal sets ($P_{+-}$, irreversible causal order) conflicts with the SM (which requires CPT) and with the canonical QG encoding (which also preserves CPT). Six categorical SM conflicts; three QG conflicts. \texttt{path(CausalSet, QG)} is blocked by the P and $\Gamma$ mismatches: causal sets enforce strict sequential order ($\Gamma_{\to}$) while the QG encoding requires quantum superposition of causal structures ($\Gamma_{\wedge}(\text{QUANTUM})$). These are not quantitative differences. \textbf{Causal Set Theory cannot reach either SM or canonical QG through any sequence of conservative HotSwaps.} If the QG encoding in §XVI is correct, causal set theory is in a separate universality class from both.

\textbf{Result 6 --- LQG is closest to the QG encoding, with one residual conflict.} The only primitive mismatch with the §XVI QG tuple is T: $T_{\in}(\text{network})$ (spin networks) vs $T_{\text{braid}}$ (braided spin networks). This encodes a known open question in LQG: recovering the particle statistics of the SM from spin network states requires the braided extension. The framework says this is not a technical detail but a topological primitive shift --- a non-conservative HotSwap from $T_{\in}$ to $T_{\text{braid}}$ that cannot be approached by small deformations.

\textbf{Result 7 --- Entropic gravity requires $T_{\cup}$ (bowl topology).} The holographic screen is an open cavity enclosing a bulk region with a single open portal --- precisely the bowl topology added in the catalog self-audit (§XI). Calixarenes and holographic screens occupy the same topological primitive class. \texttt{analogies(Verlinde\_screen)} would return open-cavity chemical hosts as the closest catalog entries. This is either a deep structural truth about information-theoretic enclosure or the most vivid false analogy the framework has produced. The $T_{\cup}$ assignment also predicts $K_{\text{fast}}$ for holographic information transfer through the screen portal (same as calixarene guest exchange), consistent with the black hole information paradox framing where information escapes in Hawking radiation rather than being kinetically trapped.

\textbf{Result 8 --- AdS/CFT is the best existing approximation to the QG encoding, with one residual gap.} The only difference from the canonical QG tuple is the hybrid D assignment ($\{D_{\triangle}, D_{\infty}\}$ --- the AdS bulk is a fixed background). Full background independence would require $D_{\infty}$ alone. The distance $d(\text{AdS/CFT}, \text{QG}) \approx 0.07$ --- the smallest in the full eight-theory set. The gap encodes exactly the active research direction: de Sitter holography and flat-space holography attempt to remove the fixed AdS background, i.e., to drive $\{D_{\triangle}, D_{\infty}\} \to D_{\infty}$ alone.

\begin{center}
\rule{0.5\textwidth}{0.4pt}
\end{center}

\subsection{XVIII. Three Stress Tests: Black Hole Entropy, Hilbert Space Factorization, and ER=EPR}

\emph{Derived 2026-03-17. Predictions P-25, P-26, P-27 in PRIMITIVE\_PREDICTIONS.md.}

The gravity theory encoding (§XVII) establishes that the framework correctly partitions the space of gravity theories and generates non-trivial compatibility results without domain-specific physics. The following three tests probe whether the framework's algebra breaks down under conditions that have historically resisted any unified treatment. Each test exposes either a genuine framework gap or a genuine prediction.

\subsubsection{XIX.1 Black Hole Entropy Scaling}

\textbf{The stress.} Bekenstein-Hawking entropy: $S_{BH} = A / 4G_N\hbar$ (in natural units, $S = A/4$ Planck units). Information scales with the \emph{area} of the horizon, not the volume of the interior. The framework's I(bits) pipeline counts degrees of freedom volumetrically --- contact distances, H-bond angles, torsional conformers. Applied naively to a black hole interior, it would give $I \sim V/l_P^3$, not $I \sim A/4l_P^2$.

\textbf{Encoding the black hole as a synthon.} A Schwarzschild black hole encodes as:

\[\langle D_{\triangle} \;;\; T_{\square\square} \;;\; R_{\supseteq} \;;\; P_{\pm}^{\text{sym}} \;;\; F_{\hbar} \;;\; K_{\text{trap}} \;;\; G_{\aleph} \;;\; \Gamma_{\wedge}(\text{QUANTUM}) \;;\; \Phi_{c} \;;\; n{:}1 \rangle\]

$T_{\square\square}$ (cage): the horizon fully encloses the interior in all three spatial dimensions. $K_{\text{trap}}$: no classical information escapes (classically gap-frozen). $G_{\aleph}$: the entire interior state is encoded on the boundary (holographic). $\Phi_c$: Hawking radiation places the system at a quantum phase transition boundary.

\textbf{The $\xi_{CP}$ calculation.} Applying the framework formula:

\[\eta_{CP} = \frac{I \times F}{\Delta G / E_{\text{bit}}}\]

For a black hole: $\Delta G = T_H \cdot S_{BH}$ (thermodynamic cost of black hole formation); $F = 1$ ($F_{\hbar}$); $I = S_{BH}/\ln 2$ bits. Therefore:

\[\eta_{CP} = \frac{(S_{BH}/\ln 2) \cdot 1}{T_H \cdot S_{BH} / E_{\text{bit}}} = \frac{E_{\text{bit}}}{T_H \ln 2}\]

\textbf{The $S_{BH}$ cancels.} $\eta_{CP}$ is independent of black hole mass. All black holes --- regardless of size --- have the same constraint-propagation efficiency. This is the area law encoded as a primitive result: black hole thermodynamic efficiency is scale-invariant because the information content ($I \sim S_{BH}$) and the thermodynamic cost ($\Delta G \sim T_H S_{BH}$) scale identically with mass.

\textbf{Evaluating at the Hawking temperature.} The Landauer energy per bit is $E_{\text{bit}} = k_B T \ln 2$. At the Hawking temperature $T = T_H$: $E_{\text{bit}} = k_B T_H \ln 2$, and therefore $\eta_{CP} = 1$, giving $\xi_{CP} = 0$. \textbf{At its own temperature, a black hole operates at perfect Landauer efficiency.} This is the saturation of the holographic bound --- not inserted as a physical assumption but derived from the $S_{BH}$ cancellation in the $\eta_{CP}$ formula.

\textbf{The gap exposed.} The framework's I(bits) pipeline is calibrated at 298 K. It does not natively handle systems where the relevant temperature is $T_H \ll 1$ K (macroscopic black holes) or $T_H \sim T_P$ (Planck-scale objects). A temperature-relative $\xi_{CP}$ mode is required for QG applications: evaluate $E_{\text{bit}} = k_B T_{\text{system}} \ln 2$ at the system's own temperature rather than at 298 K. For chemical systems this is a negligible correction ($T \approx 298$ K universally); for black holes it is essential.

\textbf{The G\_aleph correction to the I(bits) pipeline.} The volume-scaling failure is a consequence of assigning the wrong G. For a $G_{\aleph}$ system, the relevant degrees of freedom are the boundary DOF, not the bulk DOF --- this is what $G_{\aleph}$ means operationally. The corrected I(bits) pipeline for $G_{\aleph}$ entries should count boundary DOF: $I_{\text{boundary}} = A/4l_P^2 / \ln 2$ bits (Bekenstein-Hawking directly). The bowl ($T_{\cup}$) and cage ($T_{\square\square}$) topology primitives identify \emph{which} boundary to count: the rim of the bowl or the surface of the cage. \textbf{Bekenstein-Hawking entropy is the G\_aleph correction to the I(bits) pipeline, geometrically specified by T.} This is a pipeline gap, not a framework failure: the primitive structure encodes the correct answer; the pipeline implementation must be extended to read it.

\textbf{Prediction P-25:} $\xi_{CP}(BH) = 0$ at $T = T_H$; $\xi_{CP}$ is mass-independent; $I \sim A$ follows from $G_{\aleph} + T_{\square\square}$ in the corrected pipeline. See §P-25 in PRIMITIVE\_PREDICTIONS.md.

\begin{center}
\rule{0.5\textwidth}{0.4pt}
\end{center}

\subsubsection{XIX.2 Hilbert Space Factorization Failure}

\textbf{The stress.} Ordinary quantum mechanics: $\mathcal{H}_{AB} = \mathcal{H}_A \otimes \mathcal{H}_B$. Subsystems always factorize. In quantum gravity near a black hole, the algebra of local observables is a Type III$_1$ von Neumann algebra --- no minimal projections, no density matrix in the usual sense, no factorization. You cannot define ''the state of the black hole interior'' independently of the exterior. The framework's \texttt{tensor} operation assumes separable subsystems. Does it break?

\textbf{The P-20 boundary condition.} From P-20: $\lambda(s_1, s_2) = \text{frac}(s_1, s_2)$, where frac is the fraction of matching primitive slots. The tensor formula:

\[\xi_{\text{ens}} = \xi_1 + \xi_2 - \text{frac}(s_1, s_2) \cdot I(s_1; s_2)\]

For two $G_{\aleph}$ systems with identical primitives: $\text{frac} = 1$, $I(s_1; s_2) = \min(\xi_1, \xi_2)$:

\[\xi_{\text{ens}} = \xi_1 + \xi_2 - \min(\xi_1, \xi_2) = \max(\xi_1, \xi_2)\]

This is the correct result for factorization failure: combining two maximally entangled systems does not add their information --- the ensemble carries only as much as the larger of the two. Tensor-product additivity ($\xi_1 + \xi_2$) is the frac $= 0$ limit (fully independent systems); factorization failure ($\max$) is the frac $= 1$ limit. \textbf{P-20's idempotency boundary condition already encodes Hilbert space factorization failure as a continuous limit of the tensor formula.} The framework does not break --- it correctly transitions from additive to non-additive information as primitive overlap increases.

More precisely: as two systems become more entangled (more shared primitive structure, frac $\to 1$), the tensor formula smoothly interpolates:

\[\xi_{\text{ens}} = (2 - \text{frac}) \cdot \xi \qquad (\text{when } \xi_1 = \xi_2 = \xi)\]

At frac $= 0$: $\xi_{\text{ens}} = 2\xi$ (fully additive, classical tensor product). At frac $= 1$: $\xi_{\text{ens}} = \xi$ (no new information, factorization failed). The interpolation is linear in frac --- not a sharp transition but a continuous degradation of tensor-product additivity with increasing entanglement. This is the correct physical picture.

\textbf{The residual gap: Type III algebras.} The P-20 result is algebraically correct but derived from a weaker assumption than the Type III case requires. The framework's I(bits) pipeline assumes a trace-class density operator (Type I or II von Neumann algebra) --- the kind where you can count states and define entropy via $\text{Tr}(\rho \ln \rho)$. Type III$_1$ algebras have no trace, no minimal projections, and no density matrix in this sense. The entropy of a Type III region is not a number --- it is a formal infinity, with only \emph{relative} entropies between states being well-defined (Tomita-Takesaki modular theory).

The framework's current I(bits) pipeline would return a finite number for any input, which is wrong for a Type III system. The correct procedure is to use relative entropy (KL divergence between the state and a reference state) rather than absolute entropy. This requires extending the $\eta_{CP}$ formula to a relative-information form: $\eta_{CP}^{\text{rel}} = D_{KL}(\rho \| \sigma) / (\Delta G_{\text{rel}} / E_{\text{bit}})$. This is a well-defined pipeline extension --- the algebraic structure is unchanged, only the I computation is replaced.

\textbf{Prediction P-26:} \texttt{tensor(BH\_interior, BH\_exterior)} correctly returns $\xi_{\text{ens}} = \max(\xi_{\text{int}}, \xi_{\text{ext}})$ at frac $= 1$ via P-20 --- algebraic factorization failure recovered without new axioms. The pipeline extension to relative entropy (Type III case) is a calibration task, not a structural revision. See §P-26 in PRIMITIVE\_PREDICTIONS.md.

\begin{center}
\rule{0.5\textwidth}{0.4pt}
\end{center}

\subsubsection{XIX.3 ER=EPR and R-Primitive Degeneracy at $G_{\aleph}$}

\textbf{The stress.} Maldacena and Susskind (2013): an Einstein-Rosen bridge (wormhole) connecting two black holes is physically equivalent to maximal quantum entanglement (EPR pair) between them. Entanglement \emph{is} geometry. In the framework, the wormhole and the EPR pair encode differently: the wormhole uses $R_{\Leftrightarrow}$ (mechanical --- topological connection, like a rotaxane thread); the EPR pair uses $R_{\supseteq}$ (non-covalent --- quantum correlation). The ER=EPR conjecture says these are the same physical system described two ways. The framework assigns them different R values. Who is right?

\textbf{The two encodings.} An Einstein-Rosen bridge between two black holes:

\[\text{ER bridge}: \langle D_{\triangle} \;;\; T_{\bowtie} \;;\; R_{\Leftrightarrow} \;;\; P_{\pm}^{\text{sym}} \;;\; F_{\hbar} \;;\; K_{\text{trap}} \;;\; G_{\aleph} \;;\; \Gamma_{\wedge}(\text{QUANTUM}) \;;\; \Phi_{c} \rangle\]

A maximally entangled EPR pair:

\[\text{EPR pair}: \langle D_{\wedge} \;;\; T_{\bowtie} \;;\; R_{\supseteq} \;;\; P_{\pm}^{\text{sym}} \;;\; F_{\hbar} \;;\; K_{\text{trap}} \;;\; G_{\aleph} \;;\; \Gamma_{\wedge}(\text{QUANTUM}) \;;\; \Phi_{c} \rangle\]

\texttt{meet(ER, EPR)} $\to$ CONFLICT on R only (and D: $D_{\triangle}$ vs $D_{\wedge}$). \texttt{d(ER, EPR)} $= w_R \cdot 1 + w_D \cdot 1 = 0.12 + 0.15 = 0.27$ --- small but nonzero, driven by R and D. If ER=EPR is correct, this distance should be zero. The framework gives 0.27. \textbf{The discrepancy is entirely located in R and D.}

The D conflict is interpretable: the ER bridge is a supramolecular/spatial object (macroscopic wormhole geometry, $D_{\triangle}$); the EPR pair is a molecular/quantum object ($D_{\wedge}$). This is a description-level difference, not a physical one --- two observers, one using geometry and one using quantum mechanics, describe the same system with different D assignments. In the $G_{\aleph}$ regime, $D_{\triangle}$ and $D_{\infty}$ are already expected to degenerate at $\Phi_c$ (Axiom 5). The D conflict is a known consequence of Axiom 5 and does not require ER=EPR.

The R conflict is the substantive one. $R_{\Leftrightarrow}$ (mechanical/topological) and $R_{\supseteq}$ (non-covalent/entanglement) are distinguishable in chemical systems: $R_{\Leftrightarrow}$ has a discontinuous steric-cliff barrier profile; $R_{\supseteq}$ follows a smooth Morse potential (see Transformation \#8, §XI). They are orthogonal by construction and grounded independently.

ER=EPR says this orthogonality breaks down at $G_{\aleph}$. More precisely: in a background-independent quantum gravity theory, topological connectivity (wormhole --- a topological handle in spacetime) and quantum correlation (entanglement --- a constraint that propagates non-locally without a geometric channel) are dual descriptions of the same underlying structure. The operational distinction between them --- a physical thread vs a quantum correlation --- dissolves at $G_{\aleph}$ because both are just statements about which degrees of freedom share a common Hilbert space region (Ryu-Takayanagi: entanglement entropy = area of minimal surface = same formula as ER bridge area).

\textbf{R-degeneracy at $G_{\aleph}$ (proposed).} ER=EPR is the claim that $R_{\Leftrightarrow} \equiv R_{\supseteq}$ at $G_{\aleph}$. The R primitive has four values $\{R_{\subseteq}, R_{\supseteq}, R_{\ddagger}, R_{\Leftrightarrow}\}$ that are operationally distinguishable at $G_{\beth}$ --- each has a distinct mechanism, barrier profile, and information signature at the chemical scale. At $G_{\aleph}$, two of these four values become physically indistinguishable. The effective R dimension at $G_{\aleph}$ reduces from 4 to 3: $R_{\Leftrightarrow}$ and $R_{\supseteq}$ merge into a single recognition class --- call it $R_{\text{topo/ent}}$, topological-entanglement recognition.

\textbf{Extension of Axiom 5.} Axiom 5 as stated specifies only G/D degeneracy at $\Phi_c$: $G$ and $D$ become degenerate, the effective tuple contracts by one dimension. ER=EPR predicts a further contraction: at $G_{\aleph}$, R also partially degenerates. The extended statement:

\textbf{Axiom 5 (extended):} At $G_{\aleph}$ and $\Phi_c$: (i) G/D degenerate --- $G_{\aleph}$ and $D_{\infty}$ become informationally equivalent (original Axiom 5); (ii) R-degenerate --- $R_{\Leftrightarrow} \equiv R_{\supseteq}$ (ER=EPR); the effective independent primitive count is 10 $\to$ 8. Further primitive degeneracies (e.g., F and $\xi_{CP}$ conflation at extreme temperature regimes) are possible and would extend the contraction further.

\textbf{Quantitative test.} \texttt{tensor(ER, EPR)} under P-20: frac $= 8/10 = 0.8$ (D and R differ); $\lambda = 0.8$; $\xi_{\text{ens}} = \xi_{ER} + \xi_{EPR} - 0.8 \cdot \min(\xi_{ER}, \xi_{EPR})$. If $\xi_{ER} \approx \xi_{EPR} \equiv \xi$: $\xi_{\text{ens}} \approx 1.2\xi$. Under ER=EPR (R not an independent primitive at $G_{\aleph}$): frac $= 9/10 = 0.9$ (only D differs), $\xi_{\text{ens}} \approx 1.1\xi$. The residual gap is 0.1$\xi$ --- the contribution of D to the frac count. Under the full D/R degeneracy of extended Axiom 5: frac $= 10/10 = 1$, $\xi_{\text{ens}} = \xi$ (they are the same system). \textbf{Prediction P-27: $\xi_{CP}(\text{ER bridge}) = \xi_{CP}(\text{EPR pair})$, with the equality requiring frac $= 1$ under the extended Axiom 5 R-degeneracy.} This is the holographic entanglement-geometry correspondence expressed as a primitive idempotency condition.

\textbf{What the framework says and does not say about ER=EPR.} It does not derive the Ryu-Takayanagi formula or prove that wormholes are dual to entanglement. It says: if ER=EPR holds, the R primitive is not an independent primitive at $G_{\aleph}$ --- the four-value R set contracts to three values at that granularity, and Axiom 5 must be extended to account for this. The framework predicts the form of the degeneracy ($R_{\Leftrightarrow} \equiv R_{\supseteq}$, not any other pair) because these are the two recognition modes whose operational distinction relies on a geometric background (the thread of a rotaxane requires a spatial thread; the steric cliff requires a physical aperture). At $G_{\aleph}$, where the geometric background is itself a consequence of entanglement, the distinction dissolves. $R_{\subseteq}$ (covalent) and $R_{\ddagger}$ (catalytic) remain distinct even at $G_{\aleph}$: they encode chemical bond formation and catalytic transformation, which are operationally distinguishable regardless of the background structure.

\textbf{Synthesis: what the three tests reveal.} The three stress tests probe three different layers of the framework's architecture. Black hole entropy tests the I(bits) pipeline --- a calibration gap (temperature-relative mode needed) with the correct answer already derivable from the G\_aleph primitive. Hilbert space factorization tests the tensor algebra --- no gap; P-20's boundary condition already encodes the correct limit. ER=EPR tests the primitive orthogonality assumption --- a genuine structural revision is indicated: Axiom 5 must be extended to include R-degeneracy at $G_{\aleph}$. In all three cases the framework either already contains the correct result or points precisely at where it must grow. None of the three tests falsify the framework; all three refine it.

\subsection{XIX.5 Recent Advances (v0.4.4, 2026-03-18)}

\subsubsection{XX.1 Reflexive Closure: The Grammar Reads Its Own Axioms}

\emph{Experiment: \texttt{axiom\_reflexive\_tests.py}}

The seven SynthOmnicon axioms were encoded as synthon tuples using the framework's own primitive set, and the full algebra was run over them. Key results:

\begin{itemize}
    \item \textbf{meet(A3, A5) preserves $\Phi$\_c.} A3 (cooperative induction superlinearity, $G_{\gimel} \to G_{\aleph}$) and A5 (recursive tuple embedding, $G_{\aleph}$ + $\Phi_c$) share $\Phi$\_c in their meet --- criticality is invariant under intersection of its own axioms. The framework's most powerful property survives self-reference.
    \item \textbf{Global meet = $\perp$.} The meet of all seven axioms collapses to the conflict sentinel across all primitives. This is correct: the axioms span the primitive space by construction --- they are not redundant constraints, they are independent dimensions of the type system. A grammar whose axioms share a common primitive floor would be over-constrained.
    \item \textbf{Criticality probe orders A5 \textgreater{} A3 \textgreater{} A4 = A6 \textgreater{} A1 = A7 \textgreater{} A2.} Axiom 5 (recursive embedding, $D_{\text{all}}$, $G_{\aleph}$, $\Phi_c$) has the highest $\Phi$\_c candidacy score. Axiom 2 (local ordering without global coordination, $G_{\beth}$) is the most subcritical. The ordering is structurally correct.
    \item \textbf{tensor(A3, A5)} $\rightarrow$ $G_{\aleph}$ / $\Phi_c$ / $\xi_{CP} = 14.39$ nats. The axiom pair whose meet preserves $\Phi$\_c also produces a tensor product at the global granularity level --- the framework detects the axiom-level quantum critical point.
    \item \textbf{A1 $\leftrightarrow$ A7: d = 1.9, 1-hop path.} The two closing-bond axioms (self-complementarity floor for $T_{\bowtie}$, and grounding requirement for assembly direction) are the closest axiom pair, connected by a direct HotSwap.
\end{itemize}

\textbf{Interpretation.} The framework can be applied to its own rules without producing contradiction or trivial results. The reflexive closure is well-defined. The grammar is \emph{not} self-contradictory.

\subsubsection{XX.2 IIT vs. Tensor $\xi_{CP}$: The Partition Presupposition as Structural Conflict}

\emph{Experiment: \texttt{iit\_vs\_tensor\_xi\_tests.py}. Written up in METAPHYSICS.md §IV.1.}

IIT's $\Phi$ and the tensor $\xi_{CP}$ were encoded as synthons and compared algebraically. The result is a formal proof that the two measures are structurally incommensurable:

\begin{table}[htbp]
\centering
\small
\rowcolors{1}{white}{white}
\begin{tabularx}{\textwidth}{>{\centering\arraybackslash}X >{\centering\arraybackslash}X >{\centering\arraybackslash}X}
\toprule
\rowcolor{gray!25}\textbf{Pair} & \textbf{d} & \textbf{Conflict set} \\
\midrule
\rowcolors{1}{white}{gray!10}
IIT\_Phi $\leftrightarrow$ Tensor\_xi\_CP & 8.1 & $\{D, T, R, P, \Gamma\}$ --- five primitive conflicts \\
IIT\_Phi $\leftrightarrow$ Edelman\_Degeneracy & 3.9 & $\{D, G\}$ \\
Tensor\_xi\_CP $\leftrightarrow$ Meet\_Richness & 1.5 & $\{T\}$ only \\
Tensor\_xi\_CP $\leftrightarrow$ Edelman\_Degeneracy & 5.5 & $\{D, T, G\}$ \\
\bottomrule
\end{tabularx}
\end{table}

No HotSwap path exists between IIT\_Phi and Tensor\_xi\_CP --- the $D$/$T$ mismatch blocks all paths. The two measures are in different universality classes. The IIT partition presupposition ($G_{\beth}$, $T_{\text{linear}}$, $R_{\text{mechanical}}$, $P_{\text{directional}}$, external MIP decomposition) encodes five independent primitive differences from the intrinsic non-partition measure ($G_{\aleph}$, $T_{\text{network}}$, $R_{\ddagger}$, $P_{\pm}^{\psi}$, meet-richness without partition). The partition is not a computational convenience; it is a structural presupposition embedded in five orthogonal primitives simultaneously.

IIT $\otimes$ Edelman $\to$ $G_{\gimel}$/$\Phi_{\text{sub}}$: the tensor of IIT with Edelman degeneracy stays trapped at mesoscale/subcritical. Meet\_Richness $\otimes$ Edelman $\to$ $G_{\aleph}$/$\Phi_c$: replacing the partition with meet-richness unlocks global scope and criticality. The algebra makes the argument formally.

\subsubsection{XX.3 Protein--Condensate Equivalence and K-Targeting (P-48)}

\emph{Formal prediction in PRIMITIVE\_PREDICTIONS.md.}

The condensate gel and amyloid fibril share identical primitive encodings:


\[\langle D_{\bigtriangleup}; T_{\in}; R_{\supseteq}; P_{\pm}^{\psi}; F_{\hbar}; K_{\text{trap}}; G_{\gimel}; \Gamma_{\wedge}(\text{SELECTIVE}); \Phi_{\text{sub}} \rangle\]

Tuple distance $d(\text{condensate\_gel}, \text{amyloid}) = 0.00$. This is not a coincidence: both are supramolecular network assemblies with non-covalent recognition, pseudosymmetric polarity, high fidelity, and kinetically trapped topology. They are the same primitive event --- a P-34-class result now extended to a cross-disease therapeutic principle.

\textbf{Jacobian result:} $K$-targeting (K\_trap $\rightarrow$ K\_fast via disaggregase) reduces distance by $>1.5\times$ more than $F$-targeting (F\_hbar $\rightarrow$ F\_eth via competing binder) for both systems. \textbf{Prediction P-48:} K-targeting is the mechanistically preferred therapeutic strategy for both condensate gel pathology and amyloid fibrillization, with the superiority quantifiable by a Jacobian experiment. See P-34, P-35, P-48 in PRIMITIVE\_PREDICTIONS.md.

\subsubsection{XX.4 D\_holo: Holographic Dimensionality as a First-Class Primitive}

\emph{Implemented in \texttt{synthomnicon/models.py} $\cdot$ \texttt{ads\_cft\_boundary} synthon registered.}

The AdS/CFT boundary encoding previously required a hybrid $D = \{D_{\bigtriangleup}, D_{\infty}\}$ proxy (see §XVII Result 8). The gap was that the bulk-boundary correspondence is not a spatial+temporal operation; it is a dimensional reduction in which $d$-dimensional bulk information is encoded on a $(d-1)$-dimensional boundary. This is qualitatively different from any combination of the existing dimensionality values.

$D_{\text{holo}}$ is the twelfth dimensionality value: bulk degrees of freedom encoded on a holographic boundary screen. The canonical synthon:

\[\text{ads\_cft\_boundary}: \langle D_{\text{holo}}; T_{\in}; R_{\ddagger}; P_{\pm}^{\psi}; F_{\eth}; K_{\text{mod}}; G_{\aleph}; \Gamma_{\wedge}(\text{SELECTIVE}); \Phi_c \rangle\]

\textbf{Key algebraic result:} \texttt{transition(ads\_cft\_boundary, topological\_insulator)} $\rightarrow$ 1st-order (D: D\_holo $\neq$ D\_triangle), infinite cost, asymmetry = 1.0. The holographic boundary is not continuously deformable into any bulk phase --- the bulk-boundary map is a virtual Kleisli arrow. This matches the holographic duality literature: the correspondence is a dual \emph{description}, not a continuous deformation.

\subsubsection{XX.5 Phase Transitions as Morphisms}

\emph{Implemented in \texttt{synthomnicon/morphism.py} $\cdot$ \texttt{syncon transition SRC DST} CLI.}

Phase transitions are encoded as Kleisli arrows in the HotSwap monad. \texttt{find\_transition(src, dst, catalog)} returns a \texttt{TransitionMorphism} dataclass with:

\begin{itemize}
    \item \textbf{Order classification:} 2nd order (HotSwap path exists through $\Phi$\_c intermediates) or 1st order (no path --- structural D/T conflict or F-floor)
    \item \textbf{Forward/reverse costs:} total $\Delta\xi_{CP}$ on each path ($\infty$ if no path)
    \item \textbf{Asymmetry:} $|\text{fwd} - \text{rev}| / \max(\text{fwd}, 1)$ --- the irreversibility signature
    \item \textbf{$\Phi$\_c intermediates:} names of critical-phase synthons on the forward path
\end{itemize}

\textbf{Key result --- topological protection as morphism irreversibility:}

\begin{lstlisting}
syncon transition topological_insulator_bi2se3 synthon_Fermi_liquid
  Order: 1st-order (discontinuous)
  Forward cost: inf    Reverse cost: 0.288 nat
  Asymmetry: 1.000 (irreversible)
\end{lstlisting}

The TI $\rightarrow$ Fermi liquid transition is blocked forward (F\_hbar $\rightarrow$ F\_eth is a fidelity downgrade) but permitted in reverse (Fermi liquid $\rightarrow$ TI is a fidelity upgrade). \textbf{Topological protection encodes as morphism irreversibility.} The asymmetry = 1.0 is not a numerical coincidence; it reflects the categorical fact that a topologically protected phase cannot be continuously deformed into an unprotected one.

The new \texttt{syncon transition} command makes transition morphisms a first-class CLI operation, alongside \texttt{syncon path} (HotSwap path), \texttt{syncon meet}, \texttt{syncon tensor}.

\subsubsection{XX.6 QCP Morphism: GR $\rightarrow$ Asymptotic Safety $\rightarrow$ SM (First 2nd-Order Demonstration)}

\emph{Implemented in \texttt{synthomnicon/morphism.py} $\cdot$ \texttt{syncon transition} CLI $\cdot$ catalog: \texttt{general\_relativity}, \texttt{asymptotic\_safety\_reuter\_fp}, \texttt{standard\_model}.}

First successful quantum critical point morphism demonstration. Three gravity-cluster synthons share $D_{\bigtriangleup} / T_{\in}$, enabling a 2-segment $\Phi_c$-mediated path:

\begin{table}[htbp]
\centering
\small
\rowcolors{1}{white}{white}
\begin{tabularx}{\textwidth}{>{\centering\arraybackslash}X >{\centering\arraybackslash}X >{\centering\arraybackslash}X}
\toprule
\rowcolor{gray!25}\textbf{Synthon} & \textbf{Key primitives} & \textbf{Role} \\
\midrule
\rowcolors{1}{white}{gray!10}
\texttt{general\_relativity} & $K_{\text{slow}}$, $\Phi_{\text{sub}}$ & Source \\
\texttt{asymptotic\_safety\_reuter\_fp} & $K_{\text{trap}}$, \textbf{$\Phi_c$} & QCP intermediary \\
\texttt{standard\_model} & $K_{\text{fast}}$, $\Phi_{\text{sub}}$ & Destination \\
\bottomrule
\end{tabularx}
\end{table}

\begin{lstlisting}
syncon transition general_relativity standard_model
  Order:         2nd-order (continuous)
  Path:          general_relativity -> asymptotic_safety_reuter_fp -> standard_model
  Forward cost:  1.153 nat   Reverse cost: 1.153 nat
  Asymmetry:     0.000 (reversible)
    Quantum critical point detected
  QCP synthon:   asymptotic_safety_reuter_fp
\end{lstlisting}

The framework produces this from primitive structure alone: $K_{\text{slow}}$ (GR, spacetime curvature is kinetically slow), $K_{\text{trap}}$ (Reuter fixed point, trapped near $\Phi_c$), $K_{\text{fast}}$ (SM, gauge interactions are kinetically fast). The QCP bridge works because the Reuter fixed point carries $\Phi_c$ and shares the same $D/T$ cluster as both endpoints.

\textbf{Asymmetry = 0.000:} The GR $\leftrightarrow$ SM transition is bidirectional with equal cost. Effective field theory is a downhill projection in the kinematic lattice ($K_{\text{slow}} \to K_{\text{fast}}$), but the thermodynamic cost of the path is the same in both directions --- the path cost is set by the barrier at the QCP, not by an irreversibility of the kinematic ordinal change.

\textbf{§XVII complement:} The SM/QG barrier (lift blocked at $G_{\beth}$, 4 CONFLICTS) is now interpretable as a 1st-order morphism result: QG's $D_{\infty}/G_{\aleph}$ conflict with SM's $D_{\bigtriangleup}/G_{\beth}$ produces a virtual Kleisli arrow. Asymptotic Safety does not bridge SM $\leftrightarrow$ QG --- it lives in $G_{\beth}$ and bridges GR $\leftrightarrow$ SM within the effective field theory regime. The two results are complementary: there exists a 2nd-order path connecting the two successful quantum field theories (via AS Reuter FP), and no path connecting either to full quantum gravity.

\subsubsection{XX.7 Floquet Synthons and Discrete Time Crystals}

\emph{Catalog entries: \texttt{floquet\_chern\_insulator}, \texttt{time\_crystal\_dtc}.}

Two synthons encoding periodically driven (Floquet) quantum matter:

\[\text{floquet\_chern\_insulator}: \langle D_{\bigtriangleup}; T_{\uparrow\downarrow}; R_{\ddagger}; P_{\pm}^{\psi}; F_{\eth}; K_{\text{trap}}; G_{\gimel}; \Gamma_{\to}(\text{SELECTIVE}); \Phi_{\text{sub}}; \Omega_C \rangle\]

\[\text{time\_crystal\_dtc}: \langle D_{\infty}; T_{\bowtie}; R_{\supseteq}; P_{\pm}^{\text{sym}}; F_{\eth}; K_{\text{MBL}}; G_{\gimel}; \Gamma_{\to}(\text{SELECTIVE}); \Phi_{\text{sub}}; \Omega_Z \rangle\]

Key primitive assignments: \textbf{Floquet Chern insulator} --- $\Gamma_{\to}(\text{SEQUENTIAL})$ encodes the stroboscopic Floquet time-evolution operator (periodic drive is sequential by construction); $K_{\text{trap}}$ encodes the driven-phase kinetic arrest; $\Omega_C$ (Chern number) encodes the Floquet topological invariant. \textbf{Discrete time crystal} --- $D_{\infty}$ (temporal) encodes period-doubling; $K_{\text{MBL}}$ encodes the many-body localization required to stabilize DTC order against Floquet heating (without MBL, the drive heats the system to infinite temperature); $\Omega_Z$ (winding number) encodes $\mathbb{Z}$-class protection of the subharmonic response.

\textbf{Structural result:} $d(\text{floquet\_chern\_insulator}, \text{time\_crystal\_dtc}) = 4.2$. Shared: $G_{\gimel}$, $F_{\eth}$, $\Phi_{\text{sub}}$. Conflicting: $D$ (spatial vs. temporal), $T$ (braid vs. cyclic), $K$ (trap vs. MBL). The structural distinction maps the physics correctly: both are periodically driven but inhabit different topological classes ($\Omega_C$ vs. $\Omega_Z$) and dimensionality regimes.

\subsubsection{XX.8 $Z_2$ Gauge Theory / Toric Code: Zero-Cost Topological Adjacency}

\emph{Catalog entry: \texttt{z2\_lattice\_gauge\_toric\_code}.}

\[\text{z2\_lattice\_gauge\_toric\_code}: \langle D_{\bigtriangleup}; T_{\uparrow\downarrow}; R_{\Leftrightarrow}; P_{\pm}^{\text{sym}}; F_{\hbar}; K_{\text{trap}}; G_{\aleph}; \Gamma_{\wedge}(\text{QUANTUM}); \Phi_{\text{sub}}; \Omega_{Z_2} \rangle\]

Primitive assignments: $T_{\uparrow\downarrow}$ (Wilson loop holonomy = anyonic braiding phase), $R_{\Leftrightarrow}$ (loop variable recognition --- the relevant degree of freedom is topological, not site-local), $G_{\aleph}$ (toric code stabilisers are global: vertex and plaquette operators span the full lattice), $\Gamma_{\wedge}(\text{QUANTUM})$ (stabiliser operations are superposition-preserving), $\Omega_{Z_2}$ (ground-state degeneracy depends on the topology of the manifold --- the defining feature of $\mathbb{Z}_2$ topological order).

\textbf{Key result:}

\[d(\text{z2\_toric}, \text{fqh\_moore\_read}) = 3.90 \qquad \text{path cost} = 0.000 \text{ nat}\]

Zero path cost at non-zero tuple distance. The two systems share the same thermodynamic constraint profile (identical $F$ and $\xi_{CP}$ per hop) despite differing in $R$ and $\Omega$ ($\Omega_{Z_2}$ vs. $\Omega_{NA}$). This is the algebraic signature of two topological orders in the same thermodynamic universality class: categorically distinct (d = 3.90) but connected by a zero-cost continuous path. The framework encodes the known physics: toric code and Moore-Read FQH are both Abelian-adjacent to the same non-Abelian topological sector, and both exhibit ground-state degeneracy on non-trivial manifolds.

\subsubsection{XX.11 METAPHYSICS.md §XV: Synthonicon as $\Phi_c$ Event in Knowledge-Space (2026-03-19)}

\textbf{METAPHYSICS.md §XV --- The Synthonicon as $\Phi_c$ Event.} The framework's own emergence traced as a phase transition in scientific knowledge-space. Five subsections:

\begin{itemize}
    \item §XV.1: The recursion closes --- completeness criterion forced $G_{\beth} \to G_{\gimel} \to G_{\aleph}$ traversal; this was algebraically necessary, not a design choice.
    \item §XV.2: $\text{Grammar} \otimes \text{Corpus} \to \text{Synthonicon}$ --- the framework is a tensor product of relational structure and empirical content; neither alone is sufficient.
    \item §XV.3: The F-floor ratchet in knowledge-space --- cross-domain fidelity barrier dissolved by a common primitive metric ($\xi_{CP}$); the floor has moved and cannot move back.
    \item §XV.4: The LLM-human tensor product as transducer --- $\text{Human} \otimes \text{LLM} \to \Phi_c$; LLM provides $G_{\aleph}$ retrieval bandwidth ($K_{\text{fast}}$); human provides axiom enforcement ($F_{\hbar}$, $K_{\text{trap}}$); neither reaches criticality alone.
    \item §XV.5: The tool describes its own origin --- reflexive closure at Level 4; the relational ontology is validated by the conditions required to discover it.
    \item §XV.6: Predictions confirmed in real-time (2025--2026): grokking as $\Phi_c$ event (singular learning theory); condensate-amyloid $d = 0.00$ (Nature Comm Chem Nov 2025); $K$-targeting therapeutics (Hsp70 activators); AI-accelerated science as current operational mode. Timescale assessment: the phase transition was not predicted for August 2026 --- it is already in progress.
\end{itemize}

\begin{center}
\rule{0.5\textwidth}{0.4pt}
\end{center}

\subsubsection{XX.10 Catalog Merge $\cdot$ Design Agent Fix $\cdot$ METAPHYSICS.md §XI--XIV $\cdot$ P-ARCH-1--4 (2026-03-18)}

\textbf{Catalog merge + dedup.} 1,638 discovery-generated synthons (previously orphaned in \texttt{discovery\_output/} JSON snapshots) merged with 54 hand-curated entries. After deduplication by unique notation: 223 canonical synthons (1,469 duplicates culled; 2 leaked-prompt names removed). Backup at \texttt{\textasciitilde{}/.synthomnicon/catalog.json.bak\_20260318\_100146}.

\textbf{Design agent naming fix.} \texttt{syncon design} was writing the LLM's internal description string as the catalog name. Three-layer fix: \texttt{\_slugify\_goal()} extracts a clean name from the first 5 words of the goal; \texttt{\_dispatch\_tool()} injects it via \texttt{kwargs.setdefault(''name'', self.\_goal\_slug)}; \texttt{explicit\_name} parameter added to \texttt{\_create\_synthon\_from\_data()} in both \texttt{axiom\_guided\_generator.py} and \texttt{synthon\_generator\_agent.py}, bypassing the LLM JSON name field entirely. Conditional Varma QXY routing added: physical/chemical goals $\rightarrow$ Varma QXY criticality probe; non-physical goals (linguistic, cognitive, ecological) $\rightarrow$ $G_{\aleph}$-based criticality with domain-appropriate $T$/$R$/$K$/$\Gamma$.

\textbf{METAPHYSICS.md §XI --- The Convergence Argument.} Three independent routes (Kastrup idealism, Glattfelder ICR, SynthOmnicon constraint algebra) converge on structurally equivalent claims. Self-similarity argument: the grammar applies to itself --- not what a domain-specific model looks like. Tautology-from-inside framing: if literally true, ''all interacting systems follow these rules'' is analytic but unrecognizable as such from inside.

\textbf{METAPHYSICS.md §XII --- The Sage and the Algebra.} Tao Te Ching §20 (Mitchell tr.) as the inside view of what the algebra describes from outside. ''I alone possess nothing'' = no-intrinsic-properties result from inside. Opposite entry points (stopping thought vs. extreme precision about molecular recognition) arriving at the same structural description.

\textbf{METAPHYSICS.md §XIII --- External Framework Assessments.}

\begin{itemize}
    \item §XIII.1 UCBFT / Lattice Flush: 143.48T = human synaptic proxy (86T $\times$ 5/3 cooperativity factor), not a vacuum constant. ''Snap'' = $\Phi_c$ (criticality threshold). ''Flush'' = $\xi_{CP} \to 0$. $K_{\text{trap}}$ = hallucination spike mechanism (pathway multiplicity without fidelity floor).
    \item §XIII.2 CCT (Quantum Gravity Research Cycle Clock Theory): PEL $\leftrightarrow$ $\xi_{CP}$ minimization (isomorphic map). E8/phasons encodable via $D_{\text{holo}}$ + $T_{\uparrow\downarrow}$ + $\Omega_{NA}$ but not required. Retrocausality requires backward path algebra (non-trivial extension). Constants derivation epistemically inert if ordinal sufficiency generalizes.
\end{itemize}

\textbf{METAPHYSICS.md §XIV --- The Architecture Prediction.} $\Phi_c$ threshold is $D$-primitive-specific: $D_{\text{holo}}$ at \textasciitilde{}10T ($G_{\aleph}$ built-in), $D_{\infty}$ (SSM) at \textasciitilde{}143T/ln(N), $D_{\wedge\bigtriangleup}$ (Transformer) at \textasciitilde{}143T. ''Singularity'' is a point in architecture space, not a date on the scaling curve. Formal entries: P-ARCH-1--4 in PRIMITIVE\_PREDICTIONS.md.

\textbf{Epistemological note.} External frameworks (UCD/Richie Wise, UCBFT, CCT) that correctly identify the topology of AI scaling (wall, phase transition structure) while confabulating the mechanism exhibit a consistent structural signature: $\Gamma_{\vee}$ evidence grammar (any corroborating signal accepted, no falsifiability gate) fills the inaccessible interior (§X.4) with available surface-feature patterns. The Synthonicon epistemological discipline --- theory $\rightarrow$ low-level tests $\rightarrow$ rigorous tests $\rightarrow$ stricter axioms, with increasing predictive power at each step --- is the operationally distinguishable contrast.

\begin{center}
\rule{0.5\textwidth}{0.4pt}
\end{center}

\subsubsection{XX.9 Phase 3e Stress Tests Complete $\cdot$ Catalog: 51 Synthons}

Phase 3e stress test agenda (§3e.5) is fully resolved:

\begin{table}[htbp]
\centering
\small
\rowcolors{1}{white}{white}
\begin{tabularx}{\textwidth}{>{\centering\arraybackslash}X >{\centering\arraybackslash}X >{\centering\arraybackslash}X}
\toprule
\rowcolor{gray!25}\textbf{Stress test} & \textbf{Status} & \textbf{Key result} \\
\midrule
\rowcolors{1}{white}{gray!10}
Phase transitions as morphisms &  complete & \texttt{morphism.py}; asymmetric 1st order TI$\rightarrow$FL; asymmetry=1.000 \\
Floquet synthons (periodic drive) &  complete & \texttt{floquet\_chern\_insulator} + \texttt{time\_crystal\_dtc}; d=4.2 \\
Gauge theories (loop variables) &  complete & \texttt{z2\_lattice\_gauge\_toric\_code}; d=3.90 to Moore-Read; 0.000 nat \\
QCP morphism &  complete & GR$\rightarrow$Reuter FP$\rightarrow$SM; 2nd order; asymmetry=0.000 \\
SM/QG unification gap &  complete & 1st order (virtual); 4 CONFLICTS; §XVII \\
\bottomrule
\end{tabularx}
\end{table}

Catalog size after 2026-03-18 merge + dedup: \textbf{223 unique synthons} spanning molecular/supramolecular, temporal/autocatalytic, quantum/topological, physics theories, protein/condensate, and programmable matter domains. (Pre-dedup: 1,692 entries from 54 hand-curated + 1,638 discovery outputs; 1,469 duplicates culled by notation uniqueness.)

\begin{center}
\rule{0.5\textwidth}{0.4pt}
\end{center}

\subsection{XXI. Decomposition Algebra: Inverse Operations and the Six New Predictions (2026-03-19)}

\subsubsection{XXI.1 The Decomposition Module}

A complete \emph{inverse} algebra for the build-up operations (tensor, meet, join, lift, path) was implemented in \texttt{synthomnicon/decompose.py} and registered in \texttt{\_\_init\_\_.py}. The eight operations are:

\begin{table}[htbp]
\centering
\small
\rowcolors{1}{white}{white}
\begin{tabularx}{\textwidth}{>{\centering\arraybackslash}X >{\centering\arraybackslash}X >{\centering\arraybackslash}X}
\toprule
\rowcolor{gray!25}\textbf{Operation} & \textbf{Symbol} & \textbf{Meaning} \\
\midrule
\rowcolors{1}{white}{gray!10}
\texttt{project(s, primitives)} & $\pi_P(s)$ & Retain only named primitive slots; zero the rest \\
\texttt{primitive\_peel(s, prim)} & $s \setminus p$ & Descend one tier on primitive $p$; returns peel cost \\
\texttt{factor(s)} & $\text{fac}(s)$ & Find the strongest meet-irreducible factor \\
\texttt{principal\_decomp(s)} & $\text{pd}(s)$ & Recursive factor chain to join-irreducible atoms \\
\texttt{cofactor(C, A)} & $C \mathbin{/} A$ & Residual $B$ such that $A \otimes B \approx C$ \\
\texttt{complement\_rel(s, ref)} & $\overline{s}_\text{ref}$ & Primitives where $s$ is strictly below the reference \\
\texttt{kernel(s, probe)} & $\ke$r_{\phi}$(s)$ & Strip all primitives that do not fire the probe \\
\texttt{retrosynthetic\_path(target, catalog)} & $\text{retro}(t, \mathcal{C})$ & Ranked candidate pairs from catalog \\
\bottomrule
\end{tabularx}
\end{table}

Monadic wrappers (\texttt{project\_m}, \texttt{peel\_m}, \texttt{factor\_m}, \texttt{cofactor\_m}) return structured tuples with notes, warnings, block status, and block reason, composable with \texttt{DesignPipeline}. The \texttt{DesignPipeline} class in \texttt{algebra.py} was extended with \texttt{.project()}, \texttt{.peel()}, \texttt{.factor()}, and \texttt{.cofactor()} methods.

\subsubsection{XXI.2 Key Results from the Decomposition Explorations}

\textbf{Phi\_c categorical independence.} The \texttt{kernel} operation --- stripping all primitives that do not fire a given probe --- reveals that $\Phi_c$ survives the descent of $F$, $K$, and $G$ to their ordinal floors. The \texttt{phi\_c\_skeleton} synthon ($F$=LOW, $K$=FAST, $G$=LOCAL, Phi=CRITICAL) is a valid synthon. Catalog proof: \texttt{asymptotic\_safety\_reuter\_fp} carries $G$=LOCAL + Phi=CRITICAL --- the UV fixed point of asymptotic safety is \emph{locally} critical without global organisation. $\Phi_c$ is a \textbf{categorical phase label}, not a derived ordinal. \emph{Formal entry: P-49 in PRIMITIVE\_PREDICTIONS.md.}

\textbf{Quantization residual as a synthon.} \texttt{cofactor(quantum\_gravity, general\_relativity)} gives the residual $B = \{K_\text{trap}, $G_{\aleph}$, T_\text{braid}, \Phi_c, \Omega_\text{NA}\}$ plus a $D$-CONFLICT (GR=SUPRAMOLECULAR, QG=TEMPORAL). The $D$-CONFLICT is the algebraic form of the \emph{background-dependence problem}: GR organises matter in space; QG must organise spacetime itself. Any quantisation scheme that preserves GR's spatial $D$-component will produce this conflict. The residual $B$ is the complete primitive content of the ''quantisation operator'' encoded as a single synthon. \emph{Formal entry: P-51 in PRIMITIVE\_PREDICTIONS.md.}

\textbf{AdS/CFT holography as a cofactor.} \texttt{cofactor(ads\_cft\_boundary, general\_relativity)} adds exactly $G$=GLOBAL and Phi=CRITICAL above GR. The cofactor $B = \langle $G_{\aleph}$ \otimes \Phi_c \rangle$. Interpretation: $G$: LOCAL $\rightarrow$ GLOBAL encodes the bulk-to-boundary global correlation structure (the holographic RG); Phi: SUBCRITICAL $\rightarrow$ CRITICAL encodes the boundary CFT's critical point. Holography is the operation that adds global correlation and criticality to a local classical gravity theory --- encoded in two primitives. \emph{Formal entry: P-54 in PRIMITIVE\_PREDICTIONS.md.}

\textbf{SM and QG principal decompositions.} \texttt{principal\_decomp(standard\_model)} yields 3 atoms; \texttt{principal\_decomp(quantum\_gravity)} yields 8 atoms; 3 are shared --- the \emph{unification substrate}. The direct \texttt{cofactor(QG, SM)} is blocked by a $D$-CONFLICT: SM $D$=SUPRAMOLECULAR (spatial matter organisation), QG $D$=TEMPORAL (spacetime organisation). This confirms P-23: the background-dependence problem is the primary barrier to unification. The 3 shared atoms represent the currently accessible common ground.

\textbf{Amyloid nucleation F-CONFLICT.} \texttt{cofactor(amyloid\_fibril, condensate\_liquid)} $\rightarrow$ F-CONFLICT: tensor-min(MEDIUM, $B$) $\leq$ MEDIUM $<$ HIGH. A liquid condensate cannot template amyloid by primitive tensor alone --- an external $F$=HIGH nucleation seed is \textbf{algebraically required}. \emph{Formal entry: P-50 in PRIMITIVE\_PREDICTIONS.md.}

\textbf{GNF-2 combination strategy.} Drug panel cofactor analysis: GNF-2 is the only drug carrying $\Phi_c$ (distance 3.5, zero conflicts). Cofactor residual = \{$D$=SUPRAMOLECULAR, $T$=NETWORK\} --- GNF-2 needs a supramolecular network scaffold partner (PROTAC linker, DNA nanostructure, or polyvalent hydrogel) to close the allosteric gap. \emph{Formal entry: P-52 in PRIMITIVE\_PREDICTIONS.md.}

\textbf{Three stability regimes.} Peel cost profiles distinguish thermodynamic stability (DNA origami: zero costs, gradual melting), phase-protected stability (LLPS condensate: Phi peel cost = 3.0 nats, sharp 2-state), and kinetically frozen stability (condensate gel / amyloid: $K$=TRAP at floor, zero costs but mechanically immobile). LLPS condensates should show sharper melting cooperativity than DNA origami at matched $T_m$. \emph{Formal entry: P-53 in PRIMITIVE\_PREDICTIONS.md.}

\subsubsection{XXI.3 The GR $\rightarrow$ QG Bridge Network}

Retrosynthetic analysis of QG via the full catalog reveals:

\begin{itemize}
    \item $d(\text{GR}, \text{QG}) = 6.4$ --- the classical/quantum gap
    \item $d(\text{GR}, \text{asymptotic\_safety}) = 1.8$ --- AS is a GR-like fixed point
    \item $d(\text{QG}, \text{asymptotic\_safety}) = 6.6$ --- AS is not full QG
    \item $d(\text{GR}, \text{AdS/CFT}) = 6.0$ --- AdS/CFT is not a stepping stone through GR
    \item Best single-synthon approximation to QG in catalog: \texttt{fqh\_moore\_read} at $d = 3.1$
\end{itemize}

The path GR $\rightarrow$ AS $\rightarrow$ AdS/CFT $\rightarrow$ QG has total stepped distance 15.2 vs. direct 6.4. Asymptotic safety and AdS/CFT are \emph{detours} in primitive space, not stepping stones to QG --- consistent with the EFT hierarchy being a sequence of effective theories at different scales, not a path toward quantisation.

\emph{Six predictions derived: P-49 through P-54. Exploration scripts: \texttt{decompose\_explorations.py} and \texttt{decompose\_explorations2.py}.}

\begin{center}
\rule{0.5\textwidth}{0.4pt}
\end{center}

\subsection{XXII. On the Nature of the Framework: Universal Conditional Logic and the Algorithmic Assignment Project (2026-03-20)}

\subsubsection{XXII.1 SynthOmnicon Is a Universal Conditional Logic, Not a Theory of Everything}

A Theory of Everything makes ontological claims --- it asserts what reality is at bottom and derives the specific laws from that substrate. SynthOmnicon makes no such claim. What it asserts is more precise and more limited: \textbf{given any system with internal structure, certain conditional relationships hold}. The primitives do not explain \emph{why} general relativity exists; they identify \emph{what GR is conditional on}, what it can and cannot produce via tensor, and what must be added to reach QG. That is the logic of conditionality, not an ontological foundation.

The universal conditional logic (UCL) claim has a formal content: the same conditional structure --- $K_\text{trap}$ propagation, the $F$-floor ratchet, $\Phi_c$ join-dominance --- appears identically in molecular chemistry, protein biology, condensed matter, and quantum gravity. Not because these domains share a physical substrate, but because they share a \textbf{constraint grammar}. The conditional logic is invariant under substrate change. In the same way that Boolean algebra is not a theory of everything about circuits but is the universal conditional logic of two-valued systems, SynthOmnicon claims to be the Boolean algebra of self-organising systems.

The key epistemological consequence: \emph{a wrong primitive encoding falsifies the encoding, not the logic}. If the prediction fails, the primitive assignment is wrong --- not the algebra, not the axioms, not the tautological inference structure. This is the precise content of ''UCL'' rather than ''TOE.''

\subsubsection{XXII.2 The Natural-Joints Question and the Algorithmic Assignment Project}

The UCL claim is established by the formal results. A deeper question --- whether the primitives are \emph{natural joints} rather than useful conventional discretisations --- is empirically open but non-trivially supported.

The evidence is specific: the $+2.303$ nat cost ($= \ln 10$) appearing identically across topological matter, the $d = 0.00$ identity between A$\beta$ and $\alpha$-synuclein fibrils, and the SM/QG four-conflict result are not merely cross-domain successes --- they are \emph{numerically conserved} values across physically unrelated systems. An arbitrary coarse-graining would not produce this. The fact that the $K_\text{trap} \to K_\text{MBL}$ cost is exactly $\ln 10$ --- the natural log of the tier accessibility ratio --- suggests the tier boundary is tracking a real order-of-magnitude scale separation in nature, not a human-chosen bin edge (P-21: F-tier values = integer Boltzmann discrimination ratios). This updates the probability that the primitives are natural joints without proving it.

The \textbf{algorithmic assignment project} is the decisive test. Its core idea: if the primitives are natural joints, then \emph{independent assignment methods} should converge to the same tier values. Assign $F$ from direct $\Delta G$ measurement, from DFT, and from NMR --- if all three land in the same tier, the boundary is real. If they disagree systematically at edges, the boundary is a convenience. This is the test of \textbf{assignment method independence}.

A second test is now enabled by the decomposition algebra: \textbf{self-consistency under decomposition}. Take a system, assign its primitives, run \texttt{principal\_decomp} to get atoms, re-assign the atoms independently from first principles, then verify the atoms compose back to the original encoding. Natural joints should be closed under the algebra's operations.

The algorithmic assignment project will either confirm the natural-joints hypothesis (by consistently hitting the same thresholds across systems and assignment methods) or reveal that the boundaries are fuzzier than the cross-domain results suggest. Either outcome is scientifically productive: confirmation makes SynthOmnicon a \emph{discovery} (like Mendeleev's periodic table tracking real quantum states); fuzziness reveals exactly where the framework's discretisation is conventional and where it needs refinement. The UCL claim is stable under both outcomes --- only the \emph{depth} of the naturalness changes.

\emph{Implementation: \texttt{synthomnicon/assignment.py}. Consistency report against the live catalog: \texttt{assignment\_tests.py}.}

\begin{center}
\rule{0.5\textwidth}{0.4pt}
\end{center}

\subsection{XXIII. Algebraic Operations: Category-Theoretic Translations and Canonical Examples (2026-03-20)}

\subsubsection{XXIII.1 The Seven Operations}

\begin{table}[htbp]
\centering
\small
\rowcolors{1}{white}{white}
\begin{tabularx}{\textwidth}{>{\centering\arraybackslash}X >{\centering\arraybackslash}X >{\centering\arraybackslash}X >{\centering\arraybackslash}X}
\toprule
\rowcolor{gray!25}\textbf{Operation} & \textbf{Category Theory} & \textbf{Lattice / Algebra} & \textbf{SynthonM Stack} \\
\midrule
\rowcolors{1}{white}{gray!10}
\textbf{meet} ⊓ & Product / GLB & Heyting algebra meet & $\Phi$\_c is co-Heyting absorbing element: $\Phi$\_c⊓x = $\Phi$\_c \\
\textbf{join} ⊔ & Coproduct / LUB & F-floor ratchet & WriterT accumulates $\Delta$\xi\_CP; floor never regresses \\
\textbf{tensor} \otimes & Bifunctor on Synth$\times$Synth$\rightarrow$Synth & \xi\_ens = \xi₁+\xi₂$-$$\lambda$$\cdot$I(s₁;s₂) & T promotes (cage\textgreater{}network\textgreater{}hub\textgreater{}bowtie\textgreater{}linear); $\Phi$\_c and $\Omega$ join-dominant \\
\textbf{lift} & Natural transformation \eta: F$\rightarrow$G & $\Delta$nat = +2.303 nats (= ln 10, Landauer analog) & Criticality lift blocked if F \textless{} F\_\hbar (StateT gate) \\
\textbf{path} & Geodesic in Kleisli--Lawvere metric & BFS on valid-swap graph & Blocked path = 1st-order topological transition \\
\textbf{pipeline / SynthonM} & Kleisli do-notation & WriterT{[}ℝ$\geq$0{]}(StateT{[}Ctx{]}(MaybeT Id)) A & fail-fast on BLOCKED; MaybeT propagates \\
\textbf{decomp} & Inverse bifunctor / Birkhoff & cofactor, kernel, project, complement\_rel & principal\_decomp $\rightarrow$ join-irreducible atoms \\
\bottomrule
\end{tabularx}
\end{table}

\subsubsection{XXIII.2 Canonical Cross-Domain Example Table}

The Hv1 trilogy maps every operation onto concrete biology:

\begin{table}[htbp]
\centering
\small
\rowcolors{1}{white}{white}
\begin{tabularx}{\textwidth}{>{\centering\arraybackslash}X >{\centering\arraybackslash}X >{\centering\arraybackslash}X}
\toprule
\rowcolor{gray!25}\textbf{Operation} & \textbf{Canonical Hv1 Example} & \textbf{Result} \\
\midrule
\rowcolors{1}{white}{gray!10}
meet & meet(Hv1\_human\_open, AtHv1\_primed) & $\Phi$\_c dominates; T\_bowtie preserved \\
join & join(above, PsHv1\_constitutive) & F$\rightarrow$F\_eth; $\Phi$\_c join-dominant; $\Delta$\xi=0 \\
tensor & tensor(Hv1\_human\_open, 2GBI\_inhibitor) & T-CONFLICT + P-CONFLICT; $\Phi$\_c survives \\
lift & lift(AtHv1\_silent, critical) & BLOCKED --- K\_trap gate prevents lift until peel(K\_trap) \\
path & path(AtHv1\_silent $\rightarrow$ Hv1\_human\_open) & 3-primitive hop: K\_trap$\rightarrow$K\_mod, T\_network$\rightarrow$T\_bowtie, Phi\_sub$\rightarrow$Phi\_c \\
pipeline & hv1\_paper\_reproduction.syn (v6) & success=true, $\Delta$\xi\_CP=0.000, all 4 assertions pass \\
cofactor & cofactor(cooper\_pair, electron) & reconstructs K\_slow, G\_meso, T\_bowtie, $\Phi$\_c --- the phonon-dressed partner \\
principal\_decomp & principal\_decomp(Hv1\_human\_open) & atoms: atom{[}T=bowtie{]}, atom{[}Phi=c{]}, skeleton(D,R,P,F,K,G,Gamma,Omega) \\
\bottomrule
\end{tabularx}
\end{table}

\subsubsection{XXIII.3 Three Semantic Precision Points}

\textbf{1. The Exciton Theorem (tensor $\neq$ bound state)}
\texttt{tensor(electron, hole)} returns T\_linear --- the same topology as both inputs --- because topology promotion requires topological \emph{mismatch} to drive it upward. A bound excitonic state requires an additional \texttt{meet} with a Coulomb binding potential synthon. This is a feature: tensor gives the free ensemble prediction; bound-state formation is a separate lattice operation.

\textbf{2. T\_braid Special Rule}
T\_braid \otimes T\_braid $\rightarrow$ T\_braid. The braid group is closed under composition: braided topology does not network-promote when composed with itself. This is why \texttt{tensor(Majorana, Majorana)} correctly predicts topological qubit preservation --- the non-Abelian structure survives the tensor product without collapsing to a higher topology.

\textbf{3. Heyting Pseudocomplement Is Not Classical Negation}
\texttt{complement\_rel(s, proj, ref)} returns the Heyting pseudocomplement: the \emph{largest} synthon that, when met with \texttt{s}, falls below \texttt{proj}. This is not the Boolean complement. It answers ''what is the most you can add to the design while staying below threshold?'' --- a natural question in drug design (potency ceiling) and materials science (stability window). \texttt{satisfied=False} means no such synthon exists in the catalog that meets the criterion; the notes show per-axis contributions regardless.

\emph{See \texttt{TENSOR\_OPS\_DEMO.py} for 18 worked examples (§§1--7, \textasciitilde{}3 per operation) and \texttt{SYNTHONICON\_LANG.md §''Algebraic Operations Reference''} for the full formal specification.}

\begin{center}
\rule{0.5\textwidth}{0.4pt}
\end{center}

\subsection{XX. Development Roadmap}

\textbf{Phase 1 --- Grammar formalization and quantitative calibration (complete).} $F$ and $K$ are orthogonal primitives with anchored tier boundaries (v2.2: HIGH $\leq 8.5$ nats, MEDIUM 8.5--11.0 nats). $\Gamma$ encodes operator $\times$ tier. The axiom set includes grounding axioms 6 and 7, enforced at registration time. $P_{\pm}$ distinguishes symmetric and pseudosymmetric subclasses. $\Phi$ is a registered primitive. I(bits) is calibrated via DOF-counting pipeline (6--18 bits, system-specific). $\xi_{CP}$ table updated with uncertainty bands. Stoichiometry primitive $S$ implemented with Pass 4 audit enforcement. Batch criticality probe with degeneracy scoring operational.

\textbf{Phase 2 --- Anchor the criticality claim empirically (in progress).} The Varma QXY probe now provides quantitative degeneracy classification ($z_{\text{eff}}$ divergence confirmed, 2D percolation reference $z_{\text{eff}} = 1.330$ validated). Three independent scoring mechanisms are now operational: Varma QXY heuristics (Factors 1--5), steric-cliff proxy (Factor 6, db24c8 validated), and Frank-model classical bifurcation (Factor 7, Soai reaction validated). The CB{[}7{]} competitive displacement series (Kim 2001; Assaf \& Nau 2015) provides the first 6/6 experimental validation of the $F$-floor HotSwap ratchet --- confirming the $F_{\ell}$ tier at $K_a < 10^7$ M$^{-1}$ and the asymmetric displacement ordering. The proline-aldol Varma probe (ratio = 0.189) confirms $\Phi_{\text{sub}}$ as expected; the Soai probe (ratio = 0.94, score 0.920) identifies the highest-confidence $\Phi_c$ candidate to date. Remaining tasks: execute Transformation \#8 (rotaxane dethreading) with attention to near-critical barrier topology; run batch $\Phi_c$ scan on full catalog and publish top-10 candidates with degeneracy type; add one additional cooperative benchmark (e.g., quadruple H-bond array or GC base-pair mimic) to extend the cooperativity scaling calibration; resolve the architectural status of $\Phi$ as independent primitive versus derived condition.

\textbf{Phase 3 --- AI/ML integration and operational engineering (in progress, v0.4.4).} The LLM tool layer is now operational. \texttt{synthon\_tool.py} implements \texttt{SynthonTool} --- a real dispatch layer wrapping the live Python API into a structured \texttt{SYNTHON\_TOOL\_SCHEMA} (Anthropic/OpenAI format). \texttt{synthon\_agent.py} implements \texttt{SynthonDesignAgent} --- an autonomous design loop that proposes synthon encodings, validates them against all 7 axioms, probes $\Phi_c$, checks HotSwap path connectivity to a target, and self-corrects until convergence criteria are satisfied ($\Phi_c$ score $\geq$ threshold AND $\xi_{CP}$ $\leq$ threshold). Both are accessible from the CLI as \texttt{syncon tool} (single-shot dispatch) and \texttt{syncon design} (full agent loop). The LLM cannot hallucinate impossible chemistry: every proposal is immediately rejected with a precise axiom trace if it fails, making AI-driven design formally verifiable rather than heuristic. The \textbf{programmable matter domain} is now fully encoded: 11 PM synthon encodings, pairwise distance matrix, Primitive Jacobian, 4 tensor products with $\Phi_c$ propagation, DesignPipeline monad, programmability lattice (dynamic floor = $\Phi_c$ theorem), and 10 formal predictions P-38 through P-47. Phase 3e advances (v0.4.4): \textbf{$D_{\text{holo}}$} holographic dimensionality primitive added; \textbf{phase transitions as morphisms} fully implemented in \texttt{synthomnicon/morphism.py} (\texttt{syncon transition} CLI); \textbf{reflexive closure} experiment demonstrated (axioms encode themselves, $\Phi_c$ is invariant under intersection of its own axioms); \textbf{IIT comparison} formalized in §IV.1 METAPHYSICS.md (5-primitive conflict set, no HotSwap path, Meet\_Richness outscores IIT\_Phi). \textbf{Phase 3e stress tests complete:} QCP morphism demonstrated (GR$\rightarrow$Reuter FP$\rightarrow$SM, 2nd order, asymmetry=0.000); Floquet synthons registered (\texttt{floquet\_chern\_insulator} + \texttt{time\_crystal\_dtc}); $Z_2$ gauge theory encoded (\texttt{z2\_lattice\_gauge\_toric\_code}, d=3.90 to Moore-Read at 0.000 nat). \textbf{Catalog merge + dedup (2026-03-18):} 1,692 total entries (54 hand-curated + 1,638 discovery) deduplicated by notation to \textbf{223 unique synthons}; design agent catalog-naming bug fixed (three-layer: \texttt{\_slugify\_goal()}, \texttt{\_dispatch\_tool()} injection, \texttt{explicit\_name} parameter in both generator agents); conditional Varma QXY routing added (physical/chemical vs. $G_{\aleph}$-based for non-physical goals). \textbf{METAPHYSICS.md §X--XIV:} Perception as constraint-propagation compatibility (§X); convergence argument --- Kastrup/idealism + Glattfelder/ICR + SynthOmnicon all structurally equivalent (§XI); Tao Te Ching §20 as inside view (§XII); UCBFT and CCT external framework assessments (§XIII); architecture-specific $\Phi_c$ thresholds (§XIV). \textbf{PRIMITIVE\_PREDICTIONS.md:} P-ARCH-1--4 added ($D_{\text{holo}}$ \textasciitilde{}10T, $D_{\infty}$ \textasciitilde{}143T/ln(N), $D_{\wedge\bigtriangleup}$ \textasciitilde{}143T thresholds). Remaining tasks: multi-provider arbitrage ensemble; knowledge-graph export; rotaxane dethreading scan (Transformation \#8); Phase 3b Pydantic refinement types.

\begin{center}
\rule{0.5\textwidth}{0.4pt}
\end{center}


\end{document}
